% !TEX root = ../print.tex
% Chapter 4: the Matern pyramid. Draft section 4, lines 533-730 of
% draft/spatial-hemigroup-pyramids.md; its "Checks for section 4" (lines 731-742) are records, not
% transcribed. Moved in (adr/0004, decision 3): Lemma 6.2 of the draft (lines 1082-1089), the
% machine-checked lem:no-first-order-factorization of the affine blueprint, now in section 4.1.
% Moved out to chapter 5: the paragraph after Proposition 4.7 that lists the masks it covers, and
% the grid case of Theorem 4.10(4) (prop:grid-frame-step-not-selected).
%
% CHANGED (wave 3b, 2026-11-15, a strengthening): prop:thorin-in-dimension-d moves to [T] (it was
% [A] on A55); lem:cm-log-convex and lem:cm-exponential-rate are new [T] lemmas.
% CHANGED (wave 3b, 2026-11-15, a strengthening): prop:matern-increments no longer cites A55 (it
% was [A] on A39, A55); lem:cm-uncancelled-atom is a new [T] lemma. No node of the chapter cites A55.
% Nodes. Definition: def:matern-pyramid-data, def:frame-matern-pyramid,
% def:implicit-euler-principle. [T]:
% lem:bessel-profile, lem:no-first-order-factorization (Lean), prop:matern-generator-octave,
% prop:laplacian-pyramid, prop:frame-matern-pyramid-exists, prop:creation-plane, prop:creation-grid, prop:stage-lemma,
% prop:principle-other-readings, prop:principle-after-decimation (review round 2, batch C),
% prop:thorin-in-dimension-d, lem:cm-log-convex,
% lem:cm-exponential-rate, lem:cm-uncancelled-atom. [A]: lem:bessel-k (A51), thm:matern-pyramid
% (A35, A31, A23), prop:matern-increments (A39), thm:selection-stratum-two (A39).
%
% What stays prose, and why (one line each):
% - the chapter opening and the three section leads: they say what the sections prove;
% - the paragraph introducing L_r before its definition: it places the law in the literature
%   (Stein's Bessel potential, the isotropic paper's Matern kernel) and claims nothing it proves;
% - the comment after Theorem 4.1: it compares with the line and names the causal twin;
% - the paragraph on Bernstein functions before Proposition 4.2: it teaches the terms, and its one
%   claim, the subordination step, is ledger A57 itself, which no node here spends;
% - the comment after Proposition 4.2 and the paragraph before Proposition 4.3: explanation;
% - the generalized-gamma-convolution paragraph after Proposition 4.3: it places the mixing law
%   in Bondesson's class, a one-line observation about a definition, and no node rests on it;
% - the paragraphs around lem:no-first-order-factorization: they say what the lemma excludes and
%   what it does not; "degrees add" is a one-line observation, not used by any proof;
% - the comments after Propositions 4.4, 4.5, 4.6 and 4.7: explanation and pointers;
% - the remark after Definition 4.8 and the paragraph before Theorem 4.10: explanation;
% - the comments after Theorem 4.10 (the grid case's pointer, the picture of the Gaussian's
%   failure, the summary): the picture is explicitly not a proof, and nothing rests on it;
% - Remark 4.11(2) and (3): they say what the principle does not give; (3) stated an open question,
%   answered since review round 2 (batch C) by prop:principle-after-decimation.
%   Remark 4.11(1) argues, so it is the proposition prop:principle-other-readings.

\chapter{The \Matern{} pyramid}
\label{ch:matern}

The \Matern{} pyramid is the pyramid whose octave step is one Helmholtz solve. Each step from level
$t$ to level $qt$ applies the resolvent $(I - c^2t^2\Delta)^{-1}$. The operator $I - c^2t^2\Delta$ is
the modified Helmholtz, or screened Poisson, operator; this chapter calls it the Helmholtz operator
for short, and its inverse the Helmholtz resolvent. One such step is one implicit Euler step of the
heat equation, with a step size $c^2t^2$ that grows by $q^2$ per octave. This chapter constructs
the pyramid in $\RR^d$ and says why it is the pyramid to prefer.

Section~\ref{sec:matern-levels} gives its levels, its increments, its generator and its octave
step. In the plane the octave step admits no factorization into first-order passes, and its kernel
has no finite value at the centre, so it cannot be sampled. Section~\ref{sec:matern-counting} shows
that the line's reason for preferring the pyramid, that it creates no extrema, does not carry to
$d \ge 2$. Section~\ref{sec:matern-selection} gives the reason that does carry. Non-enhancement of
local extrema, read at one finite step, selects the Helmholtz resolvent among the rotation-invariant
laws (Proposition~\ref{prop:stage-lemma}). Required of the octave stages of an isotropic family
covariant under one ratio, it selects the \Matern{} pyramid up to the gauge
(Theorem~\ref{thm:selection-stratum-two}).

\section{The levels and the octave step}
\label{sec:matern-levels}

% draft: the paragraph before (4.1), line 539.
The octave step of the \Matern{} pyramid is the law $L_r$ with transform $(1 + r^2|\omega|^2)^{-1}$.
It is the \Matern{} law of index $1$ at range $r$, the kernel of
\prog{iso v0.1}{Prop.~3.12(2), p.~16} at scale $\sqrt2\,r$. Its density is $r^{-d}G_2(x/r)$, with
$G_2$ the Bessel potential of order $2$, the kernel of $(I - \Delta)^{-1}$
[\sscite{stein1970singular}, Ch.~V, \S 3, (26) and Prop.~2, p.~132], read in this paper's
convention for the transform. So $L_r$ is the Green's function of $I - r^2\Delta$. It is also a
Gaussian variance mixture: Brownian motion run for an exponential time, the simplest case of
subordination. The definition below takes that mixture as its starting point. The exponent of the
\Matern{} pyramid is the sum of the exponents of its octave steps at the finer scales, and the
polar profile of one octave step is the function $b_d$ below.

% draft: the definitions of L_r, F_q and b_d, line 539 and (4.1), line 540.
% CHANGED: L_r is defined as the Gaussian variance mixture, which the draft states as a property of
% the law with transform (1 + r^2|omega|^2)^{-1}; the transform becomes Lemma 4.2(1)
% (lem:bessel-profile), so that the definition presupposes no existence. The object is the same.
\begin{definition}[the octave step and the \Matern{} exponent]
\label{def:matern-pyramid-data}
\lean{Pyramids.maternStep, Pyramids.maternExponent, Pyramids.besselProfile}\leanok
Fix $d \ge 1$. For $r > 0$ let
\[
  L_r := \int_0^\infty e^{-s}\,N(0, 2sr^2I)\,ds ,
\]
the law of $r\sqrt{2X}\,Z$ with $X$ standard exponential and $Z$ standard normal on $\RR^d$
independent of it. Call $L_r$ the \emph{\Matern{} law of index $1$ at range $r$}, and write $L_r$
also for the operator $f \mapsto L_r * f$. For $q > 1$ and $c > 0$ put
\[
  F_q(r) := \sum_{i \ge 1}\log\bigl(1 + c^2q^{-2i}r^2\bigr), \qquad
  b_d(y) := \frac{2}{\Gamma(d/2)}\int_0^\infty v^{d/2-1}e^{-v-y^2/(4v)}\,dv ,
  \qquad r, y \ge 0 .
\]
The series converges, each term being at most $c^2q^{-2i}r^2$.
\statusT\quad\emph{(Vocabulary. The facts about these objects are Lemma~\ref{lem:bessel-profile}
and Lemma~\ref{lem:bessel-k}. The ranges $q > 1$ and $c > 0$ are not hypotheses of
\texttt{maternExponent}, a real \texttt{tsum}, which is the junk value $0$ where the series diverges,
as it does for $0 < q \le 1$ and $cr \ne 0$, nor fields of \texttt{IsMaternPyramid}; every statement that reads the series
carries $1 < q$, and the uses of \texttt{IsMaternPyramid} without it do not read the series.)}
% CHANGED (fidelity review, 2026-11-15, R8): the status line says which range hypotheses of the
% definition are not fields of its declarations, and that every consumer carries them.
\end{definition}

% draft: the claims of the paragraph before Theorem 4.1, lines 539-541 (the transform and density
% of L_r, the Green's function, the polar profile b_d(|x|/r), b_d(0) = 2, monotonicity, the bound,
% b_1), with the profile computation of the proof of Theorem 4.1(1), line 553.
% CHANGED: the paragraph's claims made a lemma, with their proofs. The bound's constant is
% C_d := 2^{1+d/2}, as in the draft; the two other constants the draft called C_d are renamed
% (BLUEPRINT-INPUTS; REFEREE-draft-sections-3-4.md, item 2, finding 2).
% CHANGED: clause (3) adds b_3(y) = 2(1+y)e^{-y}, which the draft states in the proof of
% Proposition 4.3(2); it is proved here from the integral, without K_{3/2}.
\begin{lemma}[the octave step and its profile]
\label{lem:bessel-profile}
\lean{Pyramids.bessel_profile_1, Pyramids.bessel_profile_1_green, Pyramids.bessel_profile_2,
Pyramids.bessel_profile_3, Pyramids.bessel_profile_3_closed}\leanok
\uses{def:matern-pyramid-data}
Let $d \ge 1$ and $r > 0$.
\begin{enumerate}
\item $\widehat{L_r}(\omega) = (1 + r^2|\omega|^2)^{-1}$. The law $L_r$ has the density
$r^{-d}G(x/r)$, where
\[
  G(x) := \int_0^\infty e^{-s}(4\pi s)^{-d/2}e^{-|x|^2/(4s)}\,ds \in (0, \infty]
\]
is continuous on $\RR^d \setminus \{0\}$. For every $g \in L^1(\RR^d)$,
$(I - r^2\Delta)(L_r * g) = g$ in the sense of distributions.
\item For every $\beta > 0$ and $\omega \in \RR^d$,
\[
  \log\bigl(1 + \beta|\omega|^2\bigr)
   = \int_{\RR^d}\bigl(1 - \cos\langle\omega, x\rangle\bigr)\,
     \frac{b_d(|x|/\sqrt\beta)}{|S^{d-1}|\,|x|^d}\,dx .
\]
\item $b_d(0) = 2$, and $b_d$ is continuous and strictly decreasing on $[0, \infty)$, with
$b_d(y) \le C_de^{-y/2}$, $C_d := 2^{1+d/2}$. Further $b_1(y) = 2e^{-y}$ and
$b_3(y) = 2(1 + y)e^{-y}$.
\end{enumerate}
% CHANGED (fidelity review, 2026-11-15, R26/R50/R69): the status line lists the declarations' hypotheses beside the node's (LINKAGE rule 6).
\statusT\quad\emph{(Cites nothing: Tonelli's theorem, Frullani's integral, a Gaussian integral,
the uniqueness of the transform of a probability measure and Fourier inversion on the Schwartz
space. Clause (2) says that $L_r$,
read through the polar decomposition of Definition~\ref{def:isotropic-origin-data}, has Gaussian coefficient
$0$ and polar profile $b_d(\cdot/r)$; that reading is used in Theorem~\ref{thm:matern-pyramid}(1).
The closed form of $b_d$ through $K_{d/2}$ cites entry A51 and is Lemma~\ref{lem:bessel-k}(1).
% CHANGED (2026-11-14, a strengthening): unused hypotheses dropped from the Lean; the print keeps them.
\textbf{Hypotheses} (the node's, each with what its declarations require): $d \ge 1$: same in
\texttt{bessel\_profile\_1}, \texttt{\_2}, \texttt{\_3}; weaker in \texttt{bessel\_profile\_1\_green},
which does not require it (dropped on 2026-11-14, unused); \texttt{bessel\_profile\_3\_closed} is stated
at $d = 1$ and $d = 3$ only, as the clause is; $r > 0$: same in \texttt{bessel\_profile\_1}; weaker
in \texttt{bessel\_profile\_1\_green}, which does not require it; absent from (2) and (3), which do not
read $r$; (1) $g \in L^1(\RR^d)$: same, \texttt{Integrable g} for real $g$, the distributional
identity tested against every real smooth compactly supported $\varphi$, as its definition requires;
(2) $\beta > 0$: same; (3) $y \ge 0$: same.)}
\end{lemma}

% CHANGED (wave 1c, 2026-11-13): the checked route. In (1) the density is identified through the transform of G and its uniqueness, G is shown smooth off the origin, and the Green identity is proved by moving L_r onto the test function and Fourier inversion on the Schwartz space, in place of the transform of a tempered distribution; in (2) Frullani's integral is proved by Tonelli's theorem; in (3) the closed forms are computed in the variable w with v = w^2, and b_3 by one integration by parts in w, in place of the recursion for Q_nu.
\begin{proof}
\leanok
\emph{(1)} The law $L_r$ is the image of $L_1$ under $x \mapsto rx$, and $L_1$ is the law of
$\sqrt{2X}\,Z$. Given $X = s$, the vector $\sqrt{2s}\,Z$ has the law $N(0, 2sI)$, with the transform
$e^{-s|\omega|^2}$ and the density $g_{2s}(x) := (4\pi s)^{-d/2}e^{-|x|^2/(4s)}$. Integrating against
$e^{-s}\,ds$ gives, by Tonelli's theorem, the transform
$\int_0^\infty e^{-s(1 + |\omega|^2)}\,ds = (1 + |\omega|^2)^{-1}$ of $L_1$. The function
$G = \int_0^\infty e^{-s}g_{2s}\,ds$ is integrable, and by Fubini's theorem its transform is the same
integral, $(1 + |\omega|^2)^{-1}$. By the uniqueness of the transform, $L_1$ has the density $G$.
The image under $x \mapsto rx$ then has the density $r^{-d}G(x/r)$ and the transform
$(1 + r^2|\omega|^2)^{-1}$. For $|x| \ge \delta > 0$ the integrand defining $G$ is at most
$e^{-s}(4\pi s)^{-d/2}e^{-\delta^2/(4s)}$, which is integrable in $s$. So $G$ is finite off the
origin, and, differentiating under the integral sign in $|x|^2$, it is smooth there.

For the Green's function let $g \in L^1(\RR^d)$, let $\varphi$ be real, smooth and of compact
support, and put $\psi := \varphi - r^2\Delta\varphi$. First, for every $z$,
\[
  \int\psi(z + y)\,L_r(dy) = \varphi(z) .
\]
Read $\varphi$ as a complex Schwartz function, and write $\mathcal{F}$ for the transform of the
$2\pi$ convention, $\mathcal{F}\varphi(\xi) := \int e^{-2\pi i\langle\xi, x\rangle}\varphi(x)\,dx =
\widehat\varphi(-2\pi\xi)$. Fourier inversion gives
$\varphi(x) = \int e^{2\pi i\langle\xi, x\rangle}\mathcal{F}\varphi(\xi)\,d\xi$, and the transform of
$\Delta\varphi$ is $-4\pi^2|\xi|^2\mathcal{F}\varphi$, so
$\psi(x) = \int e^{2\pi i\langle\xi, x\rangle}(1 + 4\pi^2r^2|\xi|^2)\mathcal{F}\varphi(\xi)\,d\xi$, an
absolutely convergent integral. By Fubini's theorem,
\[
  \int\psi(z + y)\,L_r(dy)
   = \int e^{2\pi i\langle\xi, z\rangle}(1 + 4\pi^2r^2|\xi|^2)\mathcal{F}\varphi(\xi)\,
     \widehat{L_r}(2\pi\xi)\,d\xi
   = \int e^{2\pi i\langle\xi, z\rangle}\mathcal{F}\varphi(\xi)\,d\xi = \varphi(z) .
\]
Second, $\psi$ is bounded, so $(x, y) \mapsto g(x - y)\psi(x)$ is integrable for
the product of $dx$ and $L_r(dy)$. Fubini's theorem and the translation $x = z + y$ give
\[
  \int(L_r * g)\,\psi = \int\int g(x - y)\psi(x)\,dx\,L_r(dy)
   = \int\int g(z)\psi(z + y)\,dz\,L_r(dy) = \int g(z)\varphi(z)\,dz .
\]
This is $(I - r^2\Delta)(L_r * g) = g$ in the sense of distributions.

\emph{(2)} Frullani's integral gives, for $u \ge 0$,
\[
  \log(1 + \beta u) = \int_0^\infty\bigl(1 - e^{-su}\bigr)\,e^{-s/\beta}\,\frac{ds}{s} .
\]
Both sides vanish at $u = 0$. For $u > 0$ put $a := 1/\beta$ and $b := u + 1/\beta$. The integrand is
$(e^{-as} - e^{-bs})/s = \int_a^be^{-ts}\,dt$, and Tonelli's theorem gives
$\int_a^b\int_0^\infty e^{-ts}\,ds\,dt = \int_a^bdt/t = \log(b/a) = \log(1 + \beta u)$.
Put $u = |\omega|^2$. Let $g_{2s}(x) := (4\pi s)^{-d/2}e^{-|x|^2/(4s)}$ be the density of
$N(0, 2sI)$. It is symmetric with mass $1$, so
$1 - e^{-s|\omega|^2} = \int(1 - \cos\langle\omega, x\rangle)\,g_{2s}(x)\,dx$. All integrands are
nonnegative, and Tonelli's theorem gives
$\log(1 + \beta|\omega|^2) = \int(1 - \cos\langle\omega, x\rangle)\,f_\beta(|x|)\,dx$ with
\[
  f_\beta(\rho) := \int_0^\infty(4\pi s)^{-d/2}e^{-\rho^2/(4s)}e^{-s/\beta}\,\frac{ds}{s} ,
  \qquad \rho > 0 .
\]
The substitution $s = \rho^2/(4v)$ has $ds/s = -dv/v$ and $(4\pi s)^{-d/2} = \pi^{-d/2}\rho^{-d}v^{d/2}$,
so $f_\beta(\rho) = \pi^{-d/2}\rho^{-d}\int_0^\infty v^{d/2-1}e^{-v-\rho^2/(4\beta v)}\,dv$. With
$|S^{d-1}| = 2\pi^{d/2}/\Gamma(d/2)$ this reads $|S^{d-1}|\,\rho^df_\beta(\rho) = b_d(\rho/\sqrt\beta)$.

\emph{(3)} At $y = 0$ the integral is $\Gamma(d/2)$, so $b_d(0) = 2$. For each $v > 0$ the integrand
is continuous and strictly decreasing in $y \ge 0$ and at most $v^{d/2-1}e^{-v}$, so $b_d$ is
continuous by dominated convergence and strictly decreasing. For the bound, split the exponent in
halves. The inequality of the means gives $\tfrac12\bigl(v + y^2/(4v)\bigr) \ge y/2$, and dropping
a term gives $\tfrac12\bigl(v + y^2/(4v)\bigr) \ge v/2$. So the integrand is at most
$e^{-y/2}v^{d/2-1}e^{-v/2}$, and
$b_d(y) \le e^{-y/2}\,\frac{2}{\Gamma(d/2)}\int_0^\infty v^{d/2-1}e^{-v/2}\,dv = 2^{1+d/2}e^{-y/2}$.

For the closed forms let $y > 0$; at $y = 0$ both read $b_d(0) = 2$. Put $\eta := y/2$ and
$E(w) := e^{-w^2 - \eta^2/w^2}$ for $w > 0$. The substitution $v = w^2$ gives
\[
  \int_0^\infty v^{\nu-1}e^{-v-y^2/(4v)}\,dv = 2\int_0^\infty w^{2\nu-1}E(w)\,dw .
\]
Since $w^2 + \eta^2/w^2 = 2\eta + (w - \eta/w)^2$, $E(w) = e^{-2\eta}e^{-(w - \eta/w)^2}$. Put
\[
  J_1 := \int_0^\infty e^{-(w - \eta/w)^2}\,dw , \qquad
  J_2 := \int_0^\infty \eta w^{-2}e^{-(w - \eta/w)^2}\,dw .
\]
The substitution $w \mapsto \eta/w$ maps $(0, \infty)$ onto itself and leaves $(w - \eta/w)^2$
unchanged, so $J_1 = J_2$. The map $\sigma(w) := w - \eta/w$ is an increasing bijection of
$(0, \infty)$ onto $\RR$ with the derivative $1 + \eta w^{-2}$, so
$J_1 + J_2 = \int_{\RR}e^{-\sigma^2}\,d\sigma = \sqrt\pi$. Hence $J_1 = J_2 = \tfrac12\sqrt\pi$, and
\[
  \int_0^\infty E(w)\,dw = \tfrac12\sqrt\pi\,e^{-2\eta} , \qquad
  \int_0^\infty w^{-2}E(w)\,dw = \frac{\sqrt\pi}{2\eta}\,e^{-2\eta} .
\]
At $\nu = 1/2$, with $\Gamma(1/2) = \sqrt\pi$, this gives
$b_1(y) = \frac{2}{\sqrt\pi}\cdot2\int_0^\infty E = 2e^{-y}$. At $\nu = 3/2$ the integral
$\int_0^\infty w^2E(w)\,dw$ is needed. The function $F(w) := wE(w)$ has the derivative
$F'(w) = (1 - 2w^2 + 2\eta^2w^{-2})E(w)$, and $F(w) \to 0$ as $w \to 0$ from above and as
$w \to \infty$, since $0 \le E(w) \le e^{-w^2}$. So $\int_0^\infty F' = 0$, that is
\[
  2\int_0^\infty w^2E(w)\,dw = \int_0^\infty E(w)\,dw + 2\eta^2\int_0^\infty w^{-2}E(w)\,dw
   = \tfrac12\sqrt\pi\,e^{-2\eta}(1 + 2\eta) .
\]
With $\Gamma(3/2) = \sqrt\pi/2$ this gives
$b_3(y) = \frac{4}{\sqrt\pi}\cdot\tfrac12\sqrt\pi\,(1 + y)e^{-y} = 2(1 + y)e^{-y}$.
% CHANGED (blueprint round, 2026-11-16, review flag "is itself proved from the", its similar
% places): the closing sentence "Every step is machine-checked, with no cited fact." is dropped; the
% proof cites nothing, and the provenance is the trust-base subsection's. In (1) the 2 pi transform
% is named \mathcal{F} and related to the functions' transform (review flag on the sign).
\end{proof}

% draft: the closed form of b_d through K_{d/2}, line 541, and the Bessel facts the draft uses in
% the proofs of Propositions 4.3(2) and 4.4(3), lines 590 and 602 (catalogue item 6; the
% librarian's finding in records/pyramids/VERBATIM-literature.md, section A51).
% CHANGED: a lemma of its own, so that the one cited fact behind the K_nu evaluations sits in one
% node; the draft's sentence "Only the integral form is used below" is dropped, as it is not true of
% Propositions 4.3(2) and 4.4(3). Every fact is derived here from DLMF 10.32.10; the DLMF equations
% that state the others directly are named in the proof.
% CHANGED (2026-11-13, approved by the author): Pyramids.bessel_k_1_profile is stated for every
% natural d, without the hypothesis d >= 1 of def:matern-pyramid-data, which its proof never read.
% The printed statement is unchanged.
% CHANGED (fidelity review, 2026-11-15, R66): \uses drops lem:bessel-profile, which the status line
% names only to say that K_{1/2}, K_{3/2} are not needed; neither the proof nor its declarations read
% it. b_d is def:matern-pyramid-data's.
\begin{lemma}[the profile through Bessel functions]
\label{lem:bessel-k}
\lean{Pyramids.bessel_k_1, Pyramids.bessel_k_1_profile, Pyramids.bessel_k_2, Pyramids.bessel_k_3}\leanok
\uses{def:matern-pyramid-data, thm:cited-A51}
Let $K_\nu$ be the modified Bessel function of the second kind of order $\nu$, and let $y > 0$.
\begin{enumerate}
\item For every real $\nu$,
\[
  \int_0^\infty v^{\nu-1}e^{-v-y^2/(4v)}\,dv = 2(y/2)^\nu K_\nu(y) .
\]
So $b_d(y) = \frac{4}{\Gamma(d/2)}(y/2)^{d/2}K_{d/2}(y)$, and $b_2(y) = 2yK_1(y)$.
\item $K_0$ is differentiable on $(0, \infty)$ with $K_0' = -K_1$, and $K_0(y) \to 0$ as
$y \to \infty$.
\item $K_0(y) = \int_1^\infty e^{-ys}(s^2 - 1)^{-1/2}\,ds$, and $K_0(y) - \log(1/y)$ stays bounded
as $y \downarrow 0$.
\end{enumerate}
% CHANGED (wave 2b, 2026-11-13): the annotation names what the checked proof reads (10.32.10 at
% the order -nu, and the definition 10.32.9 of the entry's addendum).
% CHANGED (fidelity review, 2026-11-15, R26/R50/R69): the status line lists the declarations' hypotheses beside the node's (LINKAGE rule 6).
\statusA\quad\ledger{A51}\quad\emph{(\textbf{Assignment.} \ledger{A51} carries one formula,
DLMF~10.32.10: $K_\nu(z) = \tfrac12(\tfrac12z)^\nu\int_0^\infty\exp\bigl(-t - z^2/(4t)\bigr)\,
t^{-\nu-1}\,dt$ for $|\mathrm{ph}\,z| < \pi/4$, with no condition on $\nu$. It is read at real
$z = y > 0$, where the phase condition holds, and at every real $\nu$, as clause (1) states; the
proof reads it at the order $-\nu$. The entry's addendum transcribes DLMF~10.32.9,
$K_\nu(z) = \int_0^\infty e^{-z\cosh t}\cosh(\nu t)\,dt$, the representation the proofs of (2) and (3)
start from. The chapter then uses (1) at $\nu = 0$, $1$ and $d/2$ only. Everything else is proved:
the evenness of $K_\nu$ in $\nu$ for (1); differentiation under the integral sign for (2); a
substitution and two elementary bounds for (3). DLMF~10.29.3 ($K_0' = -K_1$), 10.32.8 at
$\nu = 0$ (the integral of (3)), 10.30.3 ($K_0(z) \sim -\log z$) and 10.39.2 with 10.29.1
($K_{1/2}$, $K_{3/2}$) state those facts directly; the entry does not carry them, the proof does
not read them, and
$K_{1/2}$, $K_{3/2}$ are not used anywhere, since Lemma~\ref{lem:bessel-profile}(3) gives $b_1$ and
$b_3$ from the integral. \textbf{Not carried}: nothing about $K_\nu$ beyond the representations of
DLMF~10.32.10 and 10.32.9 at real positive argument. \textbf{The grade.} In the Lean $K_\nu$ is
\texttt{ScaleSpace.besselK}, defined at real argument by the integral of DLMF~10.32.9, not taken
from a library's classical $K_\nu$; the declaration of \ledger{A51},
\texttt{cited\_A51\_integral}, is a Lean theorem proved from that definition, and declares no
axiom. The grade [A] records that the identity is the cited DLMF~10.32.10, read from the page, and
that the definition is DLMF~10.32.9's. \textbf{Hypotheses} (the node's, each with what its
declarations require): $K_\nu$ the classical function: the declarations read
\texttt{besselK}, DLMF~10.32.9's integral at real argument, identified with $K_\nu$ by the citation;
$y > 0$: same in \texttt{bessel\_k\_1}, \texttt{bessel\_k\_1\_profile}, and in (2), (3) as their
statements place it; (1) $\nu$ real: same; (1) $d \ge 1$, from
Definition~\ref{def:matern-pyramid-data}, read in $b_d$: weaker,
\texttt{bessel\_k\_1\_profile} holds for every natural $d$, junk at $d = 0$, where
$\Gamma(0) = 0$ in Lean, by the author's decision of 2026-11-13.)}
\end{lemma}

% CHANGED (wave 2b, 2026-11-13): the proof follows the machine-checked route. (1) reads DLMF 10.32.10
% at the order -nu with K_{-nu} = K_nu, not the substitution t = y^2/(4v); (2) differentiates the
% definition DLMF 10.32.9 under the integral sign; (3) bounds K_0 near the origin by two elementary
% estimates on the definition, not by a comparison with the exponential integral.
% CHANGED (the theory paper, 2026-11-16): the closing sentence on checking and cited facts shortened; no step changed.
\begin{proof}
\leanok
\emph{(1)} DLMF~10.32.10 (\ledger{A51}), at $z = y$ and at the order $-\nu$, gives
$K_{-\nu}(y) = \tfrac12(\tfrac12y)^{-\nu}\int_0^\infty t^{\nu-1}e^{-t - y^2/(4t)}\,dt$. By
DLMF~10.32.9, $K_\nu(y) = \int_0^\infty e^{-y\cosh u}\cosh(\nu u)\,du$, which is even in $\nu$, so
$K_{-\nu} = K_\nu$ and the identity follows. The integral converges for every real $\nu$. The
substitution $v = \tfrac y2e^w$, under which $v + y^2/(4v) = y\cosh w$ and $dv/v = dw$, turns it into
$2(\tfrac y2)^\nu\int_0^\infty\cosh(\nu w)e^{-y\cosh w}\,dw$, and
$\cosh(\nu w)e^{-y\cosh w} \le e^{2\nu^2/y}e^{-(y/8)w^2}$ for $w \ge 0$, since
$\cosh(\nu w) \le e^{|\nu|w}$ and $\cosh w \ge w^2/4$. At $\nu = d/2$ the identity is the closed form
of $b_d$. At $d = 2$, $\Gamma(1) = 1$ gives $b_2(y) = 2yK_1(y)$.

\emph{(2)} By DLMF~10.32.9, $K_0(z) = \int_0^\infty e^{-z\cosh u}\,du$. For $|z - y| < y/2$ the
$z$-derivative of the integrand, $-\cosh u\,e^{-z\cosh u}$, is at most
$\cosh u\,e^{-(y/2)\cosh u}$ in absolute value, and the integral of this bound is $K_1(y/2)$, which is
finite by the bound in (1). So $K_0$ is differentiable at $y$ with
$K_0'(y) = -\int_0^\infty\cosh u\,e^{-y\cosh u}\,du = -K_1(y)$. For $y \ge 1$,
$(y - 1)(\cosh u - 1) \ge 0$ gives $y\cosh u \ge (y - 1) + \cosh u$, so
$0 \le K_0(y) \le e^{1-y}K_0(1)$, which tends to $0$ as $y \to \infty$.

\emph{(3)} The substitution $s = \cosh u$, with $ds = \sinh u\,du$ and
$\sinh u = (s^2 - 1)^{1/2}$, maps $(0,\infty)$ onto $(1,\infty)$ and turns
$K_0(y) = \int_0^\infty e^{-y\cosh u}\,du$ into the integral of (3), which therefore converges. For the
origin let $0 < y < 1$ and $L := \log(1/y) > 0$, so that $ye^L = 1$. On $(0, L]$ the integrand lies
between $1 - y\cosh u \ge 1 - ye^u$ and $1$, since $\cosh u \le e^u$. On $(L, \infty)$ it is positive
and, since $y\cosh u \ge \tfrac12ye^u = \tfrac12e^{u-L} \ge \tfrac12(u - L)$, at most
$e^{-(u-L)/2}$, whose integral over $(L,\infty)$ is $2$. Hence
\[
  L - (1 - y) = \int_0^L\bigl(1 - ye^u\bigr)\,du \le K_0(y) \le L + 2 ,
\]
and $|K_0(y) - \log(1/y)| \le 2$ for $0 < y < 1$.

The one cited fact is DLMF~10.32.10 (\ledger{A51}), read in (1); DLMF~10.32.9 is the definition of
$K_\nu$.
% CHANGED (blueprint round, 2026-11-16, review flag "is itself proved from the"): the closing
% sentence says what the proof cites, consistently with (1). It read "Every step is machine-checked,
% with no cited fact, as A51(1) is itself proved from the definition of K_nu, DLMF 10.32.9.";
% that A51's declaration is a Lean theorem is the status line's and the trust-base
% subsection's.
\end{proof}

% draft: Theorem 4.1 (section 4.1), lines 543-549 of draft/spatial-hemigroup-pyramids.md.
% CHANGED: clause (4) names its reading, f in the Schwartz class and the series converging in L^1 on
% W^{2,1}, as Proposition 4.4(1) states it; the pointer "(Proposition 4.4(1))" is dropped and the
% clause is proved here, so that the two nodes do not use each other.
% CHANGED: clause (1)'s density is claimed for t > 0 (the kernel at t = 0 is the point mass).
% CHANGED: clause (3) names the projection of f, the density of the image of f(x)dx under
% x -> <e, x>.
% CHANGED: clause (2) names the knot kernels as the kernels mu'_{0,q^n}, since the knots are
% defined in the canonical discrete gauge and the clause assumes only full covariance (blind
% review, flag 2).
% CHANGED: clause (3) narrowed and named. The draft says that the number of sign changes of each
% projection does not increase "from scale zero and from rung to rung". Both counts are stated with
% the hypothesis under which they are proved (a version of the projection of f the primitive of an
% integrable function), as the crossings of every level and the changes of monotonicity. The count
% from scale zero rests on [V-selection v0.1, Prop. 4.5(2)] (entry A23), the author's approval of
% 2026-11-12. The rung count rests on prop:ladder-non-creation-line(2). The scales are quantified,
% t > 0 and s > 0 (the kernel at t = 0 is the point mass; first blind review, flag 16). The record
% and the projected record are defined in the clause (second blind review, flags 3, 11, 12).
% CHANGED (fifth blind pass, 2026-11-15, flag 3): the rung-to-rung count in (3) says that it is
% under the same hypothesis on the projection of f, as its declaration matern_pyramid_3_counts
% takes it for both counts. The old wording began "For every s > 0 the projected record".
% CHANGED (fifth blind pass, 2026-11-15, flag 4): clause (4) names its reading, almost everywhere,
% the reading of its declaration matern_pyramid_4; the annotation said "more than the almost
% everywhere the clause states" of a clause that stated no reading.
% CHANGED (fidelity review, 2026-11-15, R24): the tag adds matern_pyramid_5_limit_exists, which
% states that the limit law of (5) exists at every scale t, and that every choice of the laws with
% the exponents l F_q(t|.|) converges weakly to it; matern_pyramid_5_weak took the limit law as a
% hypothesis. The statement is unchanged.
% CHANGED (fidelity review, 2026-11-15, R60): the tag adds Pyramids.bessel_profile_1, the transform
% (1 + r^2|omega|^2)^{-1} of the octave step L_r that clause (1) reads; it is lem:bessel-profile(1),
% which the node uses. The statement is unchanged.
\begin{theorem}[the \Matern{} pyramid in $\RR^d$]
\label{thm:matern-pyramid}
\lean{Pyramids.matern_pyramid_5_limit_exists, Pyramids.bessel_profile_1,
Pyramids.matern_pyramid_1_admissible, Pyramids.matern_pyramid_1_exists,
Pyramids.matern_pyramid_1_family, Pyramids.matern_pyramid_1_mixture,
Pyramids.matern_pyramid_2_equation, Pyramids.matern_pyramid_2_family,
Pyramids.matern_pyramid_3_marginal, Pyramids.matern_pyramid_3_polya,
Pyramids.matern_pyramid_3_counts, Pyramids.matern_pyramid_4_series, Pyramids.matern_pyramid_4,
Pyramids.matern_pyramid_5_limit, Pyramids.matern_pyramid_5_classA,
Pyramids.matern_pyramid_5_multiples, Pyramids.matern_pyramid_5_weak,
Pyramids.matern_pyramid_5_kernels}
\leanok
\uses{def:matern-pyramid-data, lem:bessel-profile, def:family-dimension-d,
def:isotropic-origin-data, def:isotropic-cascade-family, thm:cited-A35, thm:cited-A34,
thm:cited-A31, thm:matern-pyramid-line,
thm:matern-pyramid-line-noncreation, prop:ladder-non-creation-line, thm:cited-A23}
Fix $q > 1$, $c > 0$ and $d \ge 1$.
\begin{enumerate}
\item \emph{(Admissibility.)} $F_q \in \mathcal{A}_d$, with Gaussian coefficient $0$ and the
nonincreasing polar profile $k_d(r) = \sum_{i \ge 1}b_d(rq^i/c)$. So the transforms
$\exp\bigl[-\bigl(F_q(t|\omega|) - F_q(s|\omega|)\bigr)\bigr]$, $0 \le s \le t$, define an isotropic
family with $S_\lambda t = \lambda t$ for every $\lambda > 0$, the \emph{\Matern{} pyramid}. Its
octave step is one Bessel potential, $\mu_{t,qt} = L_{ct}$ for $t > 0$. Its kernel is the Gaussian variance
mixture $\mu_t = \mathrm{law}\bigl(\sqrt{2V_t}\,Z\bigr)$, with
$V_t = \sum_{i \ge 1}c^2q^{-2i}t^2X_i$, the $X_i$ independent standard exponential and $Z$ standard
normal on $\RR^d$ independent of them, and for $t > 0$ it has a $C^\infty$ density with all
derivatives bounded.
\item \emph{(Uniqueness.)} If $\widetilde F$ is continuous on $[0, \infty)$ with $\widetilde F(0) = 0$
and $\widetilde F(qr) - \widetilde F(r) = \log(1 + c^2r^2)$ for $r \ge 0$, then $\widetilde F = F_q$.
An isotropic fully covariant family $(\mu'_{s,t})$ whose kernels $\mu'_{0,q^n}$, $n \in \ZZ$, have
the transforms $\exp\bigl(-F_q(q^n|\omega|)\bigr)$ is the family of (1) read in a gauge $\chi$ with
$\chi(q^n) = q^n$ for every $n$.
\item \emph{(Projections.)} Along every unit vector $e$, the marginal family
$\bigl(\langle e, \cdot\rangle_*\mu_{s,t}\bigr)$ is the \Matern{} pyramid on the line of
Theorem~\ref{thm:matern-pyramid-line}. For $f \in L^1(\RR^d)$ and $t > 0$ call $\mu_t * f$ the
\emph{record} of $f$ at scale $t$. Its projection along $e$, the density of the image of
$(\mu_t * f)(x)\,dx$ under $x \mapsto \langle e, x\rangle$, is the \emph{projected record}. It is
the projection of $f$ convolved with a \Polya{} frequency density. Suppose that the projection of
$f$ has a version that is the primitive of an integrable function. Then for every $t > 0$ the
projected record at scale $t$ crosses each level no more often, and changes monotonicity no more
often, than that version. Under the same hypothesis, for every $s > 0$ the projected record at
scale $qs$ crosses each level no more often, and changes monotonicity no more often, than the
projected record at scale $s$.
\item \emph{(Generator.)} For $f \in \mathcal{S}(\RR^d)$ and $t > 0$,
$\partial_t(\mu_t * f) = \mathcal{H}_t(\mu_t * f)$ almost everywhere, with
$\mathcal{H}_t = \frac2t\sum_{i \ge 1}(L_{ctq^{-i}} - I)$, a series of bounded operators on
$L^1(\RR^d)$ that converges absolutely in $L^1$ on $W^{2,1}(\RR^d)$ and whose sum is not bounded in
the $L^1$ norm.
\item \emph{(The limit $q \downarrow 1$.)} With $\ell = \log q$, $\ell F_q(r) \to \Lambda(r) :=
-\tfrac12\operatorname{Li}_2(-c^2r^2)$ locally uniformly as $q \downarrow 1$, and
$\Lambda \in \mathcal{A}_d$ for every $d$, with polar profile $\int_0^\infty b_d(re^u/c)\,du$. The
laws with exponent $\ell F_q(t|\cdot|)$ exist, and they converge weakly at every scale to the law
with exponent $\Lambda(t|\cdot|)$. The kernels $\mu_t$, $t > 0$, themselves do not converge:
$\widehat{\mu_t}(\omega) = e^{-F_q(t|\omega|)} \to 0$ for $\omega \ne 0$.
% CHANGED (fifth blind pass, 2026-11-15, flags 1-2, the author's decision): "t > 0" added to the
% last sentence of (5), false at t = 0 where mu_0 = delta_0 for every q (matern_pyramid_5_kernels
% takes 0 < t), and to the octave step mu_{t,qt} = L_{ct}, L_r being defined for r > 0 only.
% CHANGED (blueprint round, 2026-11-16, review flag "is closed under positive multiples, the"): the
% two rationales of (5), "Since A_d is closed under positive multiples" and "since F_q -> infinity",
% move into its proof, which already gives both. The claims are unchanged; F_q(r) -> infinity, a
% conjunct of matern_pyramid_5_kernels, is now stated in the proof only.
\end{enumerate}
% CHANGED (wave 2g, 2026-11-14): A34 off the status line; thm:cited-A34 is regraded [T], a theorem
% with no axiom, so the node spends one cited fact fewer.
% CHANGED (fidelity review, 2026-11-15, R26/R50/R69): the status line lists the declarations' hypotheses beside the node's (LINKAGE rule 6).
\statusA\quad\ledger{A35}\quad\ledger{A31}\quad\ledger{A23}\quad\emph{(\textbf{Assignment.}
\ledger{A35}, \prog{iso v0.1}{Prop.~3.11(1)--(2), pp.~15--16}, carries the membership in
$\mathcal{A}_d$ of the exponent of an isotropic origin datum (Definition~\ref{def:isotropic-origin-data}):
it is read in (1) for $F_q$ and in (5) for $\Lambda$ and for $\ell F_q$. The realization of the
family from $(\chi, F)$ in (1), at $\chi = \mathrm{id}$, is Theorem~\ref{thm:cited-A34},
\prog{iso v0.1}{Lem.~3.8, p.~14}, proved with no axiom and cited here as no ledger entry. \ledger{A31}, \prog{iso v0.1}{Thm.~3.6, p.~13}, the classification in
$\RR^d$, carries the form of the family in the second half of (2). Proved here: the polar profile
(Lemma~\ref{lem:bessel-profile}(2)), the two integrability conditions, the telescoping, the mixture,
the density, the functional equation of (2), the generator (4) and the Riemann sum and limits of
(5), with Lévy's continuity theorem. \ledger{A23}, \prog{V-selection v0.1}{Prop.~4.5(2), p.~35}
(derived there from clause (1) on p.~39), carries the count from scale zero in (3): for the
similarity family of an exponent $a\omega^2 + \sum_i\log(1 + \omega^2/\theta_i^2)$ and a signal that is
the primitive of an integrable function, the record at every scale has no more sign changes of
the record minus a level, and of its derivative, than the signal. It is read for the line's
pyramid, $a = 0$, $\theta_i = q^i/c$, and for the signal a version of $P_ef$. The rest of clause (3) rests on the line's
Theorem~\ref{thm:matern-pyramid-line-noncreation} and Proposition~\ref{prop:ladder-non-creation-line},
which carry their own entries (A22, A24 and A23 again), reached through \texttt{uses}. The
identification of the projected records with the line's records is proved here.
\textbf{Not carried}: no class property of $\mu_t$ in $\RR^d$ beyond the Gaussian mixture; in
particular no \Polya{} frequency property in $d \ge 2$, which (3) claims of the projections only.
% CHANGED (wave 3b, 2026-11-15): every clause machine-checked; the old text priced the mixture
% in (1) and clause (4) as left.
\textbf{Machine-checked} (wave 3b): every clause, all sixteen declarations. The density in (1) is
proved from Fourier analysis on Schwartz functions, the mixture in (1) from the independence of
the $X_i$ and the line's product display, and clause (4) from the Green identity of
Lemma~\ref{lem:bessel-profile}(1) and the line's derivative of Theorem~\ref{thm:matern-pyramid-line}(3);
none of them cites anything. The proof gives the identity of (4) at every point, with the series
summed pointwise, which is more than the almost everywhere that the clause and its declaration
state.
% CHANGED (2026-11-15, a strengthening): six declarations drop binders their proofs do not read.
The print keeps the node's setting. \textbf{Hypotheses} (the node's, each with what its
declarations require): $q > 1$: same, except weaker in \texttt{matern\_pyramid\_3\_marginal}, which
does not require it, and read as the filter $q \downarrow 1$ through $q > 1$ in (5); $c > 0$: same,
except weaker in \texttt{matern\_pyramid\_3\_marginal}; $d \ge 1$: same in
\texttt{bessel\_profile\_1}, \texttt{\_1\_admissible}, \texttt{\_1\_exists}, \texttt{\_1\_family},
\texttt{\_2\_family}, \texttt{\_4\_series} (its unboundedness reads it, the test frequency lying on
an axis), \texttt{\_5\_classA}, \texttt{\_5\_multiples}, \texttt{\_5\_kernels},
\texttt{\_5\_limit\_exists}; weaker, not required, in \texttt{\_1\_mixture}, \texttt{\_2\_equation},
\texttt{\_3\_marginal}, \texttt{\_3\_polya}, \texttt{\_3\_counts}, \texttt{matern\_pyramid\_4},
\texttt{\_5\_limit}, \texttt{\_5\_weak}; the family $(\mu_{s,t})$ of (1): the declarations that
read it quantify over every $\mu$ with \texttt{IsMaternPyramid q c}, the transforms of (1), a
specification that \texttt{\_1\_exists} meets and the uniqueness of the transform makes unique:
same. (1) $t > 0$ for the octave step and the density: same; the mixture: same, for every
probability space carrying measurable $X_i$ and $Z$ with those laws and independences, at every
$t \ge 0$; no declaration constructs such a space. (2) $\widetilde F$ continuous on $[0, \infty)$,
$\widetilde F(0) = 0$, the equation for $r \ge 0$: same; $(\mu'_{s,t})$ isotropic and fully
covariant: same, \texttt{IsIsotropicCascadeFamily} for some $S$; the knot transforms for every
$n \in \ZZ$: same. (3) $e$ a unit vector: same; $f \in L^1(\RR^d)$: same, real $f$; $t > 0$,
$s > 0$: same; a version of the projection of $f$ that is the primitive of an integrable
function: same, a projection density $g$ (integrable) with $g(x) = C + \int_0^xg'$ for every $x$
and $g'$ integrable. (4) $f \in \mathcal{S}(\RR^d)$: same, real $f$; $t > 0$: same; the domain
$W^{2,1}(\RR^d)$ of the series: same, \texttt{IsW21}. (5) every scale $t$: weaker in
\texttt{\_5\_limit\_exists} and \texttt{\_5\_weak}, at every real $t$; same, $t \ge 0$, in
\texttt{\_5\_multiples}; the limit law: \texttt{\_5\_weak} takes it as a hypothesis (stronger
alone), \texttt{\_5\_limit\_exists} constructs it and requires nothing of it (same); $t > 0$ in the
last sentence: same in \texttt{\_5\_kernels}, which quantifies over every choice of the \Matern{}
pyramids $\mu^{(q)}$.)}
\end{theorem}

% CHANGED (fifth blind pass, 2026-11-15, flag 5): in (4) the three inverse-transform displays drop
% the real part, "Re", which was right only for real f; Fourier inversion gives them without it.
% CHANGED (fifth blind pass, 2026-11-15, flag 6): in (4) the passage of weak derivatives through
% L_r is said for every derivative of order at most 2, which W^{2,1} requires; it was written for the
% second derivatives along the axes only.
% CHANGED (fidelity review, 2026-11-15, R66): in (3) the marginal transform is read as the transform
% of an image measure, as the declaration matern_pyramid_3_marginal proves it (charFun_marginal);
% lem:reduction-keeps-stratum is named as the same identity for isotropic families and leaves the
% node's \uses, since neither proof reads it.
% CHANGED (the theory paper, 2026-11-16): the closing sentence on checking and cited facts shortened; no step changed.
\begin{proof}
\leanok
\emph{(1)} By Lemma~\ref{lem:bessel-profile}(2) with $\beta = c^2q^{-2i}$,
\[
  \log\bigl(1 + c^2q^{-2i}|\omega|^2\bigr)
   = \int_{\RR^d}\bigl(1 - \cos\langle\omega, x\rangle\bigr)\,
     \frac{b_d(|x|q^i/c)}{|S^{d-1}|\,|x|^d}\,dx .
\]
Summing over $i$, by Tonelli's theorem,
\[
  F_q(|\omega|) = \int_{\RR^d}\bigl(1 - \cos\langle\omega, x\rangle\bigr)\,
   \frac{k_d(|x|)}{|S^{d-1}|\,|x|^d}\,dx , \qquad k_d(r) = \sum_{i \ge 1}b_d(rq^i/c) .
\]
This is the exponent of the datum $(0, k_d)$ in the normalization of
Definition~\ref{def:isotropic-origin-data}, whose \Levy{} measure has the radial density
$k_d(|x|)\bigl(|S^{d-1}|\,|x|^d\bigr)^{-1}$. The profile $k_d$ is nonincreasing as a sum of
nonincreasing functions (Lemma~\ref{lem:bessel-profile}(3)).
% CHANGED (wave 3b, 2026-11-15)
% The checked route bounds the profile by M/r and derives both integrability conditions from that
% one bound; the old text integrated b_d(y) y and b_d(y)/y after the substitution y = rq^i/c.
Its integrability conditions follow from one bound. Since $e^{-x} \le 1/x$ for $x > 0$, the bound
$b_d(y) \le C_de^{-y/2}$ gives $b_d(y) \le 2C_d/y$, so
\[
  k_d(r) \le \sum_{i \ge 1}\frac{2C_dc}{rq^i} = \frac{M}{r} , \qquad M := \frac{2C_dc}{q - 1} ,
  \qquad r > 0 .
\]
Hence $k_d(r)\,r \le M$ on $(0, 1)$ and $k_d(r)/r \le M/r^2$ on $(1, \infty)$, and
$\int_0^1k_d(r)\,r\,dr \le M$, $\int_1^\infty k_d(r)\,dr/r \le M$. So $F_q \in \mathcal{A}_d$ by \ledger{A35}. It is not identically $0$, since $b_d > 0$, and that is the hypothesis
$F \not\equiv 0$ under which Theorem~\ref{thm:cited-A34} realizes a family.
% CHANGED (2026-11-23, review round 3, s4-5.S16): wording.
By Theorem~\ref{thm:cited-A34} with $\chi = \mathrm{id}$ there is an isotropic family with
$\widehat{\mu_{s,t}}(\omega) = \exp\bigl[-\bigl(F_q(t|\omega|) - F_q(s|\omega|)\bigr)\bigr]$. Its
covariance is read off the transforms: $\widehat{\mu_{s,t}}(\lambda\omega) =
\widehat{\mu_{\lambda s,\lambda t}}(\omega)$, so $D_\lambda\mu_{s,t} = \mu_{\lambda s,\lambda t}$ by
the uniqueness of the transform. The telescoping is
$F_q(qr) = \sum_{i \ge 1}\log(1 + c^2q^{-2(i-1)}r^2) = \log(1 + c^2r^2) + F_q(r)$. It gives
$\widehat{\mu_{t,qt}}(\omega) = (1 + c^2t^2|\omega|^2)^{-1}$, so $\mu_{t,qt} = L_{ct}$ by
Lemma~\ref{lem:bessel-profile}(1).

% CHANGED (wave 3b, 2026-11-15)
% The checked route: the law of (V_t, Z) as a product, Fubini for the transform, and dominated
% convergence on the partial sums; the old text read monotone convergence and conditioning on V_t.
For the mixture let $t \ge 0$ and put $a_i := c^2q^{-2i}t^2$. The exponential law is carried by
$(0, \infty)$ and has mean $1$, so by monotone convergence $E\sum_ia_iX_i = \sum_ia_i < \infty$, and
the series $V_t = \sum_ia_iX_i$ of nonnegative terms converges almost surely; what it is on the
exceptional null set does not change the law of $\sqrt{2V_t}\,Z$. The variable $V_t$ is a
measurable function of the sequence $(X_i)$, which is independent of $Z$, so the pair $(V_t, Z)$ has
the product law $\nu \otimes N(0, I)$, with $\nu$ the law of $V_t$. For fixed $v \ge 0$ the variable
$\sqrt{2v}\,Z$ has the transform $e^{-v|\omega|^2}$, so by Fubini's theorem the law of
$\sqrt{2V_t}\,Z$ has the transform $E\bigl[e^{-V_t|\omega|^2}\bigr]$. Let $u \ge 0$. For the
partial sums $V_t^{(N)} := \sum_{i \le N}a_iX_i$, independence and the transform
$E[e^{-uX_i}] = (1 + u)^{-1}$ of the exponential law give
$E\bigl[e^{-uV_t^{(N)}}\bigr] = \prod_{i=1}^N(1 + a_iu)^{-1}$. As $N \to \infty$ the left side tends
to $E[e^{-uV_t}]$ by dominated convergence, the integrands lying in $[0, 1]$, and the right side
tends to $\prod_{i \ge 1}(1 + a_iu)^{-1} = e^{-F_q(t\sqrt u)}$, the product (\ref{eq:line-e7}) of
Theorem~\ref{thm:matern-pyramid-line}. At $u = |\omega|^2$ this is
$e^{-F_q(t|\omega|)} = \widehat{\mu_t}(\omega)$, as $F_q(0) = 0$. So $\mu_t$ is the law of
$\sqrt{2V_t}\,Z$, by the uniqueness of the transform.

\emph{The density.} Let $t > 0$. For every integer $N \ge 1$, keeping the first $N$
terms of the series,
\[
  e^{-F_q(t|\omega|)} \le \prod_{i=1}^N\bigl(1 + c^2q^{-2i}t^2|\omega|^2\bigr)^{-1}
   \le \bigl(1 + c^2q^{-2N}t^2|\omega|^2\bigr)^{-N} .
\]
% CHANGED (wave 3b, 2026-11-15)
% The checked route proves that p_t is the density through Schwartz test functions, its
% nonnegativity by a bump, its mass by monotone convergence and the identity of the laws on
% transforms; the old text read "Fourier inversion gives the density" as one step.
So $|\omega|^m e^{-F_q(t|\omega|)}$ is integrable for every $m \ge 0$: with $b := \min(1,
c^2t^2q^{-2N})$ and $2N > m + d$, it is at most
$b^{-N}|\omega|^m(1 + |\omega|^2)^{-N} \le b^{-N}(1 + |\omega|^2)^{-(2N - m)/2}$. Put
\[
  p_t(x) := (2\pi)^{-d}\int e^{-i\langle\omega, x\rangle}e^{-F_q(t|\omega|)}\,d\omega .
\]
It is real, as the integrand is even in $\omega$ up to conjugation, and differentiation under the
integral sign gives every derivative, bounded by
$(2\pi)^{-d}\int|\omega|^{|\alpha|}e^{-F_q(t|\omega|)}\,d\omega$. For a Schwartz function $\Phi$,
Fubini's theorem, Fourier inversion on the Schwartz space, the evenness of
$e^{-F_q(t|\omega|)}$ and the self-adjointness of the transform give
$\int\Phi\,p_t\,dx = \int\Phi\,d\mu_t$. Three consequences follow. First, $p_t \ge 0$: if
$p_t(x_0) < 0$, then $p_t < 0$ on a ball around $x_0$, and a smooth bump $\Phi \ge 0$ supported in
that ball with $\Phi(x_0) = 1$ has $\int\Phi\,p_t\,dx < 0 \le \int\Phi\,d\mu_t$. Second,
$\int p_t\,dx \le 1$: smooth bumps $0 \le \Phi_n \le 1$ equal to $1$ on the ball of radius $n$
have $\int\Phi_np_t\,dx = \int\Phi_n\,d\mu_t \le 1$, and Fatou's lemma passes to the limit. Third,
$\mu_t = p_t\,dx$: both are finite measures, and testing on $\Phi_n(x)e^{i\langle\omega, x\rangle}$,
which are Schwartz functions, and letting $n \to \infty$ by dominated convergence shows that their
transforms agree at every $\omega$.

\emph{(2)} Applying the equation at $q^{-n}r$ and summing over $n = 1, \dots, N$ gives
$\widetilde F(r) = \widetilde F(q^{-N}r) + \sum_{n=1}^N\log(1 + c^2q^{-2n}r^2)$. The first term tends
to $\widetilde F(0) = 0$ by continuity, and the sum tends to $F_q(r)$. For the family, let
$(\mu'_{s,t})$ be an isotropic fully covariant family (Definition~\ref{def:isotropic-cascade-family}) with the
stated kernels $\mu'_{0,q^n}$. By \ledger{A31} there are a gauge $\chi$, an
increasing homeomorphism of $[0, \infty)$, and an $F \in \mathcal{A}_d$ with
$\widehat{\mu'_{s,t}}(\omega) = \exp\bigl(-[F(\chi(t)|\omega|) - F(\chi(s)|\omega|)]\bigr)$, the pair
unique up to $(\chi, F) \mapsto (a\chi, F(a^{-1}\,\cdot))$, $a > 0$. Choose $a$ so that
$\chi(1) = 1$. Since $\chi(0) = 0$ and $F(0) = 0$, the kernel $\mu'_{0,1}$ has the transform
$e^{-F(|\omega|)} = e^{-F_q(|\omega|)}$, so $F = F_q$. At the scale $q^n$,
$F_q(\chi(q^n)r) = F_q(q^nr)$ for every $r$. As $F_q$ is strictly increasing, $\chi(q^n) = q^n$. So
$\mu'_{s,t} = \mu_{\chi(s),\chi(t)}$.

\emph{(3)} Write $\pi := \langle e, \cdot\rangle$. The transform of an image measure is
$\widehat{\pi_*\mu}(r) = \widehat\mu(re)$, so $\pi_*\mu_{s,t}$ has the transform
$r \mapsto \widehat{\mu_{s,t}}(re) = \exp\bigl[-\bigl(F_q(t|r|) - F_q(s|r|)\bigr)\bigr]$. (For an
isotropic family this is Lemma~\ref{lem:reduction-keeps-stratum}'s identity
$\widehat{\mu_{s,t}}(\omega) = \widehat{m_{s,t}}(|\omega|)$ for the marginal laws
$m_{s,t} := \pi_*\mu_{s,t}$; the proof needs only the transform of the image.) These are the
transforms of the family of Theorem~\ref{thm:matern-pyramid-line}, with the same $q$ and $c$, and
the uniqueness of the transform identifies the laws. For $f \in L^1(\RR^d)$ write $P_ef$ for the
density of $\pi_*(f(x)\,dx)$, in $L^1(\RR)$ by Fubini's theorem. For a probability measure $\mu$
and bounded measurable $\phi$ on $\RR$, Fubini's theorem gives
\[
  \begin{aligned}
  \int\phi(\langle e, x\rangle)\,(\mu * f)(x)\,dx
   &= \iint\phi\bigl(\langle e, y\rangle + \langle e, z\rangle\bigr)f(y)\,dy\,\mu(dz) \\
   &= \iint\phi(\sigma + w)\,P_ef(\sigma)\,d\sigma\,(\pi_*\mu)(dw) ,
  \end{aligned}
\]
so $P_e(\mu * f) = (\pi_*\mu) * P_ef$. With $\mu = \mu_t$ this is the projection of $f$ convolved with
$\pi_*\mu_t$, a \Polya{} frequency density by Theorem~\ref{thm:matern-pyramid-line-noncreation}.
So the projected record at scale $t$ is $(\pi_*\mu_t) * P_ef$, the record at scale $t$ of the
signal $P_ef$ in the line's \Matern{} pyramid. For the counts let $g$ be a version of $P_ef$ that
is the primitive of an integrable $g'$. A convolution does not see the choice of version, so
$(\pi_*\mu_t) * P_ef = (\pi_*\mu_t) * g$ for every $t > 0$. The count from rung to rung is
Proposition~\ref{prop:ladder-non-creation-line}(2) for the signal $g$ at the scale $s$, whose
hypotheses hold: the line's \Matern{} pyramid has the exponent of that proposition with $a = 0$ and
$\vartheta_i = q^i/c$, so $\sum_i\vartheta_i^{-2} < \infty$ and $\mathrm{mult}(q\vartheta) -
\mathrm{mult}(\vartheta)$ is $1$ at $\vartheta = 1/c$ and $0$ elsewhere, and $g$ is the primitive of
an integrable function. The count from scale zero is the counting clause
\ledger{A23} for the same exponent and the same
signal, whose hypotheses hold as $\sum_i(q^i/c)^{-2} < \infty$ and $g$ is the primitive of an
integrable $g'$; it gives $S^-(u(t, \cdot) - y) \le S^-(g - y)$ for every level $y$ and
$S^-(\partial_xu(t, \cdot)) \le S^-(g')$ for the record $u(t, \cdot) = (\pi_*\mu_t) * g$ at every
$t > 0$, $S^-$ counting sign changes.

% CHANGED (wave 3b, 2026-11-15)
% The checked route. The bound on W^{2,1} comes from the Green identity of Lemma 4.2(1) and the
% passage of weak derivatives through L_r, not from the transform of a weak derivative; the
% unbounded sum is shown on a bump of small support, not on a modulated Gaussian; the derivative
% in t is the line's (thm:matern-pyramid-line(3)), and the generator identity is read pointwise
% through the inverse transform, not by Fourier inversion almost everywhere.
\emph{(4)} For $r > 0$ the law $L_r$ is a probability measure, so $\|L_r * g\|_1 \le \|g\|_1$ and
$L_r - I$ is bounded on $L^1$, with norm at most $2$. Let $g \in W^{2,1}(\RR^d)$, with the weak
second derivatives $v_j := \partial_j^2g \in L^1$. Weak derivatives pass through a finite measure:
for a test function $\varphi$, Fubini's theorem and the translation $x = z + y$ give
\[
  \int(L_r * g)\,\partial_j^2\varphi
   = \iint g(z)\,\partial_j^2\varphi(z + y)\,dz\,L_r(dy)
   = \iint v_j(z)\,\varphi(z + y)\,dz\,L_r(dy) = \int(L_r * v_j)\,\varphi ,
\]
a translate of a test function being a test function. The same computation with
$\partial^\alpha\varphi$ in place of $\partial_j^2\varphi$, for every multi-index $|\alpha| \le 2$,
gives $\int(L_r * g)\,\partial^\alpha\varphi = (-1)^{|\alpha|}\int(L_r * \partial^\alpha g)\,\varphi$,
with $\partial^\alpha g \in L^1$ the weak derivative of $g$. So $L_r * g$ has the weak derivatives
$L_r * \partial^\alpha g$, of $L^1$ norm at most $\|\partial^\alpha g\|_1$, and lies in
$W^{2,1}$; the argument reads of $L_r$ only that it is a probability measure. With the Green identity
$\int(L_r * g)(\varphi - r^2\Delta\varphi) = \int g\varphi$ of Lemma~\ref{lem:bessel-profile}(1),
for real $\varphi$, this gives
\[
  \int\Bigl(L_r * g - g - r^2\sum_jL_r * v_j\Bigr)\varphi = 0
\]
for every real smooth compactly supported $\varphi$. So $L_r * g - g = r^2L_r * \Delta g$ almost
everywhere, with $\Delta g := \sum_jv_j$, and $\|(L_r - I)g\|_1 \le r^2\sum_j\|v_j\|_1$. With
$r_i := ctq^{-i}$, $\sum_ir_i^2 = c^2t^2/(q^2 - 1)$, the series $\mathcal{H}_tg$ converges
absolutely in $L^1$.

The sum is not bounded. The transform of an integrable $g$ is that of the measure $g(x)\,dx$,
$\widehat g(\omega) := \int e^{i\langle\omega, x\rangle}g(x)\,dx$, so that
$|\widehat g(\omega)| \le \|g\|_1$. For $g \in W^{2,1}$ the absolute convergence lets the series
be integrated term by term, and Fubini's theorem gives $\widehat{L_r * g} = \widehat{L_r}\,\widehat g$,
so
\[
  \widehat{\mathcal{H}_tg}(\omega)
   = \frac2t\sum_{i \ge 1}\Bigl(\frac{1}{1 + r_i^2|\omega|^2} - 1\Bigr)\widehat g(\omega) .
\]
Suppose that $\|\mathcal{H}_tg\|_1 \le B\|g\|_1$ for every $g \in W^{2,1}$, and let $N \ge 1$ be an
integer. Take $\omega$ with $|\omega| = q^N/(ct)$, which exists as $d \ge 1$. Every term of the sum
is at most $0$, and for $i \le N$, $r_i|\omega| = q^{N-i} \ge 1$, so the term is at most
$-\tfrac12$. Let $g \ge 0$ be a smooth function with $\int g = 1$ supported in the ball of radius
$\varepsilon := 1/(2|\omega|)$ about the origin; it is a Schwartz function, so it is in $W^{2,1}$,
and $\|g\|_1 = 1$. Since $|e^{i\theta} - 1| \le |\theta|$,
$|\widehat g(\omega) - 1| \le \int g(x)\,|\langle\omega, x\rangle|\,dx \le \varepsilon|\omega| = \tfrac12$,
so $|\widehat g(\omega)| \ge \tfrac12$. Hence
$B \ge \|\mathcal{H}_tg\|_1 \ge |\widehat{\mathcal{H}_tg}(\omega)| \ge \frac2t\cdot\frac N2\cdot\frac12
= \frac{N}{2t}$ for every $N$, a contradiction.

For the generator let $f \in \mathcal{S}(\RR^d)$ and $v(\tau, \cdot) := \mu_\tau * f$. A Schwartz
function is in $W^{2,1}$, its derivatives being its weak derivatives by integration by parts, so
$v(t, \cdot) \in W^{2,1}$ by the passage of weak derivatives through $\mu_t$ above. Write
$\mathcal{F}$ for the transform of the $2\pi$ convention, $\mathcal{F}f(\xi) := \int e^{-2\pi
i\langle\xi, x\rangle}f(x)\,dx = \widehat f(-2\pi\xi)$, and $s := 2\pi|\xi|$. Since
$\int e^{-2\pi i\langle\xi, y\rangle}\mu_\tau(dy) = \widehat{\mu_\tau}(-2\pi\xi) = e^{-F_q(\tau s)}$,
Fourier inversion on the Schwartz space and Fubini's theorem give, for $\tau \ge 0$ and every $x$,
\[
  v(\tau, x) = \int e^{2\pi i\langle\xi, x\rangle}e^{-F_q(\tau s)}\,
    \mathcal{F}f(\xi)\,d\xi .
\]
By Theorem~\ref{thm:matern-pyramid-line}(3), $\partial_\tau e^{-F_q(\tau s)} = h_\tau(s)e^{-F_q(\tau s)}$
with $h_\tau(s) = \frac2\tau\sum_{i \ge 1}\bigl((1 + c^2\tau^2q^{-2i}s^2)^{-1} - 1\bigr)$, and
$|h_\tau(s)| \le 2\tau s^2\sum_ic^2q^{-2i}$, since $|(1 + a)^{-1} - 1| \le a$ for $a \ge 0$. So for
$\tau \in (t/2, 2t)$ the derivative of the integrand in $\tau$ is at most $A|\xi|^2|\mathcal{F}f(\xi)|$
in modulus, with $A$ a constant, an integrable bound, and differentiation under the integral sign
gives, at every $x$,
\[
  \partial_tv(t, x) = \int e^{2\pi i\langle\xi, x\rangle}h_t(s)e^{-F_q(ts)}\,
    \mathcal{F}f(\xi)\,d\xi .
\]
On the other side, the same computation with $L_{r_i}$ in place of $\mu_\tau$, applied to the
function $v(t, \cdot)$ written as above, gives
\[
  (L_{r_i} * v(t, \cdot))(x) - v(t, x)
   = \int e^{2\pi i\langle\xi, x\rangle}\Bigl(\frac{1}{1 + r_i^2s^2} - 1\Bigr)
     e^{-F_q(ts)}\,\mathcal{F}f(\xi)\,d\xi .
\]
These integrands are at most $r_i^2s^2|\mathcal{F}f(\xi)|$ in modulus, a bound summable in $i$ and
integrable in $\xi$, so the sum over $i$ may be taken under the integral sign, and $\frac2t$ times the
sum of the factors $(1 + r_i^2s^2)^{-1} - 1$ is $h_t(s)$. So
$\mathcal{H}_tv(t, \cdot)(x) = \partial_tv(t, x)$ at every $x$, and in particular almost everywhere.

\emph{(5)} For fixed $r$, $g_r(u) := \log(1 + c^2r^2e^{-2u})$ is decreasing and integrable on
$[0, \infty)$, and $\ell F_q(r) = \ell\sum_{i \ge 1}g_r(i\ell)$ is a Riemann sum of it. As in
Theorem~\ref{thm:matern-pyramid-line}(4), a statement in $r$ alone, $\ell F_q(r) \to \int_0^\infty
g_r(u)\,du = -\tfrac12\operatorname{Li}_2(-c^2r^2) = \Lambda(r)$, locally uniformly in $r$. For the
profile, Lemma~\ref{lem:bessel-profile}(2) with $\beta = c^2e^{-2u}$ and Tonelli's theorem in $u$
give $\Lambda(|\omega|) = \int(1 - \cos\langle\omega, x\rangle)\,k(|x|)\,
\bigl(|S^{d-1}|\,|x|^d\bigr)^{-1}dx$ with $k(r) := \int_0^\infty b_d(re^u/c)\,du$, nonincreasing.
% CHANGED (wave 3b, 2026-11-15)
% The integrability conditions follow from the bound k(r) <= 2C_d c/r, as for k_d in (1); the old
% text integrated b_d(y) y and b_d(y) log(cy)/y.
Its integrability conditions hold: $b_d(y) \le 2C_d/y$ gives
$k(r) \le \int_0^\infty 2C_dc\,e^{-u}\,du/r = 2C_dc/r$, and the argument of (1) applies. So
$\Lambda \in \mathcal{A}_d$ by
\ledger{A35}. A positive multiple of an origin datum
is an origin datum, with the multiple of its exponent, so $\ell F_q \in \mathcal{A}_d$, by the same
proposition, and the law with exponent $\ell F_q(t|\cdot|)$ exists by the definition of
$\mathcal{A}_d$. Its transform converges at every $\omega$ to $e^{-\Lambda(t|\omega|)}$, the transform
of a law, so the laws converge weakly by Lévy's continuity theorem. For $r > 0$,
$F_q(r) = \ell F_q(r)/\ell \to \infty$, since $\Lambda(r) > 0$ and $\ell \to 0$. So
$\widehat{\mu_t}$ tends to the indicator of $\{0\}$.
% CHANGED (wave 3b, 2026-11-15)
% The checked route reduces the last sentence to the line through the marginals of (3); the old
% text read Lévy's continuity theorem in R^d.
The kernels have no weak limit: a weak limit $\nu$ in $\RR^d$ would give, testing on
$\phi(\langle e, \cdot\rangle)$ for a unit vector $e$ and $\phi$ bounded and continuous on $\RR$,
the weak limit $\langle e, \cdot\rangle_*\nu$ of the kernels $\langle e, \cdot\rangle_*\mu_t$ of
the line's \Matern{} pyramids (clause (3)), which have none
(Theorem~\ref{thm:matern-pyramid-line}(4)).
% CHANGED (blueprint round, 2026-11-16, review flag "is itself proved from the", its similar
% places): the closing sentence "Every step is machine-checked, and the cited facts enter only
% through the statements the proof cites." is dropped; the proof names its citations where it reads
% them, and the provenance is the trust-base subsection's. In (3) the two counting results are cited
% with their hypotheses checked, one sentence each, and the page of A23's derivation is dropped
% (review flag "derivation from clause (1) is on p.~39"). The functions' transform is that of
% stage-lemma's proof, the 2 pi one is \mathcal{F} (review flag on the sign). No step changed.
\end{proof}

% draft: the comment after Theorem 4.1, line 551.
At $d = 1$ this is Theorem~\ref{thm:matern-pyramid-line}, and the comparison shows what the plane
changes. The telescoping $F_q(qr) - F_q(r) = \log(1 + c^2r^2)$ is a statement in $r$ alone, so one
octave is one Bessel potential in every dimension. The polar profile is nonincreasing and satisfies
the two integrability conditions, which is all that full covariance needs. On the line the profile
is completely monotone, and that is what makes the \Matern{} pyramid a Thorin family there
(Theorem~\ref{thm:matern-pyramid-line}(1)). For $d \ge 2$ the profile is not completely monotone,
one term or the whole sum, for every $q$ (Proposition~\ref{prop:thorin-in-dimension-d}), and that
property is lost. The counting of Proposition~\ref{prop:ladder-non-creation-line} survives on the
projections. The image itself gains local maxima: at every rung, some input does
(Proposition~\ref{prop:creation-plane}). The pyramid also has a twin in time. In the variable
$\sigma = r^2/2$ of the subordination bridge \prog{V-cone v0.1}{\S 4.2, p.~26},
$F_q(r) = \sum_{i \ge 1}\log(1 + 2c^2q^{-2i}\sigma)$. That is the exponent, in the Laplace variable
$\sigma$, of Lindeberg's time-causal limit kernel with ratio $q^2$ and constant $2c^2$, a cascade of
first-order integrators in time. Lindeberg defines that kernel by the infinite product of
first-order factors [\sscite{lindeberg2016time}, Sect.~5, (38), p.~62]; the ratio and the constant
are this reading of his product, and he realizes the kernel by recursive first-order filters
[\sscite{lindeberg2016time}, Sect.~6.2, p.~64].
% Anchored 2026-11-12 by the librarian (records/pyramids/CITATIONS-acquired.md, item 1, verdict
% SUPPORTS-NARROWER: the page prints the product, not the ratio-and-constant form).
The exponent of that kernel has the poles $\sigma = -q^{2i}/(2c^2)$, and $F_q$, as a function of
$r$, has the poles $\pm iq^i/c$ of its terms. The bridge's map $\theta \mapsto \sqrt{2\theta}$
carries the moduli $q^{2i}/(2c^2)$ of the first to the moduli $q^i/c$ of the second, which fixes the
constants of the twin. The selection by phase entropy on the line rests on a
theory of the line (Theorem~\ref{thm:phase-entropy-selection}), and
Theorem~\ref{thm:selection-stratum-two} replaces it. The projection of clause (3) is the marginal of
Lemma~\ref{lem:reduction-keeps-stratum}, the Radon projection of tomography.

The proofs of Propositions~\ref{prop:matern-increments} and~\ref{prop:thorin-in-dimension-d} below
use three facts on completely monotone functions.

% CHANGED (wave 3b, 2026-11-15, a strengthening): new lemma (new label lem:cm-log-convex), not in the
% draft. It and lem:cm-exponential-rate carry the one consequence of Bernstein's theorem (ledger
% A55) that the proof of prop:thorin-in-dimension-d reads, proved without it and without a measure
% (records/pyramids/PRICE-A55.md, Route A).
% CHANGED (fifth blind pass, 2026-11-15, flag 11): the lemma uses def:complete-bernstein, whose
% subsection defines complete monotonicity (its Lean tag names IsCompletelyMonotone); it had no
% uses.
% CHANGED (blueprint round, 2026-11-16): the three lemmas moved here from after
% prop:thorin-in-dimension-d, ahead of the propositions that use them; a pure move.
\begin{lemma}[complete monotonicity along an arithmetic progression]
\label{lem:cm-log-convex}
\lean{Pyramids.completelyMonotone_sub_shift, Pyramids.completelyMonotone_sq_le_mul}
\leanok
\uses{def:complete-bernstein}
Let $f$ be completely monotone on $(0, \infty)$ and $h \ge 0$.
\begin{enumerate}
\item $f - f(\cdot + h)$ is completely monotone on $(0, \infty)$.
\item $f(x + h)^2 \le f(x)f(x + 2h)$ for every $x > 0$.
\end{enumerate}
% CHANGED (fidelity review, 2026-11-15, R26/R50/R69): the status line lists the declarations' hypotheses beside the node's (LINKAGE rule 6).
\statusT\quad\emph{(Cites nothing; no Laplace representation of $f$ is used. Machine-checked.
\textbf{Hypotheses} (the node's, each with what its declarations require): $f$ completely monotone on
$(0, \infty)$: same, \texttt{IsCompletelyMonotone f}, smoothness and the signs of all derivatives on
$(0, \infty)$, for a real function $f$ on $\RR$ read on $(0, \infty)$ only; $h \ge 0$: same; (2)
$x > 0$: same.)}
\end{lemma}

\begin{proof}
\leanok
\emph{(1)} The $n$-th derivative of $f - f(\cdot + h)$ is $f^{(n)} - f^{(n)}(\cdot + h)$. The function
$(-1)^nf^{(n)}$ is nonincreasing on $(0, \infty)$, its derivative being
$(-1)^nf^{(n+1)} = -(-1)^{n+1}f^{(n+1)} \le 0$. So $(-1)^n\bigl(f^{(n)}(x) - f^{(n)}(x + h)\bigr) \ge 0$.

\emph{(2)} Fix $x > 0$. For a function $g$ let $L_g$ be the linear functional on $\RR[X]$ with
$L_g(X^k) = g(x + kh)$. Then $L_{g - g(\cdot + h)}(p) = L_g(p) - L_g(Xp)$, so
$L_g\bigl(X^k(1 - X)^{m+1}\bigr) = L_{g - g(\cdot + h)}\bigl(X^k(1 - X)^m\bigr)$. By (1) and induction on
$m$, simultaneously for every completely monotone $g$, $L_f\bigl(X^k(1 - X)^m\bigr) \ge 0$ for all
$k, m \ge 0$; the case $m = 0$ is $f(x + kh) \ge 0$. So the weights $w_k := L_f(b_{N,k})$ of the
Bernstein polynomials $b_{N,k} := \binom NkX^k(1 - X)^{N-k}$ are nonnegative. Applying $L_f$ to
$\sum_kb_{N,k} = 1$, $\sum_kk\,b_{N,k} = NX$ and $\sum_kk(k - 1)\,b_{N,k} = N(N - 1)X^2$ gives
$\sum_kw_k = f(x)$, $\sum_kk\,w_k = Nf(x + h)$ and $\sum_kk^2w_k = N(N - 1)f(x + 2h) + Nf(x + h)$. By
Cauchy--Schwarz, $\bigl(\sum_kk\,w_k\bigr)^2 \le \sum_kw_k\sum_kk^2w_k$, which after division by $N$
reads
\[
  N\bigl(f(x + h)^2 - f(x)f(x + 2h)\bigr) \le f(x)\bigl(f(x + h) - f(x + 2h)\bigr)
\]
for every $N \ge 1$. So $f(x + h)^2 - f(x)f(x + 2h) \le 0$.
\end{proof}

% CHANGED (wave 3b, 2026-11-15, a strengthening): new lemma (new label lem:cm-exponential-rate),
% not in the draft; see lem:cm-log-convex.
\begin{lemma}[the exponential rate of a completely monotone function]
\label{lem:cm-exponential-rate}
\lean{Pyramids.completelyMonotone_antitoneOn_exp_mul, Pyramids.not_isCompletelyMonotone_of_growth}
\leanok
\uses{lem:cm-log-convex}
Let $f > 0$ be completely monotone on $(0, \infty)$, and let $\lambda, C, R$ be real and $p \ge 0$ an
integer with $e^{\lambda r}f(r) \le C(1 + r)^p$ for $r \ge R$. Then $e^{\lambda r}f(r)$ is
nonincreasing on $(0, \infty)$. In particular, a function $f > 0$ on $(0, \infty)$ with such a bound
and with $e^{\lambda r}f(r)$ unbounded on $[1, \infty)$ is not completely monotone.
% CHANGED (fidelity review, 2026-11-15, R26/R50/R69): the status line lists the declarations' hypotheses beside the node's (LINKAGE rule 6).
\statusT\quad\emph{(Cites nothing. Machine-checked. \textbf{Hypotheses} (the node's, each with
what its declarations require): $f > 0$ on $(0, \infty)$: same in both; $f$ completely monotone: same
in \texttt{completelyMonotone\_antitoneOn\_exp\_mul}, and the negated conclusion of
\texttt{not\_isCompletelyMonotone\_of\_growth}; $\lambda, C, R$ real and $p \ge 0$ an integer: same,
$p \in \NN$; the bound $e^{\lambda r}f(r) \le C(1 + r)^p$ for $r \ge R$: same in both; $e^{\lambda
r}f(r)$ unbounded on $[1, \infty)$: same, above every real $M$ at some $r \ge 1$.)}
\end{lemma}

\begin{proof}
\leanok
\uses{lem:cm-log-convex}
Let $0 < x < y$, $h := y - x$ and $a_n := e^{\lambda(x + nh)}f(x + nh) > 0$.
Lemma~\ref{lem:cm-log-convex}(2) at $x + nh$, multiplied by
$e^{2\lambda(x + (n+1)h)} = e^{\lambda(x + nh)}e^{\lambda(x + (n+2)h)}$, gives
$a_{n+1}^2 \le a_na_{n+2}$. So the ratios $a_{n+1}/a_n$ are nondecreasing, and
$a_n \ge a_0\rho^n$ with $\rho := a_1/a_0$. Suppose $\rho > 1$. For $n \ge \max(1, R/h)$,
$a_0\rho^n \le a_n \le |C|(1 + x + nh)^p \le |C|(1 + x + h)^pn^p$, against $n^p\rho^{-n} \to 0$. So
$a_1 \le a_0$, that is $e^{\lambda y}f(y) \le e^{\lambda x}f(x)$. For the last sentence: if $f$ were
completely monotone, $e^{\lambda r}f(r) \le e^\lambda f(1)$ for every $r \ge 1$.
\end{proof}

% CHANGED (wave 3b, 2026-11-15, a strengthening): new lemma (new label lem:cm-uncancelled-atom), not
% in the draft. It carries the one step of prop:matern-increments(3) that read Bernstein's theorem
% (ledger A55), proved without it and without a measure (records/pyramids/PRICE-A55.md, A.1).
\begin{lemma}[an uncancelled atom]
\label{lem:cm-uncancelled-atom}
\lean{Pyramids.not_isCompletelyMonotone_expDiff}
\leanok
Let $0 < A < B$ and $\rho > 1$ with $A\rho^i \ne B\rho$ for every $i \ge 1$. Then
$g(x) := \sum_{i \ge 1}\bigl(e^{-A\rho^ix} - e^{-B\rho^ix}\bigr)$ is not completely monotone on
$(0, \infty)$.
% CHANGED (fidelity review, 2026-11-15, R26/R50/R69): the status line lists the declarations' hypotheses beside the node's (LINKAGE rule 6).
\statusT\quad\emph{(Cites nothing; no Laplace representation of $g$ is used. Machine-checked.
\textbf{Hypotheses} (the node's, each with what its declarations require): $0 < A < B$: same;
$\rho > 1$: same; $A\rho^i \ne B\rho$ for every $i \ge 1$: same, indexed from $0$ as
$A\rho^{i+1}$; the series $g$: the declaration writes it as the difference of the two series
$\sum_ie^{-A\rho^ix}$ and $\sum_ie^{-B\rho^ix}$, which both converge for $x > 0$, where complete
monotonicity is read: same.)}
\end{lemma}

\begin{proof}
\leanok
Put $a_i := A\rho^i$ and $b_i := B\rho^i$. For $\vartheta, x > 0$,
$\vartheta^ne^{-\vartheta x} \le (n + 2)!\,x^{-n-2}\vartheta^{-2}$, and $\sum_ia_i^{-2}$,
$\sum_ib_i^{-2}$ are geometric. So the series of $g$ and of its termwise derivatives converge
uniformly on $(x_0, \infty)$ for every $x_0 > 0$, and
\[
  (-1)^ng^{(n)}(x) = \sum_ia_i^ne^{-a_ix} - \sum_ib_i^ne^{-b_ix} .
\]
Let $x_* := b_1$ and $h(u) := ue^{1-u}$. At $x = n/x_*$, $n \ge 1$,
\[
  \Bigl(\frac{e}{x_*}\Bigr)^n(-1)^ng^{(n)}\Bigl(\frac{n}{x_*}\Bigr)
   = \sum_ih\Bigl(\frac{a_i}{x_*}\Bigr)^n - \sum_ih\Bigl(\frac{b_i}{x_*}\Bigr)^n .
\]
Since $e^{u-1} \ge u$, with equality only at $u = 1$, $0 \le h(u) \le 1$ for $u \ge 0$ and
$h(u) < 1$ for $u \ne 1$. By hypothesis $a_i \ne x_*$ for every $i$; $b_1/x_* = 1$ and
$b_i/x_* = \rho^{i-1} > 1$ for $i \ge 2$. So the terms of the first sum tend to $0$ as $n \to \infty$,
and those of the second to $1$ for $i = 1$ and to $0$ otherwise. For $n \ge 1$, $h(u)^n \le h(u)$,
and $\sum_ih(a_i/x_*)$, $\sum_ih(b_i/x_*)$ converge, $h(\vartheta) = e\,\vartheta e^{-\vartheta} \le
6e\,\vartheta^{-2}$. By dominated convergence the right side tends to $0 - 1 = -1$, so
$(-1)^ng^{(n)}(n/x_*) < 0$ for large $n$, and $g$ is not completely monotone.
\end{proof}

% draft: the paragraph before Proposition 4.2, line 561 (catalogue List D item 18; ledger A57, the
% librarian's anchor in records/pyramids/VERBATIM-Sato.md and VERBATIM-literature.md, section A57).
The octave step is a Bessel potential. The steps between other pairs of scales are not, and the
next proposition says what they are. A \emph{Bernstein function} is a function $f \ge 0$ on
$(0, \infty)$ whose derivative is completely monotone; equivalently
$f(u) = a + bu + \int_{(0,\infty)}(1 - e^{-u\tau})\,m(d\tau)$ with $a, b \ge 0$ and a measure $m$
with $\int(1 \wedge \tau)\,m(d\tau) < \infty$, its \emph{\Levy{} measure}
[\sscite{schilling2012bernstein}, Thm.~3.2, p.~21]. When $f(0{+}) = 0$, the function
$e^{-f(|\omega|^2)}$ is the transform of an infinitely divisible law in every dimension, by
subordination. Brownian motion on $\RR^d$ with the transform $e^{-|\omega|^2/2}$ at time $1$, run at
the random time of the subordinator with exponent $u \mapsto f(2u)$, which is again a Bernstein
function, has at time $1$ the transform $e^{-f(|\omega|^2)}$ (\ledger{A57}:
[\sscite{sato1999levy}, Thm.~30.1, pp.~197--198], and [\sscite{schilling2012bernstein}, Thm.~13.14]
with Example~13.13, p.~212). The function $f$ is a \emph{complete Bernstein function} if its
\Levy{} measure has a completely monotone density (Definition~\ref{def:complete-bernstein}), and a
\emph{Thorin--Bernstein function} if that density $m$ has $\tau m(\tau)$ completely monotone
(Definition~\ref{def:thorin-bernstein}).

% draft: Proposition 4.2 (section 4.1), lines 563-567 of draft/spatial-hemigroup-pyramids.md.
\begin{proposition}[the increments]
\label{prop:matern-increments}
\lean{Pyramids.matern_increments_1, Pyramids.matern_increments_2,
Pyramids.matern_increments_2_levy_exists, Pyramids.matern_increments_2_poisson,
Pyramids.matern_increments_3_levy_exists, Pyramids.matern_increments_3,
Pyramids.matern_increments_3_thorin, Pyramids.matern_increments_3_sum}
% CHANGED (wave 3b, 2026-11-15): the eight declarations are machine-checked, so statement and proof
% are marked proved.
\leanok
% CHANGED (wave 3b, 2026-11-15, a strengthening): thm:cited-A55 is no longer used; the one step of
% (3) that read Bernstein's theorem is lem:cm-uncancelled-atom, and (1)'s complete monotonicity
% reads lem:cm-log-convex(1).
% CHANGED (fidelity review, 2026-11-15, R66): \uses gains thm:matern-pyramid-line, whose octave
% identity F_q(q x) - F_q(x) = log(1 + c^2 x^2) the telescoping of (3) reads, as its declarations do;
% the proof said "as in Theorem matern-pyramid(1)".
\uses{def:matern-pyramid-data, lem:bessel-profile, thm:matern-pyramid, thm:matern-pyramid-line,
def:complete-bernstein, def:thorin-bernstein, thm:cited-A39, lem:cm-log-convex,
lem:cm-uncancelled-atom}
Let $\beta > \gamma > 0$ and $f_{\beta,\gamma}(u) := \log\frac{1 + \beta u}{1 + \gamma u}$, $u \ge 0$.
\begin{enumerate}
\item $f_{\beta,\gamma}(u) = \int_{1/\beta}^{1/\gamma}\frac{u}{u + x}\,\frac{dx}{x}
= \int_0^\infty(1 - e^{-\tau u})\,\frac{e^{-\tau/\beta} - e^{-\tau/\gamma}}{\tau}\,d\tau$. It is a
bounded complete Bernstein function and not a Thorin--Bernstein function.
\item In every dimension, the law with transform $e^{-f_{\beta,\gamma}(|\omega|^2)}$ is
$\frac\gamma\beta\delta_0 + \bigl(1 - \frac\gamma\beta\bigr)L_{\sqrt\beta}$. It is compound Poisson with
\Levy{} measure $2\int_{\sqrt\gamma}^{\sqrt\beta}L_r\,dr/r$, of total mass $\log(\beta/\gamma)$.
\item So every step $\mu_{s,t}$, $0 < s < t$, of the \Matern{} pyramid is infinitely divisible in
every dimension, with \Levy{} measure
$\nu_{s,t} = \sum_{i \ge 1}2\int_{csq^{-i}}^{ctq^{-i}}L_r\,dr/r$. As a function of $u = |\omega|^2$,
its exponent is a Thorin--Bernstein function if and only if $t/s \in q^{\NN}$, and then, with
$t = q^ks$, it is $\sum_{j=0}^{k-1}\log(1 + c^2q^{2j}s^2u)$.
\end{enumerate}
% CHANGED (wave 3b, 2026-11-15, a strengthening): the node no longer cites A55. The status line
% read: "[A] A39 A55 (Assignment. A39, Sato Thm. 8.1(ii)-(iii), pp. 37-38, carries two steps: the
% uniqueness of the generating triplet, which makes the measures of (2) and (3) the Levy measures of
% the laws, and the existence of the law with triplet (0, nu_{s,t}/n, 0), which makes mu_{s,t}
% infinitely divisible in (3). A55, Bernstein's theorem (schilling2012bernstein, Thm. 1.4, p. 3),
% carries one step of (3): a completely monotone tau m(tau) would be the Laplace transform of a
% positive measure. Proved here: the two integral forms, complete monotonicity of the density, the
% mixture of (2), the compound Poisson form, the sum of (3) and the atom count. Not carried: nothing
% about the laws beyond their transforms and Levy measures.)"
% CHANGED (fidelity review, 2026-11-15, R26/R50/R69): the status line lists the declarations' hypotheses beside the node's (LINKAGE rule 6).
% CHANGED (fidelity review, 2026-11-15, R45): the status line says that the declarations state a
% Levy measure of the law, through HasTriplet, and that "the" rests on A39's uniqueness.
\statusA\quad\ledger{A39}\quad\emph{(\textbf{Assignment.} \ledger{A39},
[\sscite{sato1999levy}, Thm.~8.1(ii)--(iii), pp.~37--38], carries two steps: the uniqueness of the
generating triplet, which makes the measures of (2) and (3) \emph{the} \Levy{} measures of the
laws, and the existence of a law with triplet $(0, \nu_{s,t}, 0)$, which makes $\mu_{s,t}$
infinitely divisible in (3). Proved here: the two integral forms, the complete monotonicity of the
density and of $f_{\beta,\gamma}'$, the mixture of (2), the compound Poisson form, the generating
triplets of (2) and (3), the sum of (3), and the failure of the Thorin property when
$t/s \notin q^{\NN}$, by Lemma~\ref{lem:cm-uncancelled-atom}, without Bernstein's theorem. Two
uniqueness facts are read, both proved: the uniqueness of the \Levy{} measure of a Bernstein
function, in (1) and (3) (the affine module's \texttt{bernstein\_triplet\_unique}), and, under it,
the uniqueness of the Laplace transform of a locally finite measure on $(0, \infty)$ (the shared
core's \texttt{laplace\_uniqueness\_locally\_finite}); clause (2) also reads the uniqueness of the
transform of a finite measure.
% Entry A57 (the subordination of the paragraph above) is not spent: clause (2) writes the law down.
\textbf{Not carried}: nothing about the laws beyond their transforms and \Levy{} measures.
Machine-checked. \textbf{The article.} ``The \Levy{} measure'' of (2) and (3) is stated in the
declarations as $\exists\,T$, \texttt{HasTriplet} of the law and $T$, with $T.\nu = \nu$: a \Levy{}
measure of the law. That it is \emph{the} one rests on the uniqueness of the generating triplet,
\ledger{A39} as above, which this line credits and the declarations do not carry.
\textbf{Hypotheses} (the node's, each with what its declarations require): $\beta > \gamma > 0$: same
in \texttt{matern\_increments\_1}, \texttt{\_2}, \texttt{\_2\_poisson}; weaker in
\texttt{\_2\_levy\_exists}, which requires nothing of $\beta, \gamma$ and only builds the measure;
$u \ge 0$: same; (2) every dimension $d \ge 1$: same in \texttt{\_2\_poisson}; weaker in
\texttt{\_2} and \texttt{\_2\_levy\_exists}, which do not require $d \ge 1$; the \Levy{} measure
$2\int_{\sqrt\gamma}^{\sqrt\beta}L_r\,dr/r$: same, every $\nu$ with that value on each measurable set,
one of which \texttt{\_2\_levy\_exists} gives; (3) the \Matern{} pyramid, $q > 1$, $c > 0$,
$d \ge 1$: same in \texttt{\_3}, over every $\mu$ with \texttt{IsMaternPyramid q c}; in
\texttt{\_3\_thorin} and \texttt{\_3\_sum} $q > 1$ and $c > 0$ same and $d$ absent, as they read the
exponent only; \texttt{\_3\_levy\_exists} requires nothing; $0 < s < t$: same in \texttt{\_3} and
\texttt{\_3\_thorin}; weaker in \texttt{\_3\_sum}, which holds for every real $s$, every $k \in \NN$
and every $u \ge 0$; $\nu_{s,t}$: same, every $\nu$ with the value of the series on each measurable
set, one of which \texttt{\_3\_levy\_exists} gives.)}
\end{proposition}

% CHANGED (wave 3b, 2026-11-15, a strengthening): the proof follows the machine-checked route. The
% old proof of (3) applied Bernstein's theorem (A55) to tau m(tau): a positive eta on [0, infinity)
% with tau m(tau) = int e^{-tau x} eta(dx), so that eta + sigma^- and sigma^+, the parts of the signed
% atomic measure sigma = sum_i (delta_{1/beta_i} - delta_{1/gamma_i}), have one Laplace transform and
% are equal by its uniqueness, while sigma^- charges 1/gamma_1, where sigma^+ has no atom. It now
% reads Lemma cm-uncancelled-atom. In (1) the old proof left the complete monotonicity of f' implicit
% (now Lemma cm-log-convex(1)), read the uniqueness of the Levy measure from f'(u) = b + int tau
% e^{-u tau} m(d tau) (now the second difference, which needs no derivative), and called tau m(tau)
% "not nonincreasing" from its limit 0 at 0 (now the bound tau m(tau) <= tau/gamma). In (2) the
% r-integral was the substitution x = r^{-2} and (1); it is now the fundamental theorem of calculus.
% In (3) the second moment read the covariance 2r^2 I of L_r; it now reads that of L_1. The infinite
% divisibility read Thm. 8.1(iii) at every triplet (0, nu_{s,t}/n, 0) and took n-th convolution
% powers; it now reads it once, at (0, nu_{s,t}, 0).
% CHANGED (fifth blind pass, 2026-11-15, flags 7-9): in (2) the compound Poisson law is shown
% infinitely divisible before Sato's Thm. 8.1(ii) is applied to it (its n-th root is the compound
% Poisson law of nu/n); in (3) the continuity of g is proved by the M-test, where it was credited to
% lem:cm-uncancelled-atom, whose statement does not say it; in (1) the forward pointer to the proof
% of prop:thorin-in-dimension-d(2) is dropped.
\begin{proof}
\leanok
\uses{def:matern-pyramid-data, lem:bessel-profile, thm:matern-pyramid, thm:matern-pyramid-line,
def:complete-bernstein, def:thorin-bernstein, thm:cited-A39, lem:cm-log-convex,
lem:cm-uncancelled-atom}
\emph{(1)} Partial fractions give $\frac{u}{x(u + x)} = \frac1x - \frac1{u + x}$, whose integral
from $1/\beta$ to $1/\gamma$ is $\log\frac{1/\gamma}{u + 1/\gamma} - \log\frac{1/\beta}{u + 1/\beta}
= f_{\beta,\gamma}(u)$. The second form is the difference of two Frullani integrals,
$\int_0^\infty(1 - e^{-\tau u})e^{-\tau/\beta}\,d\tau/\tau = \log(1 + \beta u)$ and the same for
$\gamma$, those in the proof of Lemma~\ref{lem:bessel-profile}(2). Its density
\[
  m_{\beta,\gamma}(\tau) := \frac{e^{-\tau/\beta} - e^{-\tau/\gamma}}{\tau}
   = \int_{1/\beta}^{1/\gamma}e^{-\tau x}\,dx
   = \frac1\beta\int_1^{\beta/\gamma}e^{-(\tau/\beta)s}\,ds
\]
is positive, and $(1 \wedge \tau)\,m_{\beta,\gamma}(\tau) \le \tau m_{\beta,\gamma}(\tau) \le
e^{-\tau/\beta}$, so $\int(1 \wedge \tau)\,m_{\beta,\gamma}(\tau)\,d\tau$ is finite. It is completely
monotone. A Laplace transform $y \mapsto \int_1^\infty e^{-ys}w(s)\,ds$ of a bounded $w \ge 0$ with
bounded support may be differentiated under the integral, and $(-1)^n$ times its $n$-th derivative is
$\int_1^\infty s^ne^{-ys}w(s)\,ds \ge 0$. The substitution $\tau \mapsto \tau/\beta$
multiplies the $n$-th derivative by $\beta^{-n}$. On $(0, \infty)$, $f_{\beta,\gamma} \ge 0$ and
$f_{\beta,\gamma}'(u) = (u + 1/\beta)^{-1} - (u + 1/\gamma)^{-1} = g(u) - g(u + 1/\gamma - 1/\beta)$
with $g(u) := (u + 1/\beta)^{-1}$, whose $n$-th derivative is $(-1)^nn!\,(u + 1/\beta)^{-n-1}$. So
$f_{\beta,\gamma}'$ is completely monotone by Lemma~\ref{lem:cm-log-convex}(1), and
$f_{\beta,\gamma}$ is a complete Bernstein function, with $a = b = 0$ and the \Levy{} measure
$m_{\beta,\gamma}(\tau)\,d\tau$ of the Bernstein function. It is bounded by its limit
$\log(\beta/\gamma)$ at infinity.

The density of a Bernstein function is unique among densities continuous on $(0, \infty)$. Let
$a + bu + \int(1 - e^{-u\tau})\,m(\tau)\,d\tau$ and $a' + b'u + \int(1 - e^{-u\tau})\,m'(\tau)\,d\tau$
agree and be finite for every $u > 0$, with $a, b, a', b' \ge 0$ and $m, m' \ge 0$. Pointwise,
$(1 - e^{-u\tau}) + (1 - e^{-(u+2)\tau}) + e^{-u\tau}(1 - e^{-\tau})^2 = 2(1 - e^{-(u+1)\tau})$, so the
second difference in steps of $1$ removes $a + bu$. The measures $(1 - e^{-\tau})^2m(\tau)\,d\tau$ and
$(1 - e^{-\tau})^2m'(\tau)\,d\tau$ on $(0, \infty)$ then have the same finite Laplace transform at
every $u > 0$, and they are equal by the uniqueness of the Laplace transform of a locally finite
measure. As $(1 - e^{-\tau})^2 > 0$ there, $m = m'$ almost everywhere on $(0, \infty)$, and
everywhere if both are continuous. Now let $f_{\beta,\gamma}$ have a Thorin representation, with
density $m'$. Then $\tau m'(\tau)$ is completely monotone, hence continuous and nonnegative, so
$m' = m_{\beta,\gamma}$ on $(0, \infty)$, and $h(\tau) := \tau m_{\beta,\gamma}(\tau) =
e^{-\tau/\beta} - e^{-\tau/\gamma}$ would be nonincreasing. But $h(1) > 0$, and
$h(\tau) \le 1 - (1 - \tau/\gamma) = \tau/\gamma$, so
$\tau_1 := \min\bigl(\tfrac12, \tfrac12\gamma h(1)\bigr) < 1$ has $h(\tau_1) \le \tfrac12h(1) < h(1)$.
So $f_{\beta,\gamma}$ is not a Thorin--Bernstein function.

\emph{(2)} $\frac{1 + \gamma u}{1 + \beta u} = \frac\gamma\beta + \bigl(1 - \frac\gamma\beta\bigr)
\frac{1}{1 + \beta u}$, and $\frac{1}{1 + \beta|\omega|^2}$ is the transform of $L_{\sqrt\beta}$
(Lemma~\ref{lem:bessel-profile}(1)). That gives the mixture, by the uniqueness of the transform.
Let $\nu := 2\int_{\sqrt\gamma}^{\sqrt\beta}L_r\,dr/r$, the image of $(2/r)\,dr$ on
$(\sqrt\gamma, \sqrt\beta]$ times the law of $\sqrt{2X}\,Z$ under $(r, x) \mapsto rx$. It is symmetric,
each $L_r$ having an even transform; $\nu(\{0\}) = 0$, each $L_r$ having a density; and its mass is
$2\log\sqrt{\beta/\gamma} = \log(\beta/\gamma)$. Its compound Poisson law
$e^{-\nu(\RR^d)}\sum_{n \ge 0}\nu^{*n}/n!$ has the transform
$\exp\bigl(-\int(1 - \cos\langle\omega, x\rangle)\,\nu(dx)\bigr)$, $\nu$ being symmetric. By
Tonelli's theorem, with $\int(1 - \cos\langle\omega, x\rangle)\,L_r(dx) = 1 - (1 + r^2|\omega|^2)^{-1}$,
\[
  \int\bigl(1 - \cos\langle\omega, x\rangle\bigr)\,\nu(dx)
   = \int_{\sqrt\gamma}^{\sqrt\beta}\frac{2r|\omega|^2}{1 + r^2|\omega|^2}\,dr
   = \log(1 + \beta|\omega|^2) - \log(1 + \gamma|\omega|^2) = f_{\beta,\gamma}(|\omega|^2) ,
\]
the integrand being the derivative of $r \mapsto \log(1 + r^2|\omega|^2)$. So the law of (2) is that
compound Poisson law. It is infinitely divisible: for every integer $n \ge 1$ the compound Poisson
law of $\nu/n$ has the transform $\exp\bigl(-\tfrac1n\int(1 - \cos\langle\omega, x\rangle)\,
\nu(dx)\bigr)$, so its $n$-th convolution power has the transform of the law of (2) and is that
law, by the uniqueness of the transform. Its transform is the exponential of the exponent of the triplet $(0, \nu, 0)$,
the compensator term integrating to $0$ by symmetry, so $(0, \nu, 0)$ is a generating triplet. By
its uniqueness, \ledger{A39}, $\nu$ is the \Levy{}
measure.

\emph{(3)} Put $\beta_i := c^2q^{-2i}t^2$ and $\gamma_i := c^2q^{-2i}s^2$, $i \ge 1$. The exponent of
$\mu_{s,t}$ is $F_q(t|\omega|) - F_q(s|\omega|) = \sum_if_{\beta_i,\gamma_i}(|\omega|^2)$, a
difference of two convergent sums since $\log(1 + \beta_iu) \le \beta_iu$. Let $\nu_{s,t}$ be the
sum of the measures of (2) for $\beta_i, \gamma_i$, with $\sqrt{\beta_i} = ctq^{-i}$ and
$\sqrt{\gamma_i} = csq^{-i}$. By (2) and Tonelli's theorem,
$\int(1 - \cos\langle\omega, x\rangle)\,\nu_{s,t}(dx) = \sum_if_{\beta_i,\gamma_i}(|\omega|^2)$. It is
symmetric and does not charge $0$. Its second moment is finite: that of $L_r$ is $r^2$ times that of
$L_1$, the law of $\sqrt{2X}\,Z$, which is $2E[X]\,E|Z|^2 < \infty$; so the $i$-th term has a second
moment at most $\int_{\sqrt{\gamma_i}}^{\sqrt{\beta_i}}2r\,dr \le \beta_i$ times that of $L_1$, and
$\sum_i\beta_i < \infty$. So $\nu_{s,t}$ is a symmetric \Levy{} measure, and $(0, \nu_{s,t}, 0)$ is a
generating triplet of $\mu_{s,t}$, the compensator term vanishing by symmetry. By
\ledger{A39}, an infinitely divisible law has this
triplet. It has the transform of $\mu_{s,t}$, so it is $\mu_{s,t}$, which is therefore infinitely
divisible; its \Levy{} measure is $\nu_{s,t}$ by the uniqueness of the triplet, Thm.~8.1(ii) of the
same page. This holds in every dimension.

As a function of $u$ the exponent is $\psi(u) := \sum_if_{\beta_i,\gamma_i}(u) =
\int_0^\infty(1 - e^{-\tau u})\,m(\tau)\,d\tau$ with $m := \sum_im_{\beta_i,\gamma_i}$, by (1) and
Tonelli's theorem. Here $\tau m(\tau) = g(\tau) := \sum_{i \ge 1}\bigl(e^{-\tau/\beta_i} -
e^{-\tau/\gamma_i}\bigr)$, which is continuous on $(0, \infty)$: for $\tau \ge \tau_0 > 0$ the $i$-th
term lies in $[0, e^{-\tau_0/\beta_i}]$, and $e^{-\tau_0/\beta_i} \le \beta_i/\tau_0$, as
$e^{-y} \le 1/y$ for $y > 0$, with $\sum_i\beta_i < \infty$; so the series of continuous terms
converges uniformly on $[\tau_0, \infty)$. So $m$ is continuous there.
Suppose $\psi$ were a Thorin--Bernstein function, with density $m'$. As in (1), $m' = m$ on
$(0, \infty)$, so $g$ would be completely monotone. Here $1/\beta_i = A\rho^i$ and
$1/\gamma_i = B\rho^i$, with $\rho := q^2 > 1$ and $A := (ct)^{-2} < B := (cs)^{-2}$, and
$A\rho^i = B\rho$ means $(t/s)^2 = q^{2(i-1)}$, that is $t = q^{i-1}s$. If $t/s \notin q^{\NN}$, no
$A\rho^i$ equals $B\rho$, and Lemma~\ref{lem:cm-uncancelled-atom} shows that $g$ is not completely
monotone. So $\psi$ is not a Thorin--Bernstein function. If $t = q^ks$, $k \ge 1$, then telescoping
with the octave identity $F_q(q\omega) - F_q(\omega) = \log(1 + c^2\omega^2)$ of
Theorem~\ref{thm:matern-pyramid-line}(1) gives
$\psi(u) = \sum_{j=0}^{k-1}\bigl(F_q(q^{j+1}s\sqrt u) - F_q(q^js\sqrt u)\bigr) =
\sum_{j=0}^{k-1}\log(1 + c^2q^{2j}s^2u)$. Each term $\log(1 + \theta u)$ has, by Frullani's integral,
the density $e^{-\tau/\theta}/\tau$, with $\tau m(\tau) = e^{-\tau/\theta}$ completely monotone and
$(1 \wedge \tau)e^{-\tau/\theta}/\tau \le e^{-\tau/\theta}$, and its derivative
$(u + 1/\theta)^{-1}$ is completely monotone. Densities and derivatives add, so $\psi$ is a
Thorin--Bernstein function.
\end{proof}

% draft: the comment after Proposition 4.2, line 569.
Clause (2) describes one factor of the product in clause (3). That factor keeps the image
with probability $\gamma/\beta$ and otherwise applies one Bessel potential. At $t = qs$ the factors
telescope, and the octave step is the single Bessel potential $L_{cs}$, with no atom. The next
proposition says where the line's complete monotonicity of the profile goes in the plane.

% draft: Proposition 4.3 (section 4.1), lines 575-578 of draft/spatial-hemigroup-pyramids.md.
% CHANGED (fidelity review, 2026-11-15, R47): the tag adds thorin_dimension_2_profile_integrable,
% the identity of (2) with both integrands integrable (d >= 1), so it is not read as 0 = 0. The
% statement is unchanged.
\begin{proposition}[where the Thorin property goes in $d \ge 2$]
\label{prop:thorin-in-dimension-d}
\lean{Pyramids.thorin_dimension_1, Pyramids.thorin_dimension_1_line,
Pyramids.thorin_dimension_2_profile, Pyramids.thorin_dimension_2_low,
Pyramids.thorin_dimension_2_high, Pyramids.thorin_dimension_2_profile_integrable}
% CHANGED (wave 3b, 2026-11-15): the five declarations are machine-checked, so statement and proof
% are marked proved.
\leanok
% CHANGED (wave 3b, 2026-11-15, a strengthening): thm:cited-A55 is no longer used; the two
% consequences of Bernstein's theorem the proof reads are lem:cm-log-convex and
% lem:cm-exponential-rate.
% CHANGED (fifth blind pass, 2026-11-15, flag 10): lem:cm-log-convex leaves the uses of statement
% and proof; the proof reaches it only through lem:cm-exponential-rate, which uses it.
\uses{def:matern-pyramid-data, lem:bessel-profile, lem:bessel-k, thm:matern-pyramid,
lem:laplace-atoms-admissible, lem:cm-exponential-rate}
Let $q > 1$ and $c > 0$, and let $k_d$ be the polar profile of Theorem~\ref{thm:matern-pyramid}(1).
\begin{enumerate}
\item For $d \ge 2$, neither $b_d$ nor $k_d$ is completely monotone on $(0, \infty)$. For $d = 1$
both are.
\item The limit profile of Theorem~\ref{thm:matern-pyramid}(5) is $P_d(r/c)$, where
$P_d(y) := \int_y^\infty b_d(v)\,dv/v$. The function $P_d$ is completely monotone for $d \le 3$, with
$P_1 = 2E_1$, $P_2 = 2K_0$ and $P_3 = 2E_1 + 2e^{-y}$, $E_1(y) := \int_y^\infty e^{-v}\,dv/v$ the
exponential integral. It is not completely monotone for $d \ge 4$.
\end{enumerate}
% CHANGED (wave 3b, 2026-11-15, a strengthening): the node no longer cites A55 and is [T]; it
% reaches A51 only through lem:bessel-k. The status line read: "[A] A55 (Assignment. A55,
% Bernstein's theorem (schilling2012bernstein, Thm. 1.4, p. 3), carries one step, used in (1) and
% in the case d >= 4 of (2): a completely monotone function on (0, infinity) is the Laplace
% transform of a positive measure on [0, infinity). The identification P_2 = 2K_0 and the integral
% of K_0 come from Lemma bessel-k, which carries entry A51; the rest is proved here: the two bounds
% on e^y b_d(y) and the comparison. Not carried: no class property of the kernels mu_t.)"
% CHANGED (fidelity review, 2026-11-15, R26/R50/R69): the status line lists the declarations' hypotheses beside the node's (LINKAGE rule 6).
\statusT\quad\emph{(Cites nothing of its own. Clause (1) and the case $d \ge 4$ of (2) read
Lemma~\ref{lem:cm-exponential-rate}, proved here without Bernstein's theorem, in place of a Laplace
representation. The identification $P_2 = 2K_0$ and the integral of $K_0$ come from
Lemma~\ref{lem:bessel-k}, which carries entry A51; the rest is proved here: the two bounds on
$e^yb_d(y)$ and the comparison. \textbf{Not carried}: no class property of the kernels $\mu_t$.
Machine-checked. \textbf{Hypotheses} (the node's, each with what its declarations require): $q > 1$,
$c > 0$: same in \texttt{thorin\_dimension\_1} and \texttt{\_1\_line}; $q$ absent from (2), which does
not read it, and $c > 0$ same in \texttt{\_2\_profile} and \texttt{\_2\_profile\_integrable}, absent
from \texttt{\_2\_low} and \texttt{\_2\_high}, which are about $P_d$ alone; (1) $d \ge 2$: same;
$d = 1$: same, the declaration stated at $d = 1$; (2) $d \ge 1$, the chapter's: same in
\texttt{\_2\_profile\_integrable}, weaker in \texttt{\_2\_profile}, which does not require it and states
the identity of the two integrals alone; $r > 0$, where the limit profile is read: same;
$y > 0$ for the closed forms: same; $d \le 3$: same, the declaration stated at $d = 1, 2, 3$;
$d \ge 4$: same; complete monotonicity on $(0, \infty)$: same, \texttt{IsCompletelyMonotone},
smoothness and the signs of all derivatives on $(0, \infty)$.)}
\end{proposition}

% CHANGED (wave 3b, 2026-11-15, a strengthening): the proof follows the machine-checked route. The
% old proof applied Bernstein's theorem (A55) to k_d, b_d and P_d: a representation
% k_d(r) = int e^{-rs} eta(ds) with eta positive, the bound eta([0, lambda')) e^{-r lambda'} <= k_d(r)
% <= M'(1 + r)^{d/2} e^{-lambda r} giving eta([0, lambda')) = 0 for lambda' < lambda, so that
% e^{lambda r} k_d(r) = int e^{-r(s - lambda)} eta(ds) is nonincreasing, against the lower bound.
% It now reads that monotonicity from Lemma cm-exponential-rate. The upper bound was printed as
% M_d(1 + y)^{d/2} from "(v - y/2)^2/v >= v/4 for v >= y", with the sum over j bounded by
% sum q^{jd/2} e^{-(q^j - 1)} at x >= 1; the checked bound is M_d(1 + y)^d, with the ratio of the
% geometric comparison made at most 1/2 by taking x large. For d >= 4 the old proof bounded
% e^y P_d(y) above and below directly; the checked proof passes to -P_d' = b_d(y)/y. The complete
% monotonicity for d = 1 and d <= 3 is now proved, not asserted.
\begin{proof}
\leanok
\uses{def:matern-pyramid-data, lem:bessel-profile, lem:bessel-k, thm:matern-pyramid,
lem:laplace-atoms-admissible, lem:cm-exponential-rate}
\emph{The two bounds.} Let $d \ge 2$, so $d/2 - 1 \ge 0$. \emph{Upper}: let $y > 0$. For
$0 < v \le y$, $v^{d/2-1} \le y^{d/2-1}$, $v + y^2/(4v) \ge y$ and $e^{1 - v/y} \ge 1$. For $v \ge y$,
$v + y^2/(4v) \ge y + v/4$, since $4v^2 + y^2 - 4vy - v^2 = (3v - y)(v - y) \ge 0$. So the integrand
of $b_d$ is at most $e^{-y}\bigl(e\,y^{d/2-1}e^{-v/y} + v^{d/2-1}e^{-v/4}\bigr)$, and integrating,
\[
  e^yb_d(y) \le \frac{2}{\Gamma(d/2)}\bigl(e\,y^{d/2} + 4^{d/2}\Gamma(d/2)\bigr)
  \le M_d(1 + y)^d , \qquad M_d := \frac{2}{\Gamma(d/2)}\bigl(e + 4^{d/2}\Gamma(d/2)\bigr) .
\]
\emph{Lower}: let $y \ge 1$. On $[y/2, y/2 + \sqrt y]$, $(2v - y)^2 \le 4y \le 8v$, so
$v + y^2/(4v) \le y + 2$, and $v^{d/2-1} \ge (y/2)^{d/2-1}$. So
$e^yb_d(y) \ge m_d(y/2)^{d/2-1}\sqrt y$ with $m_d := 2e^{-2}/\Gamma(d/2)$; for $y \ge 2$ this is at
least $m_d\sqrt y$. Also $b_d > 0$ on $[0, \infty)$, since $b_d \ge 0$ is strictly decreasing there
(Lemma~\ref{lem:bessel-profile}).

\emph{(1)} Let $d \ge 2$. For $b_d$, take $\lambda = 1$ in Lemma~\ref{lem:cm-exponential-rate}: the
upper bound holds for $r \ge 1$, and $e^rb_d(r) \ge m_d\sqrt r$ is unbounded. For $k_d$, put
$\lambda := q/c$ and $x := \lambda r$. Then $k_d(r) = \sum_{j \ge 0}b_d(xq^j)$, a series of positive
terms (Theorem~\ref{thm:matern-pyramid}(1)), so $k_d > 0$ and $e^{\lambda r}k_d(r) \ge e^xb_d(x)$,
which is unbounded. Above, $1 + xq^j \le q^j(1 + x)$ and $q^j - 1 \ge j(q - 1)$ give
\[
  e^xb_d(xq^j) \le M_d(1 + xq^j)^de^{-x(q^j - 1)} \le M_d(1 + x)^d\rho^j ,
  \qquad \rho := q^de^{-x(q - 1)} .
\]
For $x \ge \log(2q^d)/(q - 1)$, $\rho \le 1/2$, so $e^{\lambda r}k_d(r) \le 2M_d(1 + x)^d \le
2M_dK^d(1 + r)^d$ with $K := \max(1, q/c)$. Lemma~\ref{lem:cm-exponential-rate} applies. For $d = 1$,
$b_1 = 2e^{-y}$ on $(0, \infty)$ (Lemma~\ref{lem:bessel-profile}(3)) and
$k_1(r) = 2\sum_ie^{-rq^i/c}$, both completely monotone (Lemma~\ref{lem:laplace-atoms-admissible}(2),
with $\vartheta_i = q^i/c$ and a single atom for $b_1$).

\emph{(2)} In the profile of Theorem~\ref{thm:matern-pyramid}(5) the substitution $v = re^u/c$
gives $\int_0^\infty b_d(re^u/c)\,du = \int_{r/c}^\infty b_d(v)\,dv/v = P_d(r/c)$. For $d = 1$,
$b_1(v)/v = 2e^{-v}/v$, so $P_1 = 2E_1$. For $d = 3$, $b_3(v)/v = 2(1/v + 1)e^{-v}$
(Lemma~\ref{lem:bessel-profile}(3)), so $P_3 = 2E_1 + 2e^{-y}$. For $d = 2$,
$b_2(v)/v = 2K_1(v) = -2K_0'(v)$ and $K_0(v) \to 0$ as $v \to \infty$
(Lemma~\ref{lem:bessel-k}(1)--(2)), so $P_2 = 2K_0$. A Laplace transform
$y \mapsto \int_1^\infty e^{-ys}w(s)\,ds$ of a $w \ge 0$ whose transforms converge for every $y > 0$ is
completely monotone: for $y > y_0/2 > 0$ and $s > 1$, $s^ne^{-ys} \le C_{n,y_0}e^{-y_0s/4}$, so it may be
differentiated under the integral, with $(-1)^n$ times the $n$-th derivative equal to
$\int_1^\infty s^ne^{-ys}w(s)\,ds \ge 0$. This applies to $2E_1(y) = \int_1^\infty e^{-ys}(2/s)\,ds$, by
the substitution $v = ys$, and to $2K_0(y) = \int_1^\infty e^{-ys}\,2(s^2 - 1)^{-1/2}\,ds$
(Lemma~\ref{lem:bessel-k}(3)). Since $e^{-y}$ is completely monotone and sums of completely monotone
functions are, so is $P_3$. Let $d \ge 4$. The integrand $b_d(v)/v$ is continuous on $(0, \infty)$
and at most $2^{1+d/2}e^{-v/2}/a$ on $(a, \infty)$, so $P_d' = -b_d(y)/y$ on $(0, \infty)$. If $P_d$
were completely monotone, so would be $-P_d' = b_d(y)/y$, the derivatives of $-P_d'$ being those of
$P_d$ of one order higher. It is positive, $e^yb_d(y)/y \le M_d(1 + y)^d$ for $y \ge 1$, and, since
$d/2 - 1 \ge 1$, $e^yb_d(y)/y \ge m_d(y/2)\sqrt y/y = \tfrac12m_d\sqrt y$ for $y \ge 2$, which is
unbounded. Lemma~\ref{lem:cm-exponential-rate} with $\lambda = 1$ excludes it.
\end{proof}

% draft: the paragraph after Proposition 4.3, line 580.
For $d = 2, 3$ the failure at fixed $q$ thus disappears in the limit $q \downarrow 1$. What carries
in every dimension and for every $q$ is the mixing law of $V_t$ in
Theorem~\ref{thm:matern-pyramid}(1). A \emph{generalized gamma convolution} is a weak limit of finite
convolutions of gamma laws [\sscite{bondesson1992generalized}, \S 3.1, pp.~29--30], a class defined there
by its moment generating function (p.~29). The variable $V_t = \sum_ic^2q^{-2i}t^2X_i$ is a
convergent sum of independent exponential variables, so its law is one. The kernel is the Gaussian
mixed over that law, in every dimension. The claim is about the mixing law only: no class property
of $\mu_t$ itself in $\RR^d$ is claimed.

% CHANGED (blueprint round, 2026-11-16, review flag "three facts on completely monotone funct"): the
% three lemmas on completely monotone functions, lem:cm-log-convex, lem:cm-exponential-rate and
% lem:cm-uncancelled-atom, moved ahead of prop:matern-increments, the first proposition that uses
% them; a pure move, with a lead sentence before them.
% draft: the lead of section 6.1 before Lemma 6.2, line 1080, moved with the lemma (adr/0004).
The octave step is what a pyramid computes, and it is what changes most between the line and the
plane. On the line it is a pair of first-order recursions, one forward and one backward
(Theorem~\ref{thm:matern-pyramid-line}(3)). The two passes are the two factors of the symbol,
$1 + c^2t^2\omega^2 = (1 - ict\omega)(1 + ict\omega)$. The next lemma says that in $d \ge 2$ the symbol
has no such factors.

% draft: Lemma 6.2 (section 6.1), lines 1082-1084 of draft/spatial-hemigroup-pyramids.md, moved to
% section 4.1 (adr/0004, decision 3). Moved from the affine blueprint (06-paper-vi.tex at
% origin/main, lem:no-first-order-factorization), plan decision 5; statement and proof are the
% affine node's, which the draft prints word for word.
% CHANGED: the \lean names are this repository's, Pyramids.* in Formalization/Pyramids/NoFactorization.lean.
% CHANGED: the coefficient t is renamed r in the statement and the proof, since t is the scale
% everywhere else in the chapter (blind review, flag 13); the Lean declarations keep their t.
\begin{lemma}[no first-order factorization in dimension $d \ge 2$]
\label{lem:no-first-order-factorization}
\lean{Pyramids.no_first_order_factorization, Pyramids.first_order_factorization_line}\leanok
Let $d \ge 2$ and $r > 0$. The polynomial $1 + r^2|\omega|^2$ in $\omega \in \RR^d$ is not the
product of two polynomials of degree at most one with complex coefficients. On the line it is,
$1 + r^2\omega^2 = (1 - ir\omega)(1 + ir\omega)$.
% CHANGED (fidelity review, 2026-11-15, R26/R50/R69): the status line lists the declarations' hypotheses beside the node's (LINKAGE rule 6).
\statusT\quad\emph{(Cites nothing. \textbf{The declarations.}
\texttt{Pyramids.no\_first\_order\_factorization} carries the first sentence. Its hypotheses beside
the statement's: $2 \le d$ (same); $0 < r$ (same; the declarations call the coefficient $t$). Its
conclusion: there are no
$\alpha, \alpha' \in \CC$ and $\beta, \beta' \in \CC^d$ with
$1 + r^2\sum_i\omega_i^2 = (\alpha + \sum_i\beta_i\omega_i)(\alpha' + \sum_i\beta'_i\omega_i)$ for every
$\omega \in \RR^d$. That is the statement with ``product of polynomials'' read as an identity of
functions on $\RR^d$, which for polynomials is the same reading, and with $|\omega|^2$ written out
as $\sum_i\omega_i^2$. \texttt{Pyramids.first\_order\_factorization\_line} carries the second
sentence, for every real $r$ and $\omega$ (no hypothesis, so more general than the statement's
$r > 0$), as the identity $1 + r^2\omega^2 = (1 - i\,r\omega)(1 + i\,r\omega)$ in $\CC$. Both print
Lean core only (\texttt{Formalization/CIAxiomGuard.lean}). \textbf{Not claimed}: that no recursive
scheme of any other kind realizes the isotropic \Matern{} kernel in the plane, nor anything about
factorizations into operators that are not differential, such as the pseudo-differential square
roots. \textbf{Hypotheses} (the node's, each with what its declarations require): $d \ge 2$: same in
\texttt{no\_first\_order\_factorization}; the line, $d = 1$: same, \texttt{first\_order\_factorization\_line}
stated on $\RR$; $r > 0$: same in \texttt{no\_first\_order\_factorization}; weaker in
\texttt{first\_order\_factorization\_line}, which holds for every real $r$; $\omega \in \RR^d$: same,
as a real function on $\{0, \dots, d - 1\}$, the identity reading only its coordinates.)}
\end{lemma}

\begin{proof}
\leanok
Suppose $1 + r^2|\omega|^2 = (\alpha + \langle\beta,\omega\rangle)(\alpha' + \langle\beta',\omega\rangle)$
for every $\omega \in \RR^d$, with $\alpha,\alpha' \in \CC$ and $\beta,\beta' \in \CC^d$. On the plane
$\omega = se_1 + ue_2$ this is
$1 + r^2(s^2 + u^2) = (\alpha + \beta_1s + \beta_2u)(\alpha' + \beta'_1s + \beta'_2u)$ for all real
$s, u$, and the values at $(0,0)$, $(\pm1,0)$, $(0,\pm1)$ and $(1,1)$ fix the six coefficients:
$\alpha\alpha' = 1$, $\alpha\beta'_1 + \alpha'\beta_1 = 0$, $\alpha\beta'_2 + \alpha'\beta_2 = 0$,
$\beta_1\beta'_1 = \beta_2\beta'_2 = r^2$ and $\beta_1\beta'_2 + \beta_2\beta'_1 = 0$. Multiplying the
linear ones by $\alpha'$ gives $\beta'_k = -\alpha'^2\beta_k$, so the mixed one reads
$\alpha'^2\beta_1\beta_2 = 0$, and then
$r^4 = \beta_1\beta'_1\beta_2\beta'_2 = (\alpha'^2\beta_1\beta_2)^2 = 0$, against $r > 0$. The line's
factorization is a direct check, $(1 - ir\omega)(1 + ir\omega) = 1 - i^2r^2\omega^2$.
\end{proof}

% draft: the comment after Lemma 6.2, lines 1086-1088, without its last sentence (the measured fill
% of section 6.2, which stays with chapter 6).
The constant term is the obstruction. Without it the symbol factors in the plane,
$\omega_1^2 + \omega_2^2 = (\omega_1 + i\omega_2)(\omega_1 - i\omega_2)$. More passes do not help. A
cascade of first-order passes inverts a product of polynomials of degree at most one. Degrees add,
so for a symbol of degree two at most two factors are nonconstant. The lemma is about polynomial
symbols, which is what a cascade of first-order passes along fixed directions realizes. It does not
exclude every recursive scheme. A pseudo-differential square root of $1 - r^2\Delta$ is not covered,
and neither is a recursion that does not run along lines.

The next proposition gives the generator, the reason the octave step cannot be split, and the
reason it cannot be sampled.

% draft: Proposition 4.4 (section 4.1), lines 594-598 of draft/spatial-hemigroup-pyramids.md. Its
% clause (3) is the Matern case of prop:isotropic-stage-centre-matern of a version of the isotropic
% paper not yet released; restated here with its proof, as the draft does.
% CHANGED: the constant C_d := Gamma(d/2 - 1)/(4 pi^{d/2}) of clause (3) renamed N_d, the constant
% of the Newtonian potential (BLUEPRINT-INPUTS; REFEREE-draft-sections-3-4.md, item 2, finding 2).
% CHANGED: clause (2)'s "(Lemma 6.2)" is lem:no-first-order-factorization, moved before it.
% CHANGED: an opening line binds q, c, t, r, h (blind review, flag 11). Clause (1) no longer
% restates the generator of Theorem 4.1(4); it states the size of its symbol, "of the order of"
% read as an error O(1/t), which the proof gives (blind review, flag 12). Clause (3)'s "So ... the
% periodization diverges" reads "Also", the divergence being proved independently (flag 30).
% CHANGED: clause (2) names the octave step as L_r at r = ct, where it said "the octave step's
% inverse" with r free; clause (3)'s "like" reads "asymptotic to", the ratio tending to 1, which
% is what the proof shows (second blind review, flag 7; the overfull lines).
% CHANGED (fifth blind pass, 2026-11-15, flag 12): clause (1) states the size as the uniform bound
% its declaration matern_generator_octave_1_size states, one constant C for every t > 0 where
% ct|omega| >= 1; "O(1/t) as |omega| -> infinity" said only O(1) at fixed t.
% CHANGED (fifth blind pass, 2026-11-15, flag 16): clause (3) says what "undefined" means: a
% density is fixed only almost everywhere, and the representative r^{-d}G(x/r) of
% lem:bessel-profile(1), of which the limits are read, has no finite value at 0.
\begin{proposition}[generator, octave step, no central value]
\label{prop:matern-generator-octave}
\lean{Pyramids.matern_generator_octave_1_terms, Pyramids.matern_generator_octave_1_symbol,
Pyramids.matern_generator_octave_1_size, Pyramids.matern_generator_octave_2,
Pyramids.matern_generator_octave_3, Pyramids.matern_generator_octave_3_plane,
Pyramids.matern_generator_octave_3_high, Pyramids.matern_generator_octave_3_periodization}
% CHANGED (wave 3b, 2026-11-15): the eight declarations are machine-checked, so statement and proof
% are marked proved.
\leanok
\uses{thm:matern-pyramid, lem:bessel-profile, lem:bessel-k, lem:no-first-order-factorization}
Let $q > 1$, $c > 0$ and $d \ge 1$, let $(\mu_{s,t})$ be the \Matern{} pyramid of
Theorem~\ref{thm:matern-pyramid}, and let $t > 0$, $r > 0$ and a lattice spacing $h > 0$ be given.
\begin{enumerate}
\item \emph{(Generator.)} The generator $\mathcal{H}_t = \frac2t\sum_{i \ge 1}(L_{ctq^{-i}} - I)$ of
Theorem~\ref{thm:matern-pyramid}(4) has the symbol
$\sigma_t(\omega) = -\frac2t\sum_i\frac{\beta_i|\omega|^2}{1 + \beta_i|\omega|^2}$, $\beta_i =
c^2q^{-2i}t^2$. Each term is the symbol of a bounded operator. The sum is unbounded in $\omega$:
there is a constant $C$ with
$\bigl|\,|\sigma_t(\omega)| - \frac2t\log_q(ct|\omega|)\bigr| \le C/t$ for every $t > 0$ and every
$\omega$ with $ct|\omega| \ge 1$.
\item \emph{(The octave step.)} Let $d \ge 2$. The symbol $1 + r^2|\omega|^2$ of the inverse of
$L_r$, the octave step $\mu_{t,qt}$ when $r = ct$, is not a product of two polynomials of degree at
most one with complex coefficients (Lemma~\ref{lem:no-first-order-factorization}).
\item \emph{(No central value.)} Let $d \ge 2$. The density $r^{-d}G(x/r)$ of $L_r$ of
Lemma~\ref{lem:bessel-profile}(1) is $+\infty$ at the origin and tends to $+\infty$ there. For
$d = 2$ it is asymptotic to $(2\pi r^2)^{-1}\log(1/|x|)$. For $d \ge 3$ it is asymptotic to
$N_dr^{-2}|x|^{2-d}$, with $N_d := \Gamma(d/2 - 1)/(4\pi^{d/2})$. The periodization
$\sum_{n \in \ZZ^d}(1 + r^2|\omega + 2\pi n/h|^2)^{-1}$ of its transform diverges.
\end{enumerate}
% CHANGED (blueprint round, 2026-11-16, review flag "a density is fixed only almost everywher"): the
% rationale of (3), "So its sample on a lattice is undefined at 0: a density is fixed only almost
% everywhere, and the representative ... has no finite value there", moves to the remark after the
% proposition; the value +infinity of the representative at 0, which it asserted, is now the
% clause's first assertion, as matern_generator_octave_3 states it (maternDensityE d r 0 = top).
% The claims are unchanged.
% CHANGED (wave 3b, 2026-11-15): the status line names the facts the checked proof reads. It read
% "Cites nothing of its own. Clause (3) uses G = (2 pi)^{-1} K_0(|.|) at d = 2 and the behaviour of
% K_0 at the origin, from Lemma bessel-k, which carries entry A51; the case d >= 3 is monotone
% convergence and one Gamma integral."
% CHANGED (fidelity review, 2026-11-15, R26/R50/R69): the status line lists the declarations' hypotheses beside the node's (LINKAGE rule 6).
% CHANGED (fidelity review, 2026-11-15, R49): the status line names where the L^1 convergence, the
% evolution equation and the unboundedness of the generator are proved: Theorem matern-pyramid(4).
\statusT\quad\emph{(Cites nothing of its own. Clause (3) uses $G = (2\pi)^{-1}K_0(|\cdot|)$ at
$d = 2$ and the behaviour of $K_0$ at the origin, from Lemma~\ref{lem:bessel-k}, which carries
entry A51; the case $d \ge 3$ is dominated convergence and one Gamma integral. Machine-checked.
\textbf{Delegated.} The absolute $L^1$ convergence of the series on $W^{2,1}(\RR^d)$, the evolution
equation $\partial_t(\mu_t * f) = \mathcal{H}_t(\mu_t * f)$ and the unboundedness of
$\mathcal{H}_t$ in the $L^1$ norm are Theorem~\ref{thm:matern-pyramid}(4), with the declarations
\texttt{matern\_pyramid\_4\_series} (the convergence, and the unboundedness through
\texttt{maternGenerator\_unbounded}) and \texttt{matern\_pyramid\_4} (the equation), outside this
node's tag; clause (1) states the symbol, the bounded terms and the size of the symbol only.
\textbf{Hypotheses} (the node's, each with what its declarations require): $q > 1$: same in
\texttt{\_1\_symbol} and \texttt{\_1\_size}; weaker in \texttt{\_1\_terms}, which does not require it;
absent from (2) and (3); $c > 0$: same in \texttt{\_1\_size}; weaker in \texttt{\_1\_terms} and
\texttt{\_1\_symbol}; absent from (2) and (3); $d \ge 1$: same in \texttt{\_1\_size}; weaker in
\texttt{\_1\_terms} and \texttt{\_1\_symbol}; $(\mu_{s,t})$ the \Matern{} pyramid: absent, the
declarations reading $L_r$ and $\mathcal{H}_t$ directly, the octave step $\mu_{t,qt} = L_{ct}$ being
Theorem~\ref{thm:matern-pyramid}(1); $t > 0$: same in \texttt{\_1\_terms} and \texttt{\_1\_size};
weaker in \texttt{\_1\_symbol}, which does not require it; (1) the domain of $\mathcal{H}_t$: same,
\texttt{IsW21} in \texttt{\_1\_symbol}, and integrable $g$ for the bounded terms in
\texttt{\_1\_terms}; $ct|\omega| \ge 1$: same; (2), (3) $d \ge 2$: same; $d = 2$ and $d \ge 3$ for the
asymptotics: same; $r > 0$: same in \texttt{\_2}, \texttt{\_3}, \texttt{\_3\_plane},
\texttt{\_3\_high}; weaker in \texttt{\_3\_periodization}, which does not require it; $h > 0$: weaker
in \texttt{\_3\_periodization}, which does not require it, and no other declaration reads $h$.)}
\end{proposition}

% CHANGED (wave 3b, 2026-11-15): the proof follows the machine-checked route. The old proof cited
% the symbol from the proof of Theorem matern-pyramid(4) without saying how the series is tested
% against a character, did not derive the unboundedness from the size, read G(x) -> infinity and
% the limit of |x|^{d-2}G(x) for d >= 3 by monotone convergence, and proved the periodization's
% divergence by counting lattice points in dyadic shells 2^k <= |n| < 2^{k+1}, with the bound
% |omega + 2 pi n/h| <= A|n| for n != 0.
% CHANGED (fifth blind pass, 2026-11-15, flags 13-15): the transform of L_r is cited from
% lem:bessel-profile(1); the bound on a term is written for the operator T_i on one g,
% ||T_i g||_1 <= (4/t)||g||_1, where an operator norm was bounded by norms of g; the size paragraph
% names the constant C = 2(1 + G) of the statement; "a, b in N" reads "integers a, b >= 0".
\begin{proof}
\leanok
\uses{thm:matern-pyramid, lem:bessel-profile, lem:bessel-k, lem:no-first-order-factorization}
\emph{(1)} By Fubini's theorem $\int (L_r * g)(x)\,e^{i\langle x,\omega\rangle}\,dx =
\widehat L_r(\omega)\widehat g(\omega)$ for $g \in L^1$, with $\widehat L_r(\omega) = (1 + r^2|\omega|^2)^{-1}$
(Lemma~\ref{lem:bessel-profile}(1)),
and $(1 + \beta_i|\omega|^2)^{-1} - 1 = -\beta_i|\omega|^2/(1 + \beta_i|\omega|^2)$. So the term
$i$ is the symbol of $T_i := \frac2t(L_{ctq^{-i}} - I)$, which is bounded on $L^1$:
$\|T_ig\|_1 \le \frac2t(\|L_{ctq^{-i}} * g\|_1 + \|g\|_1) \le \frac4t\|g\|_1$. On $W^{2,1}$ the series converges
absolutely in $L^1$ (Theorem~\ref{thm:matern-pyramid}(4)), so it is integrated against the
character term by term, and its symbol is $\sigma_t$.

For the size, let $y := ct|\omega| \ge 1$, so that $\beta_i|\omega|^2 = y^2q^{-2i}$, and let
$M := \lfloor\log_q y\rfloor$, so that $q^M \le y < q^{M+1}$. For $i \le M$ the term
$\beta_i|\omega|^2/(1 + \beta_i|\omega|^2)$ lies between $1 - (\beta_i|\omega|^2)^{-1}$ and $1$,
and $(\beta_i|\omega|^2)^{-1} \le q^{-2(M-i)}$. For $i > M$ it is at most
$\beta_i|\omega|^2 < q^{-2(i-M-1)}$. Each of the two geometric sums is at most
$A := (1 - q^{-2})^{-1}$, so the sum of the terms lies between $M - A$ and $M + A$. With
$M \le \log_q y < M + 1$,
$\bigl|\,|\sigma_t(\omega)| - \frac2t\log_q(ct|\omega|)\bigr| \le \frac2t(1 + A)$
for $ct|\omega| \ge 1$, one constant for every $t > 0$: the bound of the statement with
$C := 2(1 + A)$. Along an axis, at $ct|\omega| = q^N$, this
gives $|\sigma_t(\omega)| \ge \frac2t(N - 1 - A)$, which is unbounded in $N$.

\emph{(2)} is Lemma~\ref{lem:no-first-order-factorization}, with $|\omega|^2 = \sum_i\omega_i^2$.

\emph{(3)} By Lemma~\ref{lem:bessel-profile}(1) the density of $L_r$ is $r^{-d}G(x/r)$, with $G$
the integral defined there. Let $0 < |x| < 1$. For $|x|^2 < s \le 1$ the integrand
$e^{-s}(4\pi s)^{-d/2}e^{-|x|^2/(4s)}$ is at least $\eta_ds^{-1}$, with
$\eta_d := e^{-1}(4\pi)^{-d/2}e^{-1/4}$, as $s^{-d/2} \ge s^{-1}$ for $s \le 1$ and $d \ge 2$. So
\[
  G(x) \ge \eta_d\int_{|x|^2}^1 \frac{ds}{s} = 2\eta_d\log(1/|x|) ,
\]
which tends to $\infty$ as $x \to 0$. The integrand decreases in $|x|$, so $G(0) \ge G(x)$ for
every $x$, and $G(0) = \infty$. For $d = 2$, Lemma~\ref{lem:bessel-k}(1) at $\nu = 0$ gives
\[
  G(x) = (4\pi)^{-1}\int_0^\infty s^{-1}e^{-s-|x|^2/(4s)}\,ds = (2\pi)^{-1}K_0(|x|) .
\]
By Lemma~\ref{lem:bessel-k}(3), $K_0(y) = \log(1/y) + O(1)$ as $y \downarrow 0$, so the density
$(2\pi r^2)^{-1}K_0(|x|/r)$ differs from $(2\pi r^2)^{-1}\log(1/|x|)$ by
$(2\pi r^2)^{-1}(\log r + O(1))$. That difference stays bounded while
$(2\pi r^2)^{-1}\log(1/|x|)$ tends to $\infty$, so the ratio tends to $1$. For $d \ge 3$ the
substitution $s = |x|^2v$ gives
\[
  |x|^{d-2}G(x) = \int_0^\infty e^{-|x|^2v}\,(4\pi v)^{-d/2}e^{-1/(4v)}\,dv .
\]
The integrand is at most $(4\pi v)^{-d/2}e^{-1/(4v)}$, which is integrable for $d \ge 3$, so by
dominated convergence the integral tends to $\int_0^\infty (4\pi v)^{-d/2}e^{-1/(4v)}\,dv$ as
$x \to 0$. The substitution $v = 1/w$ makes this
$(4\pi)^{-d/2}\int_0^\infty w^{d/2-2}e^{-w/4}\,dw = (4\pi)^{-d/2}4^{d/2-1}\Gamma(d/2 - 1) = N_d$.
Then $r^{-d}G(x/r) = r^{-d}|x/r|^{2-d}\bigl(N_d + o(1)\bigr)
= N_dr^{-2}|x|^{2-d}\bigl(1 + o(1)\bigr)$.

For the periodization, fix $\omega$. As $(u + v)^2 \le 2u^2 + 2v^2$,
$1 + r^2|\omega + 2\pi n/h|^2 \le B(1 + |n|^2)$ with
$B := 1 + 2r^2|\omega|^2 + 2r^2(2\pi/h)^2$, so the term at $n$ is at least $B^{-1}(1 + |n|^2)^{-1}$.
Take $n = (a, b, 0, \dots, 0)$ with integers $a, b \ge 0$, which $d \ge 2$ allows. For fixed $a$ the terms
with $b \le a$ are at least $B^{-1}(1 + 2a^2)^{-1}$, so the row $a$ sums to at least
$B^{-1}(a + 1)/(1 + 2a^2) \ge B^{-1}/(2(a + 1))$, and the sum over $a$ diverges. A sum of
nonnegative terms with a divergent part diverges.
% CHANGED (blueprint round, 2026-11-16, review flag on "G := (1 - q^{-2})^{-1}"): the local
% constants G, K and L, which collide with the density G, the Bessel function K_0 and the law L_r,
% are renamed A, eta_d and B. No step changed.
\end{proof}

% draft: the comment after Proposition 4.4, line 600.
By clause (2) the line's forward--backward pair of first-order recursions has no direct
counterpart, and in every dimension the octave step is one solve of $(I - r^2\Delta)v = g$. The
clause does not exclude every recursive scheme, only a cascade of first-order passes along fixed
directions. Clause (3) says that the line's sampled Laplace family
(Proposition~\ref{prop:sampled-laplace-exact}) has no counterpart by sampling in the plane. A
density is fixed only almost everywhere, but the failure does not depend on the version. The
continuous version off $0$ of Lemma~\ref{lem:bessel-profile}(1) tends to $+\infty$ at $0$ by clause
(3), and every version agrees with it almost everywhere. So no version is continuous at $0$ with a
finite value, and the sample of the octave step's density at the lattice point $0$ is undefined. The
octave step can be solved for on a grid, as the discrete Helmholtz family does
(Proposition~\ref{prop:discrete-helmholtz}), but not sampled. The kernels $\mu_t$ are bounded and
smooth, so their samples are defined. The level difference is a Laplacian, by the next proposition.

% draft: Proposition 4.5 (section 4.1), lines 602-604 of draft/spatial-hemigroup-pyramids.md.
% CHANGED: the input class and the reading are named: f in L^1(R^d), t > 0, and (4.2) holding at
% every point, every term being continuous. The draft says "for an input f" and reads it "on the
% transforms".
% CHANGED: the record is defined once, in Theorem 4.1(3); this node refers to it (second blind
% review, flag 12).
\begin{proposition}[the Laplacian pyramid is a Laplacian]
\label{prop:laplacian-pyramid}
% CHANGED (wave 3b, 2026-11-15): the declaration is machine-checked, so statement and proof are
% marked proved.
\lean{Pyramids.laplacian_pyramid}
\leanok
\uses{thm:matern-pyramid, lem:bessel-profile}
For an input $f \in L^1(\RR^d)$ write $u(t, \cdot) := \mu_t * f$ for its record at scale $t$
(Theorem~\ref{thm:matern-pyramid}(3)).
On the \Matern{} pyramid of Theorem~\ref{thm:matern-pyramid}, with octave step $L_{ct}$, the level
difference at $t > 0$ is
\[
  u(t, \cdot) - u(qt, \cdot) = (I - L_{ct})\,u(t, \cdot) = -c^2t^2\,\Delta\,u(qt, \cdot) ,
\]
an identity of continuous bounded functions, holding at every point.
% CHANGED (wave 3b, 2026-11-15): the status line names the facts the checked proof reads. It read
% "Cites nothing: the transforms, and the uniqueness of the Fourier transform of a tempered
% distribution."
% CHANGED (fidelity review, 2026-11-15, R26/R50/R69): the status line lists the declarations' hypotheses beside the node's (LINKAGE rule 6).
\statusT\quad\emph{(Cites nothing: the transforms of finite measures and their uniqueness,
differentiation under the integral sign, and integration by parts against a smooth function of
compact support. Machine-checked. \textbf{Hypotheses} (the node's, each with what its declarations
require): $q > 1$, $c > 0$, $d \ge 1$, from Theorem~\ref{thm:matern-pyramid}: same; the \Matern{}
pyramid: same, every $\mu$ with \texttt{IsMaternPyramid q c}; $f \in L^1(\RR^d)$: same, real $f$
with \texttt{Integrable f}; $t > 0$: same.)}
\end{proposition}

% CHANGED (wave 3b, 2026-11-15): the proof follows the machine-checked route. The old proof read
% the first equality from mu_{qt} = L_{ct} * mu_t without saying that the composition of records
% holds almost everywhere only, read the second as distributions from Lemma bessel-profile(1)
% without carrying the Laplacian from the test function to u(qt), and said that the derivatives of
% p_t * f are the convolutions of the derivatives of p_t without the domination that allows it.
% CHANGED (fifth blind pass, 2026-11-15, flag 17): the record u(t, .) is shown in L^1, which
% Lemma bessel-profile(1) requires, by ||mu_t * f||_1 <= ||f||_1.
\begin{proof}
\leanok
\uses{thm:matern-pyramid, lem:bessel-profile}
\emph{The records are smooth.} For $s > 0$ the kernel $\mu_s$ has a $C^\infty$ density $p_s$ with
every derivative bounded (Theorem~\ref{thm:matern-pyramid}(1)), so
$u(s, x) = \int f(z)\,p_s(x - z)\,dz$. For $k \le 2$ the integrand's $k$-th derivative in $x$ is
dominated by $|f(z)|\sup|D^kp_s|$, which is integrable, so differentiation under the integral sign
gives $D^ku(s, \cdot) = \int f(z)\,D^kp_s(\cdot - z)\,dz$, continuous by dominated convergence and
bounded by $\|f\|_1\sup|D^kp_s|$. So $u(t, \cdot)$ and $u(qt, \cdot)$ are continuous and bounded,
$u(qt, \cdot)$ is $C^2$, and $\Delta u(qt, \cdot)$, a sum of second derivatives, is continuous and
bounded.

\emph{The first equality.} On the transforms, $\widehat\mu_{qt} = \widehat\mu_t\,\widehat\mu_{t,qt}$, and
$\mu_{t,qt} = L_{ct}$ (Theorem~\ref{thm:matern-pyramid}(1)), so $\mu_{qt} = \mu_t * L_{ct}$ by the
uniqueness of the transform of a finite measure. By Fubini's theorem
$L_{ct}u(t, \cdot) = u(qt, \cdot)$ almost everywhere. Both sides are continuous: the left one
because $u(t, \cdot)$ is continuous and bounded and $L_{ct}$ is finite (dominated convergence). So
they agree at every point, which is the first equality.

\emph{The second.} Let $\varphi$ be smooth with compact support. For a $C^2$ function $v$, two
integrations by parts along each coordinate give $\int v\,\partial_j^2\varphi = \int
(\partial_j^2 v)\varphi$, the boundary terms vanishing with $\varphi$, so
$\int v\,\Delta\varphi = \int (\Delta v)\varphi$. The record $g := u(t, \cdot)$ is in $L^1$, with
$\|\mu_t * f\|_1 \le \|f\|_1$ by Fubini's theorem, $\mu_t$ being a probability measure.
Lemma~\ref{lem:bessel-profile}(1), applied to $g$, gives $\int (L_{ct}g)(\varphi - c^2t^2\Delta\varphi) = \int g\varphi$.
With $L_{ct}g = u(qt, \cdot)$ and $v = u(qt, \cdot)$,
\[
  \int \bigl(u(qt, \cdot) - c^2t^2\,\Delta u(qt, \cdot) - u(t, \cdot)\bigr)\,\varphi = 0 .
\]
The function in brackets is continuous, and a locally integrable function that integrates to $0$
against every such $\varphi$ vanishes almost everywhere, so it vanishes at every point. That is
$u(t, \cdot) - u(qt, \cdot) = -c^2t^2\,\Delta u(qt, \cdot)$, the second equality.
\end{proof}

% draft: the comment after Proposition 4.5, line 606.
For $d = 1$ this is the Helmholtz identity of \prog{V-selection v0.1}{Prop.~5.6(1), p.~60} for the
\Matern{} step of index $\gamma = 1$, the law $L_{ct}$ itself; that release states its jet results
on the line alone. So the ``Laplacian pyramid'' of the \Matern{} pyramid is, exactly, the Laplacian
of the next level, scale-normalized by $c^2t^2$. Burt and Adelson's Laplacian pyramid
[\sscite{burt1983laplacian}, p.~535] is a \emph{difference} of Gaussians named after the Laplacian
it approximates.
% Anchored 2026-11-12 by the librarian (records/pyramids/CITATIONS-acquired.md, item 2): p. 535
% defines L_l = g_l - EXPAND(g_{l+1}).
On the Gaussian pyramid the level difference is a difference of Gaussians. It lies in the
$2$-jet only in the limit $q \downarrow 1$, through the heat equation. The identity above holds
between the records on the continuum, before any decimation. Burt and Adelson's level subtracts the
expanded, decimated next level. On one grid the
identity holds for the discrete Helmholtz family with the five-point Laplacian $\Delta_h$ in place
of $\Delta$ (Proposition~\ref{prop:discrete-helmholtz}). What decimation does to it is
Proposition~\ref{prop:subsampling-discrepancy}. The pyramid also stores the scale-normalized jet of
\prog{V-selection v0.1}{Def.~5.1} at every level with one set of stencils: the normalized components
$t^{|\alpha|}\partial^\alpha u$ at scale $q^nt$ on the grid $q^nh\ZZ^d$ are computed by the same
finite-difference stencils, with the same normalization, at every level. The jet in $\RR^d$ is
chapter~\ref{ch:jet}.

Section~\ref{sec:matern-counting} compares the \Matern{} pyramid with a separable twin, applied one
axis at a time: the product of the line's \Matern{} pyramids along the coordinate axes.

% CHANGED (blueprint round, 2026-11-16, review flag "with level transforms $\prod_je^{-F_q(s|"): new
% definition (new label def:frame-matern-pyramid), the frame pyramid that prop:creation-plane(3)
% described by its level transforms and steps. Its declaration is the specification
% IsFrameMaternPyramid, as IsMaternPyramid is the Matern pyramid's.
\begin{definition}[the frame \Matern{} pyramid]
\label{def:frame-matern-pyramid}
\lean{Pyramids.IsFrameMaternPyramid}\leanok
\uses{def:matern-pyramid-data}
Let $q > 1$, $c > 0$ and $d \ge 1$. A \emph{frame \Matern{} pyramid} with the constants $q$ and $c$
is a collection of probability measures $\mu_{s,t}$ on $\RR^d$, $0 \le s \le t$, with the transforms
\[
  \widehat{\mu_{s,t}}(\omega)
   = \prod_{j=1}^d\exp\bigl[-\bigl(F_q(t|\omega_j|) - F_q(s|\omega_j|)\bigr)\bigr] .
\]
Its kernels are $\mu_t := \mu_{0,t}$, and its octave steps the stages $\mu_{t,qt}$, $t > 0$.
\statusT\quad\emph{(Vocabulary. Its existence and its octave steps are
Proposition~\ref{prop:frame-matern-pyramid-exists}. \texttt{IsFrameMaternPyramid q c} is the
specification of the transforms, probability measures for $0 \le s \le t$, as
\texttt{IsMaternPyramid} is the \Matern{} pyramid's; it does not take $q > 1$, $c > 0$ or
$d \ge 1$, and the statements that read it carry the first two.)}
\end{definition}

% CHANGED (blueprint round, 2026-11-16, review flag "with level transforms $\prod_je^{-F_q(s|"): new
% proposition (new label prop:frame-matern-pyramid-exists). Its statement is the conclusion of
% frame_matern_pyramid_exists, which was in the tag of prop:creation-plane; its proof is the
% paragraph that sat in the proof of prop:creation-plane(3), written out.
\begin{proposition}[the frame \Matern{} pyramid exists]
\label{prop:frame-matern-pyramid-exists}
\lean{Pyramids.frame_matern_pyramid_exists}\leanok
\uses{def:frame-matern-pyramid, thm:matern-pyramid-line}
Let $q > 1$, $c > 0$ and $d \ge 1$. There is a frame \Matern{} pyramid $(\mu_{s,t})$ in dimension
$d$ with the constants $q$ and $c$, and its octave step $\mu_{t,qt}$ has the transform
$\prod_{j=1}^d(1 + c^2t^2\omega_j^2)^{-1}$ for every $t > 0$.
\statusT\quad\emph{(Cites nothing of its own: it reads the \Matern{} pyramid on the line,
Theorem~\ref{thm:matern-pyramid-line}, which carries its own entries. Machine-checked.
\textbf{Hypotheses} (the node's, each with what its declaration requires): $q > 1$, $c > 0$: same;
$d \ge 1$: weaker, \texttt{frame\_matern\_pyramid\_exists} holds for every natural $d$, its proof
not reading $d$; $t > 0$: same.)}
\end{proposition}

\begin{proof}
\leanok
\uses{def:frame-matern-pyramid, thm:matern-pyramid-line}
Let $(\nu_{s,t})$ be the \Matern{} pyramid on the line of Theorem~\ref{thm:matern-pyramid-line}
with the constants $q$ and $c$. Its laws are probability measures with the transforms
$\exp\bigl[-\bigl(F_q(t|\omega|) - F_q(s|\omega|)\bigr)\bigr]$, $0 \le s \le t$, as $F_q$ is even.
Let $\mu_{s,t}$ be the law of $d$ independent coordinates, each of law $\nu_{s,t}$, that is the
product measure of $d$ copies of $\nu_{s,t}$. It is a probability measure, and the transform of a
product measure is the product of the transforms of its factors, so $\mu_{s,t}$ has the transform of
Definition~\ref{def:frame-matern-pyramid}. The octave identity
$F_q(qt\omega) - F_q(t\omega) = \log(1 + c^2t^2\omega^2)$ of
Theorem~\ref{thm:matern-pyramid-line}(1), at $\omega = |\omega_j|$, gives
$\exp\bigl[-\bigl(F_q(qt|\omega_j|) - F_q(t|\omega_j|)\bigr)\bigr] = (1 + c^2t^2\omega_j^2)^{-1}$,
so the octave step $\mu_{t,qt}$ has the transform $\prod_j(1 + c^2t^2\omega_j^2)^{-1}$.
\end{proof}

\section{Counting in the plane}
\label{sec:matern-counting}

% draft: the lead of section 4.2, line 615.
On the line the case for the \Matern{} pyramid was counting. It creates no zero crossings and no
extrema, from the signal and from rung to rung (Theorem~\ref{thm:matern-pyramid-line-noncreation},
Proposition~\ref{prop:ladder-non-creation-line}). In the plane the diffusion equation itself creates
local extrema \prog{V-selection v0.1}{\S 6}. One step of any symmetric absolutely continuous kernel
with a nonnegative transform turns some image with one local maximum into one with two. On records,
the records of Proposition~\ref{prop:laplacian-pyramid} read along the ladder, the count can rise
at every rung: for each rung of the \Matern{} pyramid, and of its frame twin, the product of the
line's \Matern{} pyramids along the axes, some input gains a local maximum there. On the Gaussian
ladder it rises at some rung. On a grid it rises at one step of every symmetric mask $w \ne
\delta_0$ with a nonnegative transform, before and after decimation
(Proposition~\ref{prop:creation-grid}). The witness is a ridge between two bumps: smoothing removes
the thin ridge faster than the bumps. Clause (2) of the next proposition lifts a creation in one
step to records: under a law whose transform decays more slowly than every Gaussian, each witness
is itself a record.

% draft: Proposition 4.6 (section 4.2), lines 617-626 of draft/spatial-hemigroup-pyramids.md.
% CHANGED: clause (3)'s "subtracting a level from g" is stated for the records minus a level,
% mu_s * g - C and mu_{qs} * g - C, so that no input outside the class of the clause is formed
% (second blind review, flag 8). The content is the same: mu_s * (g - C) = mu_s * g - C.
\begin{proposition}[creation in $d \ge 2$, from the signal and on records]
\label{prop:creation-plane}
\lean{Pyramids.creation_plane_1, Pyramids.creation_plane_2, Pyramids.creation_plane_3,
Pyramids.creation_plane_3_frame, Pyramids.creation_plane_4, Pyramids.creation_plane_5}
% CHANGED (wave 3b, 2026-11-15): all six declarations are machine-checked, so statement and proof
% are marked proved. The line marked the node not ready.
\leanok
% CHANGED (wave 3b, 2026-11-15): thm:matern-pyramid-line added to the uses, which the frame
% pyramid's existence reads. The line read "\uses{thm:matern-pyramid, lem:bessel-profile}".
% CHANGED (blueprint round, 2026-11-16): the frame pyramid's existence is
% prop:frame-matern-pyramid-exists, which carries frame_matern_pyramid_exists and the use of
% thm:matern-pyramid-line; the tag keeps the five declarations of the clauses.
\uses{thm:matern-pyramid, lem:bessel-profile, def:frame-matern-pyramid,
prop:frame-matern-pyramid-exists}
Let $d \ge 2$, write $x = (x_1, x')$, and for $a > 1$ and $\delta > 0$ let
\[
  f_{a,\delta}(x) := e^{-|x - ae_1|^2/2} + e^{-|x + ae_1|^2/2}
   + 8a^2\,e^{-x_1^2/(8a^2)}\,e^{-|x'|^2/(2\delta^2)} ,
\]
the \emph{witness}: two Gaussian bumps at $\pm ae_1$ joined by a thin ridge along the axis.
\begin{enumerate}
% CHANGED (2026-11-15, a strengthening)
% The symmetry of mu is not assumed: a nonnegative transform is real, so mu is symmetric. The old
% wording: "Let $\mu$ be a symmetric absolutely continuous probability measure on".
% CHANGED (review round 2, batch C, 2026-11-23, a strengthening)
% The hypothesis on the transform is dropped: W = mu * B is positive at 0 and tends to 0 for every
% probability law. The old wording: "on $\RR^d$ with $\widehat\mu \ge 0$. For suitable".
\item \emph{(One step.)} Let $\mu$ be an absolutely continuous probability measure on
$\RR^d$. For suitable $a$ and $\delta$, $f := f_{a,\delta}$ has exactly one
local maximum and $\mu * f$ has at least two. For a suitable level $\alpha$, $\{f > \alpha\}$ is
connected and $\{\mu * f > \alpha\}$ is not.
\item \emph{(Every witness is a record.)} Let $\eta$ be a probability measure with
$\widehat\eta > 0$ and $1/\widehat\eta(\omega) \le C_\varepsilon e^{\varepsilon|\omega|^2}$ for every
$\varepsilon > 0$, with constants $C_\varepsilon$. Then every $f_{a,\delta}$ is $\eta * g$ with $g$
real, bounded, continuous and in $L^2$.
\item \emph{(The \Matern{} and frame \Matern{} ladders.)} For every $s > 0$ there is a real
$g \in \mathcal{S}(\RR^d)$ such that the record $\mu_s * g$ has exactly one local maximum and the
next record $\mu_{qs} * g = L_{cs} * (\mu_s * g)$ has at least two. For a suitable level $\alpha$,
$\mu_s * g - \alpha$ has one positive region and $\mu_{qs} * g - \alpha$ has at least two. The frame
\Matern{} pyramid (Definition~\ref{def:frame-matern-pyramid}) has the same property.
\item \emph{(The Gaussian ladder.)} For every $q > 1$ there are a positive $g$, a sum of three
Gaussians, and an $n \ge 1$ such that the Gaussian record of $g$ at variance
$\tau_n := (q^{2n} - 1)^{-1}$ has exactly one local maximum, the record at variance
$q^{2n}\tau_n$ has at least two, and
% CHANGED (review round 2, batch C, 2026-11-23, a strengthening)
% The rise is placed at a rung with exactly one local maximum. The old wording: "for some integer
% $0 \le k < n$ the record at variance $q^{2(k+1)}\tau_n$ has more local maxima than the record at
% variance $q^{2k}\tau_n$."
for some integer $0 \le k < n$ the record at variance $q^{2k}\tau_n$ has exactly one local maximum
and the record at variance $q^{2(k+1)}\tau_n$ has at least two.
% CHANGED (review round 2, batch C, 2026-11-23): new clause, the rung of (4) moved to a prescribed
% variance by a dilation of the input.
\item \emph{(Any variance.)} For every $q > 1$ and $\tau > 0$ there is a positive $g$ whose
Gaussian record at variance $\tau$ has exactly one local maximum and whose record at variance
$q^2\tau$ has at least two.
\end{enumerate}
% CHANGED (blueprint round, 2026-11-16, review flags "Write $h := h_{a,\delta}$", "(Records.)", "So
% the count is not monotone along the l" and the frame pyramid's definition): the witness h_{a,delta}
% is renamed f_{a,delta} and named "the witness", and the level C is renamed alpha, so that h stays
% the lattice spacing of the neighbouring nodes and C a constant; clause (2)'s label "(Records.)"
% names its content; in (4) the inference "So the count is not monotone along the ladder of ratio q,
% at some rung below the n-th" is the clause's assertion, as creation_plane_4 already states it
% (a k < n whose count is below the count at k + 1); in (3) the frame pyramid is
% def:frame-matern-pyramid, its existence prop:frame-matern-pyramid-exists, where the clause
% described it by its level transforms and steps. The claims are unchanged.
% CHANGED (wave 3b, 2026-11-15): the status line says which clauses are machine-checked and prices
% the rest. It read "(Cites nothing: Fourier inversion, Plancherel's theorem and the uniqueness of
% the transform. The Gaussian record ... is positive.)", the last sentence kept below.
% CHANGED (wave 3b, 2026-11-15): clause (2) is machine-checked, and the price is that of (3) alone.
% The sentences read "Clauses (1) and (4) are machine-checked. Clauses (2) and (3), and the
% existence of the frame pyramid, are not yet, priced at about 1900--2700 lines: for (2), about
% 600--900, the transform of the ridge computed over the coordinates, and Plancherel's theorem for a
% function both integrable and square-integrable, through the tempered distributions, where the
% library's transform on L^2 agrees with the distributional one; for (3), ...".
% CHANGED (2026-11-15, a strengthening)
% The declaration of (2) does not take d >= 2, which its proof does not read; the print keeps the
% setting. The status line had no such sentence before.
% CHANGED (wave 3b, 2026-11-15): the existence of the frame pyramid is machine-checked, and the
% price is that of the rest of (3). The sentences read "Clause (3), and the existence of the frame
% pyramid, are not yet, priced at about 1300--1800 lines: that ... is a Schwartz function (...),
% about 700--900, the frame pyramid as the product of the line's Matern pyramids, its step a product
% of Laplace laws, about 400--600, and the assembly, about 200--300."
% CHANGED (2026-11-15, a strengthening)
% The declaration of the frame pyramid's existence does not take d >= 1, which its proof does not
% read. The status line had no such sentence before.
% CHANGED (wave 3b, 2026-11-15): clause (3) is machine-checked, and no price is left. The sentences
% read "Clauses (1), (2) and (4), and the existence of the frame pyramid, are machine-checked; ...
% The rest of clause (3) is not yet, priced at about 900--1200 lines: that h-hat e^{F_q(s|omega|)}
% and h-hat prod_j e^{F_q(s|omega_j|)} are Schwartz functions (the derivatives of the series F_q of
% every order, and their bounds), about 700--900, and the assembly, about 200--300."
% CHANGED (fidelity review, 2026-11-15, R26/R50/R69): the status line lists the declarations' hypotheses beside the node's (LINKAGE rule 6).
% CHANGED (fidelity review, 2026-11-15, R46): the status line says how the counts of local maxima
% are read, and which clauses exclude plateaus.
\statusT\quad\emph{(Cites nothing: Fourier inversion, Plancherel's theorem, the uniqueness of the
transform and Faà di Bruno's bound on the derivatives of a composition. The Gaussian record at variance $v$ is $N(0, vI) * g$, and a positive region is a
connected component of the set where the function is positive. Every clause is machine-checked;
the frame pyramid's existence is Proposition~\ref{prop:frame-matern-pyramid-exists}. \textbf{The counts.} A local maximum is read in the
declarations as a weak one, $\ge$ on a neighbourhood (\texttt{IsLocalMax}), and ``exactly one'' and
``at least two'' count such points. In (1) and (3) the further conjunct, a superlevel set that is
connected before and not preconnected after, excludes a plateau as the source of the second
maximum; in (4) the counts are of weak-maximum points with no such conjunct.
\textbf{Hypotheses} (the node's, each with what its declarations require): $d \ge 2$: same in
\texttt{creation\_plane\_1}, \texttt{\_3}, \texttt{\_3\_frame}, \texttt{\_4}; weaker in
\texttt{creation\_plane\_2}, which requires only $d \ne 0$; $a > 1$, $\delta > 0$: same in
\texttt{\_2}; in (1) they are produced, as the clause says; (1) $\mu$ an absolutely continuous
probability measure: same, $\mu \ll$ Lebesgue measure; (2) $\eta$ a probability
measure with $\widehat\eta > 0$ and the growth bound on $1/\widehat\eta$: same, the bound read on the
real part, which is $\widehat\eta$; (3) the \Matern{} pyramid, $q > 1$, $c > 0$: same, every $\mu$
with \texttt{IsMaternPyramid q c}; the frame \Matern{} pyramid: same, every $\mu$ with
\texttt{IsFrameMaternPyramid q c}, Definition~\ref{def:frame-matern-pyramid}; $s > 0$: same;
(4) $q > 1$: same; the rung $k$ of the rise: same, $k < n$ in $\NN$, the counts read as
cardinalities in $\NN \cup \{\infty\}$, $1$ at rung $k$ and at least $2$ at rung $k + 1$;
(5) $q > 1$, $\tau > 0$: same, in \texttt{creation\_plane\_5}, the counts read as in (4).)}
\end{proposition}

% CHANGED (wave 3b, 2026-11-15): clauses (1) and (4) follow the machine-checked route. In (1) the
% bound phi'(x_1) <= x_1 comes from the monotonicity of y e^{-y^2/2} + y/2 in place of
% phi'' <= 4e^{-3/2}; the ridge's L^1 norm tends to 0 by dominated convergence, the axis being
% null, in place of its closed form; the density is split as K rho <= M rho + 8a^2 (K - M)^+; and
% the symmetry of mu is read off the real transform. In (4) the identity h = N(0, tau I) * g is
% computed coordinatewise, and the record at tau + 1 is read through N(0, (tau+1)I) = N(0, tau I) *
% N(0, I).
% CHANGED (fifth blind pass, 2026-11-15, flags 19-23): the proof gets its uses, the three nodes it
% cites; (3) writes the transform of the 2 pi convention as F throughout, where it wrote the hat of
% (2) for it; Faa di Bruno's bound with G = exp names its constants at each point; the unused
% sentence "since mu-hat >= 0 is real, mu is symmetric" is dropped; the continuity of K * rho_delta
% is proved; in (4) the bumps at +-a e_1 are translates. The status line names Faa di Bruno's bound.
% CHANGED (length take (b), 2026-11-16): wording tightened, route and steps unchanged.
\begin{proof}
\leanok
\uses{thm:matern-pyramid, lem:bessel-profile, def:frame-matern-pyramid,
prop:frame-matern-pyramid-exists}
% CHANGED (blueprint round, 2026-11-16): the witness h is f, the level C is alpha, and in (3) the
% local function f of the derivative bounds is sigma; the frame pyramid's existence moves to
% prop:frame-matern-pyramid-exists, and (3) cites it; the last sentence of (4) concludes the rise at
% a rung k < n, the clause's assertion. No step changed.
\emph{(1), one local maximum.} Each of the three terms is strictly decreasing in $|x'|$ at fixed
$x_1$. On the axis, $\phi(x_1) := e^{-(x_1 - a)^2/2} + e^{-(x_1 + a)^2/2}$ has
$\phi'(x_1) = k(a - x_1) - k(a + x_1)$ with $k(y) := ye^{-y^2/2}$. The function $y \mapsto k(y) + y/2$
is nondecreasing, its derivative being $(1 - y^2)e^{-y^2/2} + 1/2 \ge 0$, since
$e^t \ge 4t - 2$ for all $t$ (from $e^{3/2} \ge 4$ and $e^{t - 3/2} \ge t - 1/2$). So
$\phi'(x_1) \le x_1$ for $x_1 \ge 0$, and $\phi' < 0$ on $(a, \infty)$, where $k(a - x_1) < 0 <
k(a + x_1)$. The third term on the axis has the derivative $-2x_1e^{-x_1^2/(8a^2)}$, which is at
most $-2 \cdot \frac{7}{8}x_1 < -x_1$ on $(0, a]$, as $e^{-u} \ge 1 - u$, and negative beyond. So $f(x_1, 0)$ is strictly decreasing in $|x_1|$, $f$ being even. From any
$x \ne 0$ one therefore strictly increases $f$ by moving $x'$ toward $0$, or $x_1$ toward $0$ when
$x' = 0$. So $0$ is the only local maximum, it is the global one, and every set $\{f > \alpha\}$ that
contains a point contains the path from it to $0$.

% CHANGED (review round 2, batch C, 2026-11-23): the route without a nonnegative transform. The old
% sentences read "Then $\widehat W = \widehat\mu\widehat B \ge 0$ is integrable and positive near
% $0$. So $W$ is continuous, tends to $0$ at infinity by dominated convergence, and by Fourier
% inversion $W(0) - W(x) = (2\pi)^{-d}\int(1 - \cos\langle\omega, x\rangle)\,\widehat
% W(\omega)\,d\omega > 0$ for $x \ne 0$. Let $m := W(0) - \max_{|x| = 1}W > 0$, and choose $a > 1$
% with $W < m/2$ on $\{|x| \ge 2a - 1\}$. With $p := ae_1$, the function $u_0 := W(\cdot - p) +
% W(\cdot + p)$ has $u_0(\pm p) \ge W(0)$, and $u_0 \le W(0) - m/2$ on the spheres of radius $1$
% about $\pm p$." and, below, "below $m/8$", "at most $W(0) - 3m/8$ on the two spheres", "the
% balls are disjoint since $a > 1$" and "$\alpha := W(0) - m/4$".
\emph{(1), two maxima.} Let
$B := e^{-|\cdot|^2/2}$ and $W := \mu * B$, which do not depend on $a$
and $\delta$. As $0 < B \le 1$, $W$ is positive, and it is continuous and tends to $0$ at infinity
by dominated convergence. Choose $r \ge 1$ with $W < W(0)/4$ on $\{|x| \ge r\}$, and put
$a := r + 1$. With $p := ae_1$, the function $u_0 := W(\cdot - p) + W(\cdot + p)$ has
$u_0(\pm p) \ge W(0)$. On the spheres of radius $r$ about $\pm p$ both of its terms are below
$W(0)/4$, as $|x \mp p| = r$ there and $|x \pm p| \ge 2a - r > r$. The
ridge $\rho_\delta$, the third term of $f$, has $0 \le \rho_\delta \le 8a^2$. For $\delta \le 1$ it
lies under $8a^2e^{-|x|^2/(8a^2)}$, and it tends to $0$ as $\delta \to 0$ off the axis $x' = 0$,
which is a null set because $d \ge 2$. So $\|\rho_\delta\|_1 \to 0$ by dominated convergence. With
% CHANGED (2026-11-23, review round 3, s4-5.S16): wording.
$K$ the density of $\mu$ and any $M > 0$, $K(y)\rho_\delta(x - y) \le M\rho_\delta(x - y) + 8a^2(K(y) - M)^+$ for all $x, y$, so $\|K * \rho_\delta\|_\infty \le M\|\rho_\delta\|_1 + 8a^2\int(K - M)^+$, and
$\int(K - M)^+ \to 0$ as $M \to \infty$, by dominated convergence. Choose $M$, then $\delta$, so
that this is at most $W(0)/4$. Then $\mu * f = u_0 + K * \rho_\delta$ is continuous, as $\rho_\delta$ is
bounded and uniformly continuous and $K$ integrable:
$|K * \rho_\delta(x) - K * \rho_\delta(x')| \le \|K\|_1\sup_y|\rho_\delta(x - y) - \rho_\delta(x' - y)|$.
As $K * \rho_\delta \ge 0$, $\mu * f$ is at least $W(0)$ at $\pm p$ and below $3W(0)/4$ on the
two spheres. So its maximum over each closed ball is attained inside, at a local maximum, and the
balls are disjoint since $a > r$. Take $\alpha := 3W(0)/4$. The points $\pm p$ lie in
$\{\mu * f > \alpha\}$, and the spheres separating them do not meet it, so it is not connected.

% CHANGED (wave 3b, 2026-11-15): clause (2) follows the machine-checked route. The decay of the
% transform is read off h = N(0, tau I) * g_0 from the proof of (4), in place of the three
% transforms computed; Plancherel's theorem is applied through the tempered distributions; g is the
% real part of the inverse transform, and eta * g = h is Fourier inversion, everywhere. The old
% proof read: "The transform of h_{a,delta} is a sum of three Gaussians, of decay at least
% e^{-epsilon_0|omega|^2} with epsilon_0 := min(1, delta^2)/2. The two bumps have the transforms
% (2 pi)^{d/2} e^{+-i a omega_1} e^{-|omega|^2/2}, and the ridge has the transform
% e^{-2a^2 omega_1^2} e^{-delta^2 |omega'|^2/2} times a constant, with 2a^2 > epsilon_0. So, taking
% epsilon = epsilon_0/2, h-hat/eta-hat is in L^1 cap L^2. Its inverse transform g is bounded and
% uniformly continuous, and in L^2 by Plancherel's theorem. It is real, since h is even, so that
% h-hat is real and even, and eta-hat is real by positivity. Finally eta * g is in L^2 with the
% transform h-hat, so eta * g = h almost everywhere, and everywhere, both being continuous."
\emph{(2)} Write $f := f_{a,\delta}$ and put $\tau := \min(1, \delta^2)/2$. As in the proof of (4)
below, $f = N(0, \tau I) * g_0$ with $g_0$ a positive sum of three Gaussians, continuous and
integrable. By Fubini, $\widehat f(\omega) = e^{-\tau|\omega|^2/2}\,\widehat{g_0}(\omega)$, so
$|\widehat f(\omega)| \le \|g_0\|_1e^{-\tau|\omega|^2/2}$. Take $\varepsilon := \tau/4$. Then
$\phi := \widehat f/\widehat\eta$ is continuous, as $\widehat\eta$ is continuous and positive, and
$|\phi(\omega)| \le \|g_0\|_1C_\varepsilon e^{-\tau|\omega|^2/4}$, so $\phi$ is bounded and lies in
$L^1 \cap L^2$. Let $G$ be the inverse transform of $\phi$, as an integral, and
$g := \operatorname{Re}G$. Then $|g| \le |G| \le (2\pi)^{-d}\|\phi\|_1$, and $g$ is continuous by
dominated convergence. It is in $L^2$ by Plancherel's theorem, read through the tempered
distributions: the inverse transform on $L^2$ of $\phi$ and $G$ have the same distribution, by
Fubini against Schwartz functions, so they agree almost everywhere. Finally
$\eta * g = \operatorname{Re}\bigl(\eta * G\bigr)$ is the real part of the inverse transform of
$\widehat\eta\,\phi = \widehat f$, which is $f$ at every point by Fourier inversion, $f$ being
continuous and integrable with $\widehat f$ integrable.

% CHANGED (wave 3b, 2026-11-15): the proof of (3) follows the machine-checked route. The old
% paragraph read "The number of i with cq^{-i}r > 1 is at most log_q(1 + cr), each such term of
% F_q(r) is at most log(1 + c^2r^2) <= 2log(1 + cr), and the others sum to at most (1 - q^{-2})^{-1}.
% So F_q(r) <= A_q(1 + log^2(1 + r)), and e^{F_q(s|omega|)} and prod_j e^{F_q(s|omega_j|)} are at
% most C_eps e^{eps|omega|^2} for every eps > 0. So (2) applies with eta = mu_s, or with its frame
% twin. Further F_q(s|omega|) = Phi(|omega|^2), where Phi(u) = sum_i log(1 + beta_i u), beta_i :=
% c^2q^{-2i}s^2, has bounded derivatives of every order k >= 1 on [0, infty), at most
% (k - 1)! sum_i beta_i^k. So every derivative of e^{Phi(|omega|^2)} is bounded by a polynomial in
% omega times e^{Phi(|omega|^2)}, and g-hat = h-hat e^{Phi(|omega|^2)} decreases rapidly with all its
% derivatives. Hence g in S(R^d), and the same holds for the frame twin." The checked route bounds
% F_q by eps|omega|^2 + A_eps directly, bounds the derivatives of each term uniformly through
% Faa di Bruno's bound at a recentred point, and builds g from Schwartz Gaussians and a multiplier of
% temperate growth, so that (2) is not invoked.
\emph{(3)} Write $b := 4\pi^2c^2s^2$ and $\beta_i := bq^{-2(i+1)}$, $i \ge 0$, so that
$F_q(2\pi s|\xi|) = S(\xi) := \sum_i\log(1 + \beta_i|\xi|^2)$, and the frame's exponent at
$2\pi\xi$ is $\sum_jS(\xi_j)$, $S$ read on the line.

\emph{Growth.} For $\varepsilon > 0$ take $N$
with $\sum_{i \ge N}\beta_i \le \varepsilon/2$. As $\log(1 + y) \le y$, the terms $i \ge N$ sum to
at most $(\varepsilon/2)|\xi|^2$; as $\log(1 + y) \le 2\sqrt y$ (from $e^{2t} \ge (1 + t)^2$), each
of the others is at most $2\sqrt b\,|\xi|$. So $S(\xi) \le \varepsilon|\xi|^2 + 2N^2b/\varepsilon$,
and the frame's exponent is at most $\varepsilon|\xi|^2 + 2dN^2b/\varepsilon$.

\emph{Derivatives.}
Fix $\xi_0$, put $P := 1 + \beta|\xi_0|^2$ and $\sigma(\xi) := (1 + \beta|\xi|^2)/P$, so that
$\log(1 + \beta|\xi|^2) = \log P + \log\sigma(\xi)$, $\sigma(\xi_0) = 1$, $|D\sigma(\xi_0)| =
2\beta|\xi_0|/P \le \sqrt\beta$, $|D^2\sigma| \le 2\beta$ and $D^k\sigma = 0$ for $k \ge 3$. Faà di
Bruno's bound, $|D^n(G \circ \sigma)(\xi_0)| \le n!\,C\,D^n$ when $|G^{(i)}(\sigma(\xi_0))| \le C$
for $i \le n$ and $|D^i\sigma(\xi_0)| \le D^i$ for $1 \le i \le n$, applied to $G = \log$ with $C = n!$ and
$D = \sqrt{2\beta}$, gives $|D^n\log(1 + \beta|\xi|^2)| \le (n!)^2(2\beta)^{n/2}$ for every
$n \ge 1$, uniformly in $\xi$. These bounds are summable in $i$, so $DS$ is the series of the
derivatives and has every derivative bounded; with $0 \le S(\xi) \le (\sum_i\beta_i)|\xi|^2$, $S$ is
of temperate growth, and so is the frame's exponent.

\emph{The input.} In this clause
$\mathcal{F}$ is the transform of the $2\pi$ convention,
\[
\mathcal{F}g(\xi) := \int e^{-2\pi i\langle\xi, x\rangle}g(x)\,dx = \widehat g(-2\pi\xi),
\]
as in the proof of Theorem~\ref{thm:matern-pyramid}(4), and
$\mathcal{F}(\mu * g)(\xi) = \widehat\mu(-2\pi\xi)\,\mathcal{F}g(\xi)$ for a finite measure $\mu$ and
an integrable $g$, by Fubini's theorem. Let $\Psi_0$ be either
exponent, so that $\widehat{\mu_s}(-2\pi\xi) = e^{-\Psi_0(\xi)}$. With $a$ and $\delta$ fixed below,
take $\tau := \min(1, \delta^2)/2$ and, as in (2), $f = N(0, \tau I) * g_0$ with $g_0$ a positive
sum of three Gaussians. Faà di Bruno's bound with $G = \exp$ at $\xi_0$,
$C := e^{\Psi(\xi_0)}$ and $D := 1 + \max_{1 \le i \le n}|D^i\Psi(\xi_0)|$ gives
$|D^ne^\Psi(\xi_0)| \le n!\,e^{\Psi(\xi_0)}D^n$, and this $D$ grows at most polynomially in
$\xi_0$ when $\Psi$ is of temperate growth. So $e^\Psi$ is a Schwartz function when $\Psi$ is of
temperate growth with $\Psi \le A - \kappa|x|^2$, $\kappa > 0$, and of temperate growth when $\Psi$
is and is bounded above. Hence each Gaussian of
$g_0$, so $g_0$ and $\mathcal{F}g_0$, is a Schwartz function, and the multiplier
$m(\xi) := e^{\Psi_0(\xi) - 2\pi^2\tau|\xi|^2}$, its exponent bounded above by the growth bound at
$\varepsilon = 2\pi^2\tau$, is of temperate growth. So $\phi := m\,\mathcal{F}g_0$ is a Schwartz
function, and so is the real part $g$ of its inverse transform. By Fubini,
$\mathcal{F}f(\xi) = e^{-2\pi^2\tau|\xi|^2}\mathcal{F}g_0(\xi)$, so
$\widehat{\mu_s}(-2\pi\xi)\,\phi(\xi) = \mathcal{F}f(\xi)$, and as in (2) $\mu_s * g$ is the real
part of the inverse transform of $\mathcal{F}f$, which is $f$.
% CHANGED (wave 3b, 2026-11-15): the frame pyramid's existence is shown, as machine-checked; the
% print described the frame pyramid without showing that its laws exist. No text was removed.
% CHANGED (blueprint round, 2026-11-16): that paragraph is the proof of
% prop:frame-matern-pyramid-exists. It read "The frame pyramid exists: if nu is a Matern pyramid
% on the line (Theorem matern-pyramid-line), the laws ... has the transform
% prod_j (1 + c^2 t^2 omega_j^2)^{-1}."
% CHANGED (wave 3b, 2026-11-15): the frame step is the law of independent coordinates of law L_{cs}
% on the line, absolutely continuous as a product of absolutely continuous laws, and the record at
% qs is read through mu_{0,qs} = mu_{0,s} * mu_{s,qs}. The sentences read "So is its frame twin, the
% product of the Laplace densities (2cs)^{-1}e^{-|x_j|/(cs)}, whose transform is
% prod_j(1 + c^2s^2omega_j^2)^{-1}. ... Then mu_s * g = h has one local maximum, and
% mu_{qs} * g = L_{cs} * h, by Theorem matern-pyramid(1), has at least two."; in the last two
% sentences L_{cs} read where mu_{s,qs} now reads, which is L_{cs} or the frame step.
The step $\mu_{s,qs} = L_{cs}$ (Theorem~\ref{thm:matern-pyramid}(1)) is symmetric
and absolutely continuous with a positive transform (Lemma~\ref{lem:bessel-profile}(1)). So is the
step of the frame pyramid (Definition~\ref{def:frame-matern-pyramid}), whose transform
$\prod_j(1 + c^2s^2\omega_j^2)^{-1}$ (Proposition~\ref{prop:frame-matern-pyramid-exists}) is that of
the law of $d$ independent coordinates, each of law $L_{cs}$ on the line (the image of $L_{cs}$ on
$\RR^1$), absolutely continuous as a product of absolutely continuous laws. So (1) applies to either
step and fixes $a$ and $\delta$. The transforms multiply, so $\mu_{0,qs} = \mu_{0,s} * \mu_{s,qs}$,
and $\mu_{qs} * g = \mu_{s,qs} * (\mu_s * g)$ for the bounded continuous $g$. Then $\mu_s * g = f$
has one local maximum, and $\mu_{qs} * g = \mu_{s,qs} * f$ has at least two. With the level
$\alpha$ of (1), $\mu_s * g - \alpha = f - \alpha$ is positive exactly on the connected set
$\{f > \alpha\}$, not empty as $f(0) = \max f \ge \max(\mu_{s,qs} * f) > \alpha$: one positive
region. And $\mu_{qs} * g - \alpha = \mu_{s,qs} * f - \alpha$ is positive exactly on
$\{\mu_{s,qs} * f > \alpha\}$, which is not connected: at least two.

\emph{(4)} Take $\mu = N(0, I)$ and fix $a$ and $\delta$ by (1). Take $n$ with
$\tau_n < \min(1, \delta^2)$. Each Gaussian of $f_{a,\delta}$ has a covariance larger than
% CHANGED (2026-11-23, review round 3, s4-5.S16): wording.
$\tau_nI$: the bumps have $I$, and the ridge $\mathrm{diag}(4a^2, \delta^2, \dots, \delta^2)$. Write $g_\Sigma$ for the centred Gaussian density with covariance $\Sigma$. For a diagonal covariance $\Sigma = \mathrm{diag}(v)$ and $\tau < \min_iv_i$, the integral
$N(0, \tau I) * g_{\Sigma - \tau I}$ factors over the coordinates into one-dimensional Gaussian
integrals, and completing the square gives $g_\Sigma$. The bumps are the translates of $g_I$ by
$\pm ae_1$, and convolution commutes with translations, so each is $N(0, \tau I)$ convolved with
the same translate of $g_{(1 - \tau)I}$. So
$f_{a,\delta} = N(0, \tau_nI) * g$ with $g$ a positive sum of three Gaussians. Since
$N(0, (\tau_n + 1)I) = N(0, \tau_nI) * N(0, I)$ (the transforms), the record of $g$ at
variance $\tau_n + 1 = q^{2n}\tau_n$ is $\mu * f_{a,\delta}$, with at least two local maxima. The
records at the variances $q^{2k}\tau_n$, $k = 0, \dots, n$, are the rungs of the ladder of ratio
$q$.
% CHANGED (review round 2, batch C, 2026-11-23): the least rung with two maxima, and the rung below
% it. The old sentence read "If the count did not rise from rung $k$ to rung $k + 1$ for any
% $k < n$, it would be nonincreasing from rung $0$ to rung $n$, and the $n$-th record would have one
% local maximum."
Rung $0$ has one local maximum and rung $n$ at least two, so there is a least rung $j \ge 1$ with
at least two, and rung $k := j - 1$ has at most one. Rung $k$ is
$N(0, (q^{2k} - 1)\tau_nI) * f_{a,\delta}$, by the same identity of Gaussians. It is continuous and
positive, and it tends to $0$ at infinity by dominated convergence, as $f_{a,\delta}$ is bounded by
$f_{a,\delta}(0)$ and each of its three Gaussians tends to $0$. So it attains its
maximum, at a local maximum, and rung $k$ has exactly one.

% CHANGED (2026-11-23, review round 3, s4-5.S16): the dilation factor is beta, freeing c.
\emph{(5)} Take $g$, $n$ and $k$ of (4), put $v := q^{2k}\tau_n$ and $\beta := (v/\tau)^{1/2}$, and
let $g_\beta(y) := g(\beta y)$. The image of $N(0, sI)$ under $y \mapsto \beta y$ is $N(0, \beta^2sI)$ (the
transforms), so
\[
  (N(0, sI) * g_\beta)(x) = \int g(\beta x - \beta y)\,N(0, sI)(dy) = (N(0, \beta^2sI) * g)(\beta x)
\]
for $s \ge 0$. The map $x \mapsto \beta x$ is a homeomorphism, so $x$ is a local maximum of
$x \mapsto F(\beta x)$ exactly when $\beta x$ is one of $F$, and the counts agree. At $s = \tau$,
$\beta^2\tau = v$, and at $s = q^2\tau$, $\beta^2q^2\tau = q^{2(k+1)}\tau_n$. So the records of the
positive input $g_\beta$ at $\tau$ and $q^2\tau$ have the counts of rungs $k$ and $k + 1$ in (4).
\end{proof}

% draft: the comment after Proposition 4.6, line 628.
Clause (1) applies to every kernel $\mu_t$ and every octave step $L_{ct}$ of the \Matern{}
pyramid, and to the Gaussian. Clauses (3) and (4) refute, as counts of extrema, the rung-to-rung
count of Proposition~\ref{prop:ladder-non-creation-line} carried to the plane. There ``zero
crossings'' have no single meaning, and the positive regions are the reading treated here. In
clause (3) the input $g$ need not be positive. Adding a constant makes it positive and bounded, and
changes each record by that constant.
% CHANGED (review round 2, batch C, 2026-11-23): the old sentences read "Clause (4) does not say
% which rung raises the count. The reciprocal ... does not apply to it. Which rung of the Gaussian
% ladder raises the count is open. Nothing here uses it." Clause (4) now places the rise at a rung
% with one local maximum, which answers the question.
Clause (4) places the rise after a rung with one local maximum, and clause (5) moves that rung to
any prescribed variance by a dilation of the input. Clause (2) does not give this directly. The
reciprocal $e^{b|\omega|^2}$ of a Gaussian transform grows faster than
$C_\varepsilon e^{\varepsilon|\omega|^2}$ for $\varepsilon < b$, so it does not apply to the
Gaussian.

% draft: Proposition 4.7 (section 4.2), lines 632-638 of draft/spatial-hemigroup-pyramids.md.
% CHANGED: renamed N to \mathcal{N} and "N-maximum" to "\mathcal{N}-maximum" (BLUEPRINT-INPUTS).
% CHANGED: clause (3) writes out the level-k record (f_j smoothed by the dilated mask and read on
% the coarser grid), which the draft names only; the reading is the draft's proof's.
% CHANGED: clause (3) quantifies k, every integer k >= 0 (second blind review, flag 9).
\begin{proposition}[creation on a grid]
\label{prop:creation-grid}
% CHANGED (wave 3b, 2026-11-15): every declaration is machine-checked, so statement and proof are
% marked proved. Clause (3) reads the existence and uniqueness of the generating triplet (A39), which
% the old proof used without citing it; the node is now [A] with A39 in its uses. The old status line
% read "Cites nothing of its own: clause (3) reads Lemma lattice-infinitely-divisible, which carries
% its entries. The rest is tightness and uniform continuity."
\lean{Pyramids.creation_grid_1, Pyramids.creation_grid_2, Pyramids.creation_grid_3}
\leanok
\uses{lem:lattice-infinitely-divisible, lem:symmetric-levy-khintchine, thm:cited-A39}
% CHANGED (2026-11-15, a strengthening)
% The hypothesis on the transform is dropped: no clause reads it. The old wording: "let
% $w \ne \delta_0$ be a symmetric probability law on $h\ZZ^d$ with $\widehat w \ge 0$, and let".
Let $d \ge 2$, let $w \ne \delta_0$ be a symmetric probability law on $h\ZZ^d$, and let
$\mathcal{N} \subset h\ZZ^d \setminus \{0\}$ be finite and symmetric and
contain the $2d$ axis neighbours. Call $x$ an \emph{$\mathcal{N}$-maximum} of $u$ if
$u(x) \ge u(x + y)$ for all $y \in \mathcal{N}$.
\begin{enumerate}
\item There is a summable $f > 0$ on $h\ZZ^d$ with exactly one $\mathcal{N}$-maximum such that
$w * f$ has at least two.
\item For an integer $q \ge 2$, $f$ can be chosen so that $(w * f)|_{qh\ZZ^d}$ has at least two
$q\mathcal{N}$-maxima as well.
\item If $w$ is infinitely divisible and $q \ge 2$ is an integer, then for every integer $k \ge 0$
every bounded array on $q^kh\ZZ^d$ is the level-$k$ record, smoothed and decimated $k$ times, of a bounded input on
$h\ZZ^d$: with $f_0$ the input, $f_j := \bigl((D_{q^{j-1}}w) * f_{j-1}\bigr)|_{q^jh\ZZ^d}$ and the
array $f_k$.
\end{enumerate}
% CHANGED (2026-11-15, a strengthening)
% The declarations require less than the print, which keeps its setting: (1) and (2) do not require N finite,
% symmetric or free of 0, only that it contain the axis neighbours; (2) and (3) hold for every
% integer q >= 1; (3) does not require d >= 2 or w other than delta_0. The status line had no such
% sentence before.
% CHANGED (fidelity review, 2026-11-15, R26/R50/R69): the status line lists the declarations' hypotheses beside the node's (LINKAGE rule 6).
\statusA\quad\ledger{A39}\quad\emph{(\textbf{Assignment.} \ledger{A39},
[\sscite{sato1999levy}, Thm.~8.1(i)--(ii), pp.~37--38], carries one step of (3), through
Lemma~\ref{lem:symmetric-levy-khintchine}: a symmetric infinitely divisible law has symmetric data
$(A, \nu)$. It does not carry the lattice: $A = 0$ and $\nu$ finite on
$h\ZZ^d \setminus \{0\}$ are Lemma~\ref{lem:lattice-infinitely-divisible}'s, which carries its own
entries; nor the inverse in $\ell^1$, proved here. Clauses (1) and (2) cite
nothing: they are the maximum of an even function that tends to $0$, and the tightness of $w$.
Machine-checked. The declarations require less than the print, which keeps its setting.
\textbf{Hypotheses} (the node's, each with what its declarations require): the lattice $h\ZZ^d$: in
(1) and (2) absent, the declarations work in the index coordinates $\ZZ^d$, which the scaling by
$h$ identifies with $h\ZZ^d$; in (3) same, $h > 0$, read through the law of $w$ on $h\ZZ^d$;
$d \ge 2$: same in \texttt{creation\_grid\_1} and \texttt{\_2}; weaker in \texttt{\_3}, which does
not require it; $w$ a symmetric probability law: same, \texttt{IsProbKernel w} (nonnegative, summable,
of sum $1$) and $w(-n) = w(n)$; $w \ne \delta_0$: same in \texttt{\_1} and \texttt{\_2}; weaker in
\texttt{\_3}, which does not require it; $\mathcal{N}$ finite: weaker, not required; $\mathcal{N}$
symmetric: weaker, not required; $0 \notin \mathcal{N}$: weaker, not required; $\mathcal{N}$ contains
the $2d$ axis neighbours: same in \texttt{\_1} and \texttt{\_2}; $\mathcal{N}$ is not read in (3);
(2) and (3) $q \ge 2$ an integer: weaker, $q \in \NN$ with $q \ge 1$; (3) $w$ infinitely
divisible: same, \texttt{IsInfinitelyDivisible} of the law of $w$ on $h\ZZ^d$; $k \ge 0$ an
integer: same; the array bounded: same.)}
\end{proposition}

% CHANGED (wave 3b, 2026-11-15): the proof follows the machine-checked route. The old proof of
% (1)-(2) used the bound P <= P(0) of the line marginal by Fourier inversion (which needed the
% transform nonnegative), a Gaussian profile, the limit profile V = phi + p_0 T, and an interior
% maximum in the box Q_R. The checked route uses only that some axis line has mass p < 1, a profile
% linear on [-1, 1], and the global maximum of the even output, compared at two points. The old proof
% of (3) wrote w = delta_{z_0} * pi_w with a drift z_0 and inverted pi_w; the checked route uses the
% symmetric data (no drift) and the uniqueness of the coefficients of a summable kernel to conclude
% w * m = delta_0.
% CHANGED (fifth blind pass, 2026-11-15, flags 24-28): in (3) the symmetric data are derived from
% A39(i)-(ii) by the reflection, printed, where they were attributed to A39 directly, and the data
% are named in this paper's normalization, A half of Sato's matrix; the circular sentence "the
% transform of w is real, as w is symmetric, so the law w is symmetric" is dropped; the axis of mass
% below one is the j-th, the argument written for j = 1, where it was relabelled Z e_1; the
% induction names the level-(k-1) array it is applied to.
% CHANGED (review response r1, batch 2b, 2026-11-23, finding B2): the data of the symmetric
% infinitely divisible w are lem:symmetric-levy-khintchine, cited in (3) as exists_levyKernel calls
% exists_symData, where the triplet, its reflection and the odd part were written out. No step
% changed.
\begin{proof}
\leanok
\uses{lem:lattice-infinitely-divisible, lem:symmetric-levy-khintchine, thm:cited-A39}
Take $h = 1$, and write $|y|_1 := \sum_j|y_j|$.

\emph{An axis line of mass below one.} If every axis line $\ZZ e_j$ had mass $1$, $w$ would be
carried by their intersection $\{0\}$, as $d \ge 2$, and $w = \delta_0$. So some axis line
$\ZZ e_j$ has mass $p < 1$. The argument is written for $j = 1$; for another $j$ it is the same with
the coordinates $x_1$ and $x_j$ exchanged in the witness below, $\mathcal{N}$ entering only through
the $2d$ axis neighbours it contains. Put $\theta := 2/(1 + p)$, so that $\theta > 1$ and
$1 - \theta p = (1 - p)/(1 + p) > 0$.

\emph{The witness.} Write $x = (x_1, x')$ and $|x'|_1 := \sum_{j \ge 2}|x_j|$. Let
$\Psi(t) := \min(1 + |t|, 2e^{1 - |t|})$. It is $1 + |t|$ for $|t| \le 1$ and $2e^{1 - |t|}$
otherwise, it lies in $(0, 2]$, and $|\Psi(s) - \Psi(t)| \le 6|s - t|$. For an integer $R \ge 1$ put
\[
  b(x) := \Psi(x_1/R)\,e^{-|x'|_1/R}, \qquad
  f(x) := b(x) + \mathbf{1}_{\{x' = 0\}}\,\theta\max(1 - |x_1|/R, 0) .
\]
It is positive and even, and summable, as $f(x) \le (2 + \theta)e\,e^{-|x|_1/R}$.

\emph{One $\mathcal{N}$-maximum.} First $f(0) = 1 + \theta \ge f(x)$ for every $x$. Off the line
$f = b \le 2$. On it, $f(x) = 1 + \theta - (\theta - 1)|x_1|/R$ when $|x_1| \le R$, and $f \le 2$
beyond. So $0$ is an $\mathcal{N}$-maximum. From every $x \ne 0$ one axis step strictly raises $f$.
If $x' \ne 0$, a step toward the line in a coordinate $j \ge 2$ with $x_j \ne 0$ lowers $|x'|_1$ by
one, which raises $b$, and the ridge term cannot fall from $0$. If $x' = 0$ and $x_1 \ne 0$, the step
toward $0$ along the line raises $f$ by $(\theta - 1)/R$ while $|x_1| \le R$. Beyond, it raises $\Psi$
strictly, as $|x_1| - 1 \ge R$ (both are integers), and the ridge term is $0$ at $x$. As
$\mathcal{N}$ contains the axis neighbours, no $x \ne 0$ is an $\mathcal{N}$-maximum.

\emph{Two for the output.} The output $u := w * f$ is even, as $w$ is symmetric and $f$ even. It is
positive, and it tends to $0$ at infinity, by dominated convergence, as $f$ does. So $u$ attains its
maximum, and a maximizer is an $\mathcal{N}$-maximum for every $\mathcal{N}$. If $u(z) > u(0)$ for
some $z$, a maximizer $x^* \ne 0$ exists, and $-x^* \ne x^*$ is a second one. Now
$|b(x - y) - b(x)| \le \min(6|y|_1/R, 2)$, so
$\varepsilon(R) := \sum_yw(y)\min(6|y|_1/R, 2)$ bounds $|w * b - b|$ at every point, and
$\varepsilon(R) \to 0$ as $R \to \infty$ by dominated convergence: this is the tightness of $w$. The
ridge term of $u$ at $0$ is at most $\theta$ times the mass $p$ of the line. So
\[
  u(0) \le b(0) + \varepsilon(R) + \theta p = 1 + \varepsilon(R) + \theta p, \qquad
  u(Re_1) \ge b(Re_1) - \varepsilon(R) = 2 - \varepsilon(R),
\]
and $u(Re_1) - u(0) \ge 1 - \theta p - 2\varepsilon(R) > 0$ for large $R$. For (2), take $R = qK$
with $K$ large. Then $Re_1 = q(Ke_1)$ is on $q\ZZ^d$, and the same argument applies to the even
function $m \mapsto u(qm)$, which tends to $0$ at infinity.

\emph{(3)} Being symmetric and infinitely divisible, $w$ has \Levy{} data $(A, \nu)$ and the
generating triplet $(2A, \nu, 0)$, by Lemma~\ref{lem:symmetric-levy-khintchine}:
$\widehat w(\omega) = e^{-\langle\omega, A\omega\rangle - \int(1 - \cos\langle\omega, x\rangle)\,\nu(dx)}$.
Being carried by $\ZZ^d$, it has $A = 0$, and $\nu$ is finite and carried by
$\ZZ^d \setminus \{0\}$ (Lemma~\ref{lem:lattice-infinitely-divisible}). Put $a(n) := \nu(\{n\})$, a
summable symmetric kernel, and $M := \sum_na(n)$. Then
$\widehat w(\vartheta) = e^{-\sum_na(n)(1 - \cos\langle n, \vartheta\rangle)} = e^{-M}e^{\widehat a(\vartheta)}$,
the transform $\widehat a$ being real. Put $m := e^M\sum_{k \ge 0}(-1)^ka^{*k}/k!$. The series
converges absolutely in $\ell^1(\ZZ^d)$, with norm at most $e^{2M}$, so
$\widehat m = e^Me^{-\widehat a} = 1/\widehat w$ and $\widehat{w * m} = 1 = \widehat{\delta_0}$. A
summable kernel is determined by its Fourier coefficients, so $w * m = \delta_0$. At $k = 0$ the
array is its own input. Let $k \ge 1$, and suppose that every bounded array on $q^{k-1}\ZZ^d$ is the
level-$(k-1)$ record of a bounded input. Given a bounded array on $q^k\ZZ^d$, extend it by $0$ to a
bounded $u$ on $q^{k-1}\ZZ^d$, and put $v := (D_{q^{k-1}}m) * u$, a bounded array on
$q^{k-1}\ZZ^d$, as $D_{q^{k-1}}m$ is summable. By the induction hypothesis $v$ is the level-$(k-1)$
record $f_{k-1}$ of a bounded input $f_0$. Its smoothing $(D_{q^{k-1}}w) * v$ is
$(D_{q^{k-1}}(w * m)) * u = u$, by associativity, so
$f_k = \bigl((D_{q^{k-1}}w) * f_{k-1}\bigr)|_{q^k\ZZ^d} = u|_{q^k\ZZ^d}$ is the given array.
\end{proof}

% draft: the paragraph after Proposition 4.7 (line 640) moves to chapter 5, beside
% Propositions 5.4 and 5.6 (adr/0004, decision 3).
Chapter~5 reads the proposition on the masks of the grid pyramids it constructs: the discrete
Helmholtz mask, the frame mask and the binomial masks.

\section{The selection}
\label{sec:matern-selection}

% draft: the lead of section 4.3, line 660.
Counting is gone in the plane. What remains of non-creation is non-enhancement: smoothing does not
raise a local maximum nor lower a local minimum. The classical form requires $\partial_tu \le 0$ at a
local maximum of the record, which needs a derivative in scale
\prog{V-selection v0.1}{Thm.~4.2(3), p.~31}. Read at one finite step, as the difference quotient
$(u_t - u_s)/(t - s) \le 0$ at a local maximum of the later record, it is the implicit Euler
discretization of that axiom. It needs no continuum of scales, so it can be required of the octave
stage of a pyramid. The two forms select different families. On the line, the derivative form required
at every scale selects the Gaussian ray \prog{V-selection v0.1}{Thm.~4.2(3) and Cor.~4.4, pp.~30,
33}. The finite-step form selects the resolvent and rejects every Gaussian step
(Proposition~\ref{prop:stage-lemma} and Theorem~\ref{thm:selection-stratum-two}(4)).

% draft: Definition 4.8 (section 4.3), line 662 of draft/spatial-hemigroup-pyramids.md.
% CHANGED: the grid form is written out, where the draft sends it to Proposition 5.11: the principle
% on a finite symmetric neighbourhood, required of bounded inputs, as Proposition 5.11 requires it.
% CHANGED (blueprint round, 2026-11-16, review flag "A local maximum is meant in the weak sen"): the
% definition says "weak local maximum" and defines it, as the predicate reads it (IsLocalMax); the
% paper's remark after the definition, which said so, is dropped. The object is the same.
\begin{definition}[the implicit-Euler maximum principle]
\label{def:implicit-euler-principle}
\lean{Pyramids.HasImplicitEulerPrinciple, Pyramids.IsNMax, Pyramids.IsNMin,
Pyramids.HasGridImplicitEulerPrinciple}\leanok
A symmetric probability measure $\mu$ on $\RR^d$ has the \emph{implicit-Euler maximum principle}
if for every real $f \in \mathcal{S}(\RR^d)$ the output $u := \mu * f$ has $u(x) \le f(x)$ at every
weak local maximum $x$ of $u$, and $u(x) \ge f(x)$ at every weak local minimum. Here $x$ is a
\emph{weak local maximum} of $u$ if $u(x) \ge u(y)$ for every $y$ in a neighbourhood of $x$, and a
\emph{weak local minimum} if it is a weak local maximum of $-u$. A stage $\mu_{s,t}$ of a
family has the principle if the measure $\mu_{s,t}$ has it. On a lattice $h\ZZ^d$, with
$\mathcal{N} \subset h\ZZ^d \setminus \{0\}$ finite and symmetric, a symmetric probability $\kappa$
on $h\ZZ^d$ has the principle \emph{on $\mathcal{N}$} if for every bounded real $f$ on $h\ZZ^d$ the
output $u := \kappa * f$ has $u(x) \le f(x)$ at every $\mathcal{N}$-maximum $x$ of $u$ and
$u(x) \ge f(x)$ at every $\mathcal{N}$-minimum. Here $x$ is an \emph{$\mathcal{N}$-maximum} of $u$ if
$u(x) \ge u(x + y)$ for every $y \in \mathcal{N}$, and an \emph{$\mathcal{N}$-minimum} if it is an
$\mathcal{N}$-maximum of $-u$.
\statusT\quad\emph{(Vocabulary. Proposition~\ref{prop:grid-implicit-euler-principle} characterizes
the principle on a finite symmetric neighbourhood $\mathcal{N} \subset \ZZ^d \setminus \{0\}$, on
$\ZZ^d$ for every $d \ge 1$, and Definition~\ref{def:neighbourhood-principle}
extends the definition to real kernels of mass one, the signed principle.
\texttt{HasGridImplicitEulerPrinciple N $\kappa$} requires nothing of $\mathcal{N}$ (finiteness, symmetry,
$0 \notin \mathcal{N}$) or of $\kappa$ (symmetry, summability, positivity), and reads $\kappa * f$ as
a \texttt{tsum}, the junk value $0$ where it diverges; the statements that use it carry those
hypotheses, $\kappa$ summable among them.)}
% CHANGED (fidelity review, 2026-11-15, R10): the status line says what the grid predicate does not
% require and that its uses carry it.
\end{definition}

% draft: the sentence after Definition 4.8, line 664.
The inputs are signed, so the clause at minima is the clause at maxima applied to $-f$.

% draft: Proposition 4.9 (section 4.3), lines 666-670 of draft/spatial-hemigroup-pyramids.md. The
% stage lemma of module C (spatial-hemigroup-scale-space, prop:stage-lemma, clauses (1), (2) and
% (4)), in a version not yet released; restated with its proof, the converse's included
% (adr/0004, decision 1). Not cited: the module is not released.
% TODO(module C): cite prop:stage-lemma when released, or keep the proof (adr/0004, Consequences).
% CHANGED: clause (2) names the range of its parameter, tau >= 0, as clause (1) does.
% CHANGED (fifth blind pass, 2026-11-15, flag 30): the proof of (1) reads the law L_1 and its
% transform from lem:bessel-profile(1), which joins the uses of statement and proof and the status
% line; the status line said "Cites nothing".
\begin{proposition}[the stage lemma]
\label{prop:stage-lemma}
\lean{Pyramids.stage_lemma_1, Pyramids.stage_lemma_1_helmholtz, Pyramids.stage_lemma_2,
Pyramids.stage_lemma_3, Pyramids.stage_lemma_3_rotation}
% CHANGED (wave 3b, 2026-11-15): all five declarations are machine-checked, so statement and proof
% are marked proved. The line marked the node not ready.
\leanok
\uses{def:implicit-euler-principle, lem:bessel-profile}
Let $\mu$ be a symmetric probability measure on $\RR^d$.
\begin{enumerate}
\item \emph{(Sufficiency.)} For every positive semidefinite matrix $C$ there is a symmetric
probability measure with transform $(1 + \omega\cdot C\omega)^{-1}$, and it has the principle. In
particular the Helmholtz resolvent $(I - \tau\Delta)^{-1}$ has it for every $\tau \ge 0$, in every
dimension.
\item \emph{(The grid.)} On $h\ZZ^d$ the resolvent $\kappa$ of the $(2d+1)$-point Laplacian, with
transform $\bigl(1 + 4\tau h^{-2}\sum_{j=1}^d\sin^2(h\omega_j/2)\bigr)^{-1}$, $\tau \ge 0$, is a
symmetric probability and has the principle on the $2d$ axis neighbours.
\item \emph{(Converse.)} If $\mu$ has the principle, then $1/\widehat\mu(\omega) = 1 + \omega\cdot
C\omega$ for a positive semidefinite $C$. If $\mu$ is rotation invariant, $C = \tau I$ and $\mu$ is
the Helmholtz resolvent, the identity when $\tau = 0$.
\end{enumerate}
% CHANGED (wave 3b, 2026-11-15): all five declarations are machine-checked. The previous
% status read: "Cites nothing: the uniqueness of the Fourier transform, Fourier inversion,
% Taylor's theorem and dominated convergence. Clauses (1) and (2), three declarations, are
% machine-checked. Clause (3), two declarations, is not, and is priced at about 3700 to 5500 lines:
% Step 1, a functional carried by a point, about 900 to 1400; Step 2 about 1200 to 1800; the jets
% and the annihilator of Step 3 about 700 to 1000; Step 4 about 600 to 900; the growth of the
% reciprocal that Step 2 requires about 150 to 300; the rotation-invariant case about 150 to 250." The
% status before wave 3b: "None of the five declarations is machine-checked yet. Clauses (1) and (2)
% need a second-order identity for mu * f and a second-derivative test at a local maximum, and on
% the grid the uniqueness of Fourier coefficients on Z^d. Clause (3) needs Step 1, the description
% of a functional carried by a point, and the jets of Step 3, which the library does not have."
% CHANGED (fidelity review, 2026-11-15, R26/R50/R69): the status line lists the declarations' hypotheses beside the node's (LINKAGE rule 6).
\statusT\quad\emph{(Cites nothing beyond Lemma~\ref{lem:bessel-profile}(1), the law $L_1$ and its
transform $(1 + |\omega|^2)^{-1}$, read in (1): the uniqueness of the Fourier transform, Fourier
inversion, the bounds $1 - y^2/2 \le \cos y \le 1 - y^2/2 + y^4$ near $0$, and dominated
convergence. All five declarations are machine-checked. The predicate
\texttt{HasImplicitEulerPrinciple} is the condition of Definition~\ref{def:implicit-euler-principle}
alone, at every weak local maximum and minimum (\texttt{IsLocalMax}, \texttt{IsLocalMin}) of the
output for every real Schwartz input; the symmetry and the probability of the definition are
carried here as separate hypotheses or conjuncts. \textbf{Hypotheses} (the node's, each with what
its declarations require): $\mu$ a symmetric probability measure on $\RR^d$: in (1) the measure is
produced, symmetric and a probability, by \texttt{stage\_lemma\_1}, whose principle conjunct and
\texttt{stage\_lemma\_1\_helmholtz} require only a probability with the transform: weaker, the symmetry
following from the real transform; same in \texttt{stage\_lemma\_3} and \texttt{\_3\_rotation};
(1) $C$ positive semidefinite: same; $\tau \ge 0$: same; every dimension: same, no hypothesis on
$d$; (2) the lattice $h\ZZ^d$, $h > 0$: same, in the index coordinates $\ZZ^d$ with $\theta = h\omega$;
$d \ge 1$: same; $\tau \ge 0$: same; the resolvent $\kappa$: same, every probability kernel with the
transform of the clause, one of which the declaration gives; (3) $\mu$ has the principle: same;
$\mu$ rotation invariant: same, \texttt{IsRotationInvariant}, under the orthogonal group;
$d \ge 1$: same in \texttt{\_3\_rotation}, weaker in \texttt{stage\_lemma\_3}, which does not
require it.)}
% CHANGED (wave 3b, 2026-11-15): the proof of (3) is no longer module C's. The old line: "The proof
% of (3) is module C's, carried over with its symbols renamed, and refereed once, blind
% (records/pyramids/REFEREE-draft-sections-3-4.md, item 5); no node of module C is used."
% The proof of (3) follows the machine-checked route, shorter than module C's; no node of module C
% is used.
\end{proposition}

% CHANGED (length take (b), 2026-11-16): wording tightened, route and steps unchanged.
\begin{proof}
\leanok
\uses{def:implicit-euler-principle, lem:bessel-profile}
% CHANGED: the proof of (2) is written out on h Z^d (the random-walk mixture and the maximum over
% the axis neighbours), where the draft cites Propositions 5.6(4) and 5.11 of chapter 5; the
% argument is the one they give, without the citation ahead.
% CHANGED (wave 3b, 2026-11-15): the proof of (1) follows the machine-checked route, which takes no
% derivative of u. The old proof took the law as the mixture of N(0, 2rC) against e^{-r} dr, showed
% u - sum_{i,j} C_ij d_i d_j u = f through the transform, and used tr(CH) <= 0 for the negative
% semidefinite Hessian H of u at a local maximum.
% CHANGED (fifth blind pass, 2026-11-15, flags 29, 31-33): in (1) the transform of f is defined,
% with the sign e^{-i}, under which the inversion displays hold (under the sign of a measure's
% transform they gave u(-z)), and the law of (1) is mu_C, where it reused the statement's mu; in
% Step 1 the parameter s is chosen in (0, 1]; Steps 1-3 say which pairs they apply Step 1 to.
% CHANGED (blueprint round, 2026-11-16, review flag "the sign opposite to that of a measure's"): one
% transform for functions across the chapter, f-hat(omega) := int e^{i omega.x} f(x) dx, the sign of
% a measure's transform; the inversion displays read e^{-i omega.z}. The 2 pi convention keeps the
% distinct symbol \mathcal{F}. The route is unchanged.
\emph{(1)} Write a positive semidefinite $C$ as $B^{\tp}B$, $B$ real, and put
$b_k := B^{\tp}e_k$, so that $\omega\cdot C\omega = |B\omega|^2 = \sum_k(\omega\cdot b_k)^2$. The
image $\mu_C$ of $L_1$ under $y \mapsto B^{\tp}y$ is a probability measure with the transform
$\widehat{L_1}(B\omega) = (1 + \omega\cdot C\omega)^{-1}$ (Lemma~\ref{lem:bessel-profile}(1)), which
is even, so $\mu_C$ is symmetric by the uniqueness of the transform.

Let $f \in \mathcal{S}(\RR^d)$ be real, and put $u := \mu_C * f$. The transform of $f$ is
$\widehat f(\omega) := \int e^{i\omega\cdot x}f(x)\,dx$, the transform of the measure $f(x)\,dx$, so
that Fourier inversion reads
$f(z) = (2\pi)^{-d}\int e^{-i\omega\cdot z}\widehat f(\omega)\,d\omega$. By Fubini's theorem $u$ has
the transform $\widehat{\mu_C}(\omega)\widehat f(\omega) =: g(\omega)$. As
$0 < \widehat{\mu_C} \le 1$ and $\widehat f$ is rapidly
decreasing, $g$ and $|\omega|^2g$ are integrable, and by Fourier inversion
\[
  u(z) = (2\pi)^{-d}\int e^{-i\omega\cdot z}g(\omega)\,d\omega, \qquad
  f(z) = (2\pi)^{-d}\int e^{-i\omega\cdot z}(1 + \omega\cdot C\omega)\,g(\omega)\,
  d\omega .
\]
These integrals are real, as $u$ and $f$ are; the real part written below only makes that
visible. So $f(x) - u(x) = (2\pi)^{-d}\sum_kQ_{b_k}$ with $Q_v := \operatorname{Re}\int e^{-i\omega\cdot x}
(\omega\cdot v)^2g(\omega)\,d\omega$. For a vector $v$ and $s \ne 0$,
\[
  \frac{u(x + sv) + u(x - sv) - 2u(x)}{s^2} = (2\pi)^{-d}\operatorname{Re}\int e^{-i\omega\cdot x}\,
  \frac{2\cos(s\,\omega\cdot v) - 2}{s^2}\,g(\omega)\,d\omega .
\]
The fraction in the integrand tends to $-(\omega\cdot v)^2$ as $s \to 0$ and is bounded by
$(\omega\cdot v)^2 \le |v|^2|\omega|^2$, as $1 - \cos y \le y^2/2$, so by dominated convergence the
quotient tends to $-(2\pi)^{-d}Q_v$. At a local maximum $x$ of $u$ its numerator is $\le 0$ for
small $s$, so $Q_v \ge 0$ for every $v$, and $f(x) \ge u(x)$. Minima follow with $-f$. The
Helmholtz resolvent is $C = \tau I$.

\emph{(2)} At $\tau = 0$, $\kappa = \delta_0$. Let $\tau > 0$, put $a := \tau h^{-2}$ and
$p := 2da/(1 + 2da) \in (0, 1)$, and let $\sigma$ be the simple random walk step, of mass $1/(2d)$ at
each $\pm he_j$, with transform $\widehat\sigma(\omega) = d^{-1}\sum_j\cos(h\omega_j)$. The
measure $(1 - p)\sum_{k \ge 0}p^k\sigma^{*k}$ is a symmetric probability on $h\ZZ^d$, a convex
combination of such, with the transform $(1 - p)/(1 - p\widehat\sigma)$. Its reciprocal is
$1 + \frac{p}{1 - p}(1 - \widehat\sigma) = 1 + 2da\cdot d^{-1}\sum_j(1 - \cos h\omega_j) =
1 + 4\tau h^{-2}\sum_j\sin^2(h\omega_j/2)$. So it is $\kappa$, by the uniqueness of Fourier
coefficients.
% CHANGED (wave 3b, 2026-11-15): the identity for kappa is read from its series and Delta_h is moved
% onto kappa by shifting the summation index, as the machine-checked proof does. The old text: the
% convolution with a signed measure lambda of transform -4h^{-2} sum_j sin^2(h omega_j/2), and
% (delta_0 - tau lambda) * kappa = delta_0 because its transform is 1.
Let $\Delta_hu(n) := h^{-2}\sum_j\bigl(u(n + he_j) - 2u(n) + u(n - he_j)\bigr)$. The series gives
$\kappa * \sigma = (\kappa - (1 - p)\delta_0)/p$, and $\Delta_h\kappa = 2dh^{-2}(\kappa * \sigma -
\kappa)$, so $\kappa - \tau\Delta_h\kappa = \delta_0$, since $(1 - p)(1 + 2da) = 1$. For bounded real
$f$, $u := \kappa * f$ is bounded and $\Delta_hu = (\Delta_h\kappa) * f$, by shifting the summation
index, $\kappa$ being summable. So $u - \tau\Delta_hu = f$. At a maximum $x$ of $u$ over the axis
neighbours every term $u(x \pm he_j) - u(x)$ is nonpositive, so $\Delta_hu(x) \le 0$ and
$f(x) \ge u(x)$. Minima follow with $-f$.

\emph{(3)}
% CHANGED (wave 3b, 2026-11-15): the proof of (3) follows the machine-checked route, which tests the
% principle on cosine sums under a wide Gaussian and takes no derivative of an output. The old proof
% (module C's, refereed blind in REFEREE-draft-sections-3-4.md, item 5) had four steps: Step 1, a
% linear functional on C_c^infty vanishing on functions that vanish near 0 and bounded by
% sup-norms of derivatives up to order n is sum c_beta d^beta phi(0) (Taylor, a cut-off
% psi(./epsilon), Leibniz); Step 2, for 1/mu-hat smooth of polynomial growth, the multiplier L
% with symbol m = 1/mu-hat - 1 and f := u + L u give L u >= 0 at maxima of u, L is local (through
% zeta = |y|^2 psi), Step 1 makes m a polynomial along phi_R(y) = e^{i omega_0 y} psi(y/R), and
% homogeneous test polynomials of degree k >= 3 and the sign leave m = omega.C omega with C
% positive semidefinite; Step 3, the 2-jets at 0 of outputs of inputs vanishing near 0 form a
% subspace missing (a_0, 0, H) with a_0 > 0 and H negative definite, an annihilating
% (alpha, xi, Theta) gives a second-order equation off the support of the input, and Step 1 at
% n = 2 with phi_R gives (alpha + i xi.omega - omega.Theta omega) mu-hat = a polynomial of degree
% two; Step 4, averaging and the projections mu_v onto lines, with the Laplace law on the line and
% the density of the directions outside U_0, gave mu-hat = (1 + omega.Y omega)/(1 + omega.B omega)
% with B, Y positive semidefinite, so that Step 2 applied. The rotation-invariant case read the
% radial form omega.C omega = tau|omega|^2. In that proof the draft's functional Upsilon was Phi,
% and in Step 3 the draft's reused Upsilon and c_beta were Omega and pi_beta
% (REFEREE-draft-sections-3-4.md, item 5, findings 4 and 5).
Write $\rho := \widehat\mu$. It is real, since $\mu$ is symmetric, continuous, and $\rho(0) = 1$.
The sine integrates to $0$ against $\mu$, so for $w, x \in \RR^d$
\[
  \int\cos\langle w, x - y\rangle\,\mu(dy) = \rho(w)\cos\langle w, x\rangle .
\]
The principle is tested on the inputs $\sum_ka_k\cos\langle w_k, x\rangle e^{-s|x|^2}$, whose
outputs tend as $s \to 0$ to $\sum_ka_k\rho(w_k)\cos\langle w_k, x\rangle$. Step~1 turns it into an
inequality between finitely many values of $\rho$, Step~2 shows that $\rho$ has no zero, and
Step~3 that $1/\rho - 1$ is a nonnegative quadratic form.

\emph{Step 1: the inequality.} Let $a_1, \dots, a_n \in \RR$ and $w_1, \dots, w_n \in \RR^d$, put
$b_k := a_k\rho(w_k)$, and suppose that $\sum_kb_k\langle w_k, x\rangle^2 \ge c|x|^2$ for every $x$,
with $c > 0$. Then $\sum_kb_k \le \sum_ka_k$. Step~2 applies this to pairs $(a_k, w_k)$ of a
weight and a frequency.

For $s \in (0, 1]$ the input $f_s(x) := \sum_ka_k\cos\langle w_k, x\rangle\,e^{-s|x|^2}$ is a real
Schwartz function, a Gaussian times functions whose derivatives are all bounded. Put
$h(x) := \sum_kb_k\cos\langle w_k, x\rangle$ and $A := \sum_k|a_k|$. By the identity above, and as
$|x - y| \le 1 + |y|$ for $|x| \le 1$,
\[
  \bigl|(\mu * f_s)(x) - h(x)\bigr| \le A\int\bigl(1 - e^{-s|x - y|^2}\bigr)\,\mu(dy)
  \le A\int\bigl(1 - e^{-s(1 + |y|)^2}\bigr)\,\mu(dy) \qquad (|x| \le 1),
\]
and the right side tends to $0$ as $s \to 0$, by dominated convergence. Next, $|\cos y - 1 +
y^2/2| \le y^4$ for $|y| \le 1$. With $K := \sum_k|b_k||w_k|^4$, as long as $|x||w_k| \le 1$ for
every $k$,
\[
  h(x) \le h(0) - \tfrac12\sum_kb_k\langle w_k, x\rangle^2 + K|x|^4
  \le h(0) - \tfrac c2|x|^2 + K|x|^4 ,
\]
so there is $r_0 \in (0, 1]$ with $h(x) \le h(0) - \frac c4|x|^2$ for $|x| \le r_0$. Let
$0 < r \le r_0$, and choose $s \in (0, 1]$ with $|\mu * f_s - h| < cr^2/8$ on the unit ball. The output
$u := \mu * f_s$ is continuous, so it attains its maximum over the closed ball of radius $r$ at some
point $x_s$. On the sphere $|x| = r$, $u(x) < h(x) + cr^2/8 \le h(0) - cr^2/8 < u(0)$, so
$|x_s| < r$, and $x_s$ is a local maximum of $u$. The principle gives
$h(0) - cr^2/8 < u(0) \le u(x_s) \le f_s(x_s)$. From $1 - y^2/2 \le \cos y \le 1$ and
$1 - e^{-t} \le t$, $f_s(x) \le \sum_ka_k + A_1|x|^2$ with $A_1 := \sum_k|a_k|(|w_k|^2/2 + 1)$.
Hence $\sum_kb_k = h(0) < \sum_ka_k + (A_1 + c/8)r^2$ for every $r \in (0, r_0]$, and
$\sum_kb_k \le \sum_ka_k$.

\emph{Step 2: $\rho$ has no zero.} By continuity there is $t > 0$ with $\rho(te_i) \ge \frac12$ for
$i = 1, \dots, d$. Suppose $\rho(\omega_0) = 0$, and apply Step~1 to the pairs $(1, te_i)$,
$i = 1, \dots, d$, and $(a_0, \omega_0)$. Whatever $a_0$, the hypothesis holds with $c = t^2/2$,
since $b_0 = a_0\rho(\omega_0) = 0$. The conclusion $\sum_i\rho(te_i) \le d + a_0$ fails for
$a_0 = -d - 1$.

\emph{Step 3: the reciprocal.} Put $m := 1/\rho - 1$. Step~1 with $a_k := b_k/\rho(w_k)$ reads: if
$\sum_kb_k\langle w_k, x\rangle^2 \ge c|x|^2$ for every $x$, with $c > 0$, then
$\sum_kb_km(w_k) \ge 0$. In the rest of this step a pair $(b_k, w_k)$ names the weight $b_k$ of this
form, not the weight $a_k$ of Step~1. If instead $\sum_kb_k\langle w_k, x\rangle^2 = 0$ for every $x$, the pairs
$(b_k, w_k)$ together with $(\varepsilon, e_i)$, $i = 1, \dots, d$, have the form $\varepsilon|x|^2$,
so $\sum_kb_km(w_k) + \varepsilon\sum_im(e_i) \ge 0$ for every $\varepsilon > 0$. The same holds
for $-b$, so $\sum_kb_km(w_k) = 0$. In the same way the pairs $(1, \omega)$ and
$(\varepsilon, e_i)$ give $m(\omega) \ge 0$. Let $C_{ij} := \frac12\bigl(m(e_i + e_j) - m(e_i) -
m(e_j)\bigr)$, a symmetric matrix. Fix $\omega$. The pair $(1, \omega)$ and, for all $i$ and $j$,
the pairs $(-\omega_i\omega_j/2, e_i + e_j)$, $(\omega_i\omega_j/2, e_i)$ and
$(\omega_i\omega_j/2, e_j)$ have the form $\langle\omega, x\rangle^2 -
\sum_{i,j}\omega_i\omega_jx_ix_j = 0$. So
$m(\omega) = \sum_{i,j}C_{ij}\omega_i\omega_j = \omega\cdot C\omega$. This is $\ge 0$ for every
$\omega$, so $C$ is positive semidefinite, and $1/\widehat\mu(\omega) = 1 + \omega\cdot C\omega$.

\emph{Rotation invariance.} If $\mu$ is rotation invariant, $\widehat\mu$ is radial, so
$\omega\cdot C\omega = \tau|\omega|^2$ with $\tau := e_1\cdot Ce_1 \ge 0$, and
$\widehat\mu(\omega) = (1 + \tau|\omega|^2)^{-1}$: $\mu$ is the Helmholtz resolvent. Any positive
semidefinite $C'$ with $1/\widehat\mu(\omega) = 1 + \omega\cdot C'\omega$ has the quadratic form
$\tau|\omega|^2$, so $C' = \tau I$, a symmetric matrix being determined by its quadratic form. When
$\tau = 0$, $\widehat\mu = 1$ and $\mu = \delta_0$, by the uniqueness of the transform.
% CHANGED (blueprint round, 2026-11-16, review flag "So (3) holds already when the principle"): the
% closing remark is dropped. It claimed (3) under the principle required only of the inputs f_s, at the
% maximizers of Step 1; stage_lemma_3 takes the full principle (HasImplicitEulerPrinciple), and no
% declaration states the weaker hypothesis. The old text: "So (3) holds already when the principle is
% required only of the inputs $f_s$, and only at a point where the output is largest on a small ball
% about a strict maximum of its limit."
\end{proof}

% draft: the comment after Proposition 4.9, line 672, without the overview of the proof, which is
% now the opening of the proof of (3).
So a rotation-invariant stage that never lets a local maximum of its output exceed its input there
is a Helmholtz resolvent, and nothing else is. Without rotation invariance it is a resolvent
$(1 + \omega\cdot C\omega)^{-1}$. The converse needs no class of laws.

% draft: the paragraph before Theorem 4.10, line 690.
Inside stratum 2 the lemma selects. Between its knots a family covariant under $q$ is not determined
by its ladder, not even up to a relabelling of the scales
(Proposition~\ref{prop:ladder-knot-kernel}(2)). Required of one octave stage, the principle fixes the
ladder through it. Required of every sliding octave stage, it fixes the family up to the gauge, and full
covariance follows.

% draft: Theorem 4.10 (section 4.3), lines 692-697 of draft/spatial-hemigroup-pyramids.md.
% CHANGED: the grid case moved to chapter 5 (adr/0004), prop:grid-frame-step-not-selected
% Clause (4) keeps the Gaussian stage, the continuum frame step and the axis passes; "on the
% continuum and on the grid" reads "on the continuum". The Gaussian stage is named as N(0, 2aI),
% a > 0, the reading of the draft's "the Gaussian stage".
% CHANGED: clause (1)'s "the stages between the rungs are free" is stated as what the proof shows,
% through prop:ladder-knot-kernel(2) (blind review, flag 7).
% CHANGED: clause (3) names the Matern pyramid mu^M with its constants q, c, as (mu_{s,t}) is the
% given family; clause (4) binds c, t > 0 (second blind review, flag 6).
% The proof of the freedom in (1) writes out the two couplings of prop:ladder-knot-kernel(2) for the
% step L_{sqrt(tau)} and proves them isotropic, which that proposition states for one family only
% (second blind review, flag 1). The statement is unchanged.
% CHANGED (fifth blind pass, 2026-11-15, flag 35): the Assignment names every node the proof reads,
% the [A] node prop:ladder-knot-kernel(2), whose two couplings the freedom writes out, and the [T]
% nodes prop:stage-lemma(1) and (3), lem:bessel-profile(1)-(3) and thm:matern-pyramid(1); it named
% the structure lemmas and lem:synthesis only.
% CHANGED (fidelity review, 2026-11-15, R22): the tag adds the converse of (2) as the node states
% it, for every isotropic q-covariant family in the canonical discrete gauge with the exponents
% F^1_q(chi(t) .), chi in G_q (selection_stratum_two_2_converse_family, Lean core), and the Matern
% pyramid read in the gauge chi as such a family (maternPyramid_gauge_family, reading A34 through
% thm:matern-pyramid(1)) with its principle (selection_stratum_two_2_converse_gauge). The
% statement is unchanged.
\begin{theorem}[the selection, inside stratum 2]
\label{thm:selection-stratum-two}
\lean{Pyramids.selection_stratum_two_1, Pyramids.selection_stratum_two_1_ladder,
Pyramids.selection_stratum_two_1_free, Pyramids.selection_stratum_two_2,
Pyramids.selection_stratum_two_2_converse, Pyramids.selection_stratum_two_3,
Pyramids.selection_stratum_two_4_gaussian, Pyramids.selection_stratum_two_4_frame,
Pyramids.selection_stratum_two_4_axis, Pyramids.selection_stratum_two_2_converse_family,
Pyramids.maternPyramid_gauge_family, Pyramids.selection_stratum_two_2_converse_gauge}
% CHANGED (wave 3b, 2026-11-15): all nine declarations are machine-checked, so statement and proof
% are marked proved. The line marked the node not ready.
\leanok
\uses{prop:stage-lemma, def:implicit-euler-principle, thm:matern-pyramid, lem:bessel-profile,
def:family-dimension-d, lem:family-structure, lem:family-dilations, lem:synthesis,
prop:ladder-knot-kernel, thm:cited-A39, def:frame-matern-pyramid}
Let $(\mu_{s,t})$ be an isotropic $q$-covariant family in dimension $d$, in the canonical discrete
gauge, with exponents $\Psi_t(|\omega|) := -\log\widehat{\mu_t}(\omega)$. Write $F^1_q$ for $F_q$ at
$c = 1$.
\begin{enumerate}
\item \emph{(One stage.)} For one $t > 0$, $\mu_{t,qt}$ has the principle if and only if
$\Psi_t(r) = \sum_{n \ge 1}\log(1 + \tau q^{-2n}r^2)$ for some $\tau > 0$. Then the whole ladder
through $t$ is that of the \Matern{} pyramid, $\mu_{q^kt,q^{k+1}t} = L_{\sqrt\tau q^k}$ for
$k \in \ZZ$, and the stages between the rungs are free: the isotropic $q$-covariant families in
the canonical discrete gauge with this ladder are not unique, and not unique up to a relabelling in
$\mathcal{G}_q$ either.
\item \emph{(Every octave stage.)} If $\mu_{t,qt}$ has the principle for all $t \in [1, q)$, then
there is $\chi \in \mathcal{G}_q$ with $\Psi_t = F^1_q(\chi(t)\,\cdot)$. So the family is the
\Matern{} pyramid of Theorem~\ref{thm:matern-pyramid} read in the gauge $\chi$, and it is fully
covariant, with $S_\lambda = \chi^{-1}(\lambda\chi)$. Conversely every such family has the
property.
\item \emph{(Which stages.)} Let $(\mu^{\mathrm M}_{s,t})$ be the \Matern{} pyramid of
Theorem~\ref{thm:matern-pyramid} with the constants $q$ and $c > 0$. Its stage
$\mu^{\mathrm M}_{s,t}$ with $0 < s < t$ has the principle if and only if $t = qs$, and
$\mu^{\mathrm M}_{0,t}$ with $t > 0$ never has it. In particular two octaves together fail, as does every stage
between other scales.
\item \emph{(Not selected.)} Let $c > 0$ and $t > 0$. A Gaussian stage $N(0, 2aI)$, $a > 0$, fails
the principle. The octave step at $t$ of the frame \Matern{} pyramid with the constant $c$
(Definition~\ref{def:frame-matern-pyramid}),
$\prod_j(1 + c^2t^2\omega_j^2)^{-1}$, fails it in $d \ge 2$, on the continuum, while each of its
axis passes has it.
\end{enumerate}
% CHANGED (fidelity review, 2026-11-15, R22 follow-up): "All nine declarations" reads "All
% twelve", the tag having gained the three declarations of R22.
% CHANGED (wave 3b, 2026-11-15): the Assignment follows the checked proof, in which entry A39
% carries one step of the freedom, the relabelling, and the two families differ because the
% identity is a relabelling; the line "All nine declarations are machine-checked" is added. The old
% wording: "carries two steps of the freedom in (1): two laws with equal exponents have equal Levy
% measures, which separates the two families built there at the scale sqrt q, and excludes a
% relabelling between them. Everything else is proved here".
% CHANGED (fidelity review, 2026-11-15, R26/R50/R69): the status line lists the declarations' hypotheses beside the node's (LINKAGE rule 6).
% CHANGED (fidelity review, 2026-11-15, R27): the status line says that the principle's predicate
% omits the definition's symmetry and probability, and that they hold at every use.
\statusA\quad\ledger{A39}\quad\emph{(\textbf{Assignment.} \ledger{A39}, [\sscite{sato1999levy},
Thm.~8.1(ii), p.~38], the uniqueness of the generating triplet, carries one step of the freedom
in (1): two laws with equal exponents have equal \Levy{} measures, which excludes a relabelling
between the two families built there, and so separates them. All twelve declarations are
machine-checked, and only the freedom reads entry A39. Everything
else is proved here: from the stage lemma, Proposition~\ref{prop:stage-lemma}, its converse (3)
for the necessity and its clause (1) for the sufficiency, whose proofs are printed in this
chapter; from the law $L_r$ and its transform, the radial density of $\log(1 + \tau|\omega|^2)$ and
the bounds on $b_d$ (Lemma~\ref{lem:bessel-profile}(1)--(3)); from the octave step
$\mu_{s,qs} = L_{cs}$ of the \Matern{} pyramid (Theorem~\ref{thm:matern-pyramid}(1)); from the
structure of a $q$-covariant family (Lemmas~\ref{lem:family-structure} and
\ref{lem:family-dilations}); and, for the freedom in (1), from the uniform and the split coupling of
Proposition~\ref{prop:ladder-knot-kernel}(2), which carries its own entries and which the proof
writes out for the step $L_{\sqrt\tau}$, and from Lemma~\ref{lem:synthesis}, which builds the two
families and carries its own use of entry A39, with their isotropy proved here.
\textbf{Not carried}: the principle on a grid after decimation, which
is open (the remarks after Proposition~\ref{prop:principle-other-readings}).
\textbf{The principle.} \texttt{HasImplicitEulerPrinciple} is the condition of
Definition~\ref{def:implicit-euler-principle} without its symmetry and probability, so the positive
conclusions (the sufficiency in (1), the converse of (2), $t = qs$ in (3), the axis passes in (4))
do not state them. They hold at every use: the stages of an isotropic family are probability
measures, symmetric by the rotation invariance, and the axis pass is a probability by hypothesis,
symmetric as its transform is real.
\textbf{Hypotheses} (the node's, each with what its declarations require): $q > 1$, part of a
$q$-covariant family: same; $d \ge 1$: same in \texttt{selection\_stratum\_two\_1}, \texttt{\_1\_free},
\texttt{\_2}, \texttt{\_2\_converse\_gauge}, \texttt{maternPyramid\_gauge\_family}, \texttt{\_3},
\texttt{\_4\_gaussian}; weaker, not required, in \texttt{\_1\_ladder}, \texttt{\_2\_converse},
\texttt{\_2\_converse\_family}, \texttt{\_4\_axis}; $(\mu_{s,t})$ isotropic and $q$-covariant in the
canonical discrete gauge: same, \texttt{IsIsotropicFamily} and \texttt{IsCovariantAt} at $q$ with
$S_q t = qt$, in \texttt{\_1}, \texttt{\_1\_ladder}, \texttt{\_2}, \texttt{\_2\_converse\_family}
(whose proof does not read the covariance); absent from \texttt{\_1\_free}, \texttt{\_3} and (4),
which do not read the family; $\Psi_t = -\log\widehat{\mu_t}$: same, stated through the transform
$\widehat{\mu_t} = e^{-\Psi_t}$, with no logarithm taken; (1) one $t > 0$: same; $\tau > 0$: same;
the representation of $\Psi_t$, for the ladder: same in \texttt{\_1\_ladder}; (2) the principle for
all $t \in [1, q)$: same; the converse, $\chi \in \mathcal{G}_q$ and $\Psi_t = F^1_q(\chi(t)\,\cdot)$
for an isotropic $q$-covariant family: same in \texttt{\_2\_converse\_family}; stronger alone in
\texttt{\_2\_converse} and \texttt{\_2\_converse\_gauge}, which require a \Matern{} pyramid $\nu$ at
$c = 1$ read through $\chi$, of which
\texttt{maternPyramid\_gauge\_family} shows the relabelling such a family; (3) the \Matern{} pyramid
with $q > 1$, $c > 0$: same, every $\mu$ with \texttt{IsMaternPyramid q c}; $0 < s < t$, and $t > 0$
for $\mu^{\mathrm M}_{0,t}$: same; (4) $a > 0$: same; $c > 0$, $t > 0$: same in \texttt{\_4\_frame};
weaker in \texttt{\_4\_axis}, which requires neither; the octave step of the frame \Matern{} pyramid:
same, every probability measure with the product transform; $d \ge 2$: same in
\texttt{\_4\_frame}; the axis pass: same, every probability measure with the transform
$(1 + c^2t^2\omega_j^2)^{-1}$.)}
\end{theorem}

% CHANGED (fifth blind pass, 2026-11-15, flags 34, 36, 37): nu* is shown a Levy measure in the
% proof, by bounding the sum of the dilates; it was attributed to prop:ladder-knot-kernel(1), which
% assumes a family. The identifications mu_{t,qt} = L_{sqrt(tau)} and mu_{t,qt} = L_{chi(t)} cite
% lem:bessel-profile(1) and the uniqueness of the transform; g(sqrt q) lies in (1, q), so lambda in
% (0, 1). The proof gets the leanok mark and the uses of the statement, which it lacked.
\begin{proof}
\leanok
\uses{prop:stage-lemma, def:implicit-euler-principle, thm:matern-pyramid, lem:bessel-profile,
def:family-dimension-d, lem:family-structure, lem:family-dilations, lem:synthesis,
prop:ladder-knot-kernel, thm:cited-A39}
% CHANGED (wave 3b, 2026-11-15): the first paragraph of (1) follows the machine-checked route, on
% the transforms, without the logarithm of mu_t-hat. The old paragraph wrote Psi_qt - Psi_t =
% log(1 + tau r^2), Psi_qt(r) = Psi_t(qr), summed the differences at q^{-n} r and let
% Psi_t(q^{-N} r) tend to 0; the converse read the stage's transform as e^{-(Psi_t(q|omega|) -
% Psi_t(|omega|))}. It did not say why the stage is symmetric, which the stage lemma requires.
\emph{(1)} The stage $\mu_{t,qt}$ is rotation invariant by (I3), and symmetric, $-I$ being
orthogonal, so Proposition~\ref{prop:stage-lemma}(3) gives
$\widehat{\mu_{t,qt}}(\omega) = (1 + \tau|\omega|^2)^{-1}$ with $\tau \ge 0$; and $\tau > 0$ by
(ND), as $\tau = 0$ makes $\mu_{t,qt} = \delta_0$. In the canonical discrete gauge
$D_q\mu_{0,t} = \mu_{0,qt}$, and $\mu_{0,qt} = \mu_{t,qt} * \mu_{0,t}$ by (I1), so
$\widehat{\mu_t}(q\omega) = (1 + \tau|\omega|^2)^{-1}\,\widehat{\mu_t}(\omega)$. Applied at
$q^{-n}\omega$, $n = 1, \dots, N$, this gives
\[
  \widehat{\mu_t}(\omega) = \widehat{\mu_t}(q^{-N}\omega)\prod_{n=1}^N
    \bigl(1 + \tau q^{-2n}|\omega|^2\bigr)^{-1} ,
\]
and $\widehat{\mu_t}(q^{-N}\omega) \to \widehat{\mu_t}(0) = 1$, as $\widehat{\mu_t}$ is
continuous. The product tends to $\exp\bigl(-\sum_{n \ge 1}\log(1 + \tau q^{-2n}|\omega|^2)\bigr)$,
the series converging since $\log(1 + x) \le x$. This is the form of $\Psi_t$, read on the
transform. Conversely, if $\Psi_t$ has that form, the series is $F^1_q(\sqrt\tau\,|\omega|)$, and
the telescoping $F^1_q(qr) = \log(1 + r^2) + F^1_q(r)$ gives
$\widehat{\mu_{t,qt}}(\omega) = \widehat{\mu_t}(q\omega)/\widehat{\mu_t}(\omega) =
(1 + \tau|\omega|^2)^{-1}$. So $\mu_{t,qt} = L_{\sqrt\tau}$ by Lemma~\ref{lem:bessel-profile}(1)
and the uniqueness of the transform, and it has the principle by
Proposition~\ref{prop:stage-lemma}(1). The stage at $q^kt$, $k \in \ZZ$, is the dilate
$D_{q^k}L_{\sqrt\tau} = L_{\sqrt\tau q^k}$, by covariance applied $|k|$ times.

For the freedom between the rungs it suffices to take $t = 1$. If $(\mu'_{s,u})$ is an isotropic
$q$-covariant family with $S_qu = qu$ and the steps $\mu'_{q^k,q^{k+1}} = L_{\sqrt\tau q^k}$,
$k \in \ZZ$, then $(\mu'_{s/t,u/t})$ is one with the ladder through $t$ of (1). A relabelling by
$g \in \mathcal{G}_q$ between two families so rescaled is a relabelling by
$u \mapsto g(tu)/t$, again in $\mathcal{G}_q$, between the families before rescaling. So it
suffices to build two isotropic families with $S_qu = qu$ and the ladder $(L_{\sqrt\tau q^k})_k$
that are not relabellings of each other by any element of $\mathcal{G}_q$.
Proposition~\ref{prop:ladder-knot-kernel}(2) gives them for the step $\varrho := L_{\sqrt\tau}$.

Let $\nu$ be the measure with the radial density
$b_d(|x|/\sqrt\tau)\bigl(|S^{d-1}|\,|x|^d\bigr)^{-1}$. By Lemma~\ref{lem:bessel-profile}(2) at
$\beta = \tau$, $\int(1 - \cos\langle\omega, x\rangle)\,\nu(dx) = \log(1 + \tau|\omega|^2)$.
In polar coordinates $\int(1 \wedge |x|^2)\,\nu(dx) = \int_0^\infty(1 \wedge r^2)\,b_d(r/\sqrt\tau)\,dr/r$,
which is finite since $b_d \le 2$ and $b_d(y) \le C_de^{-y/2}$ (Lemma~\ref{lem:bessel-profile}(3)).
So $\nu$ is a rotation-invariant \Levy{} measure, and the infinitely divisible law with the data
$(0, \nu)$ has the transform $(1 + \tau|\omega|^2)^{-1}$. It is $\varrho$, by
Lemma~\ref{lem:bessel-profile}(1), so $\varrho$ is symmetric, rotation invariant, infinitely
divisible and not $\delta_0$. Its logarithmic moment is finite, since $\log_+ r \le r$ and so
$\log_+ r\,b_d(r/\sqrt\tau)/r \le C_d e^{-r/(2\sqrt\tau)}$. Split $\nu = \nu_a + \nu_b$, with
$\nu_a$ the restriction of $\nu$ to $\{|x| \le 1\}$ and $\nu_b$ its restriction to $\{|x| > 1\}$.
Both are symmetric and rotation invariant, carried by disjoint sets, and nonzero since $b_d > 0$.
By Proposition~\ref{prop:ladder-knot-kernel}(2) there are two isotropic $q$-covariant families
with $S_qu = qu$ and the step $\mu_{1,q} = \varrho$, neither a relabelling of the other by an
element of $\mathcal{G}_q$. Covariance gives their ladder,
$\mu_{q^k,q^{k+1}} = D_{q^k}L_{\sqrt\tau} = L_{\sqrt\tau q^k}$. This is the last sentence of (1).

% CHANGED (wave 3b, 2026-11-15): (2) follows the machine-checked route. The old proof defined chi
% on [1, q) from (1) and extended it by chi(q^k t) := q^k chi(t); it read chi(t) as
% (F^1_q)^{-1}(Psi_t(1)), continuous by (I2) and nondecreasing by lem:family-structure, strictly
% increasing by (ND), and in G_q by lem:family-dilations(3) from its limits at 0 and infinity. Here
% every stage has the principle, chi is read from (1) at every t > 0, its monotonicity comes from
% (I1) and |mu-hat| <= 1, and its surjectivity from the intermediate values on [1, q].
\emph{(2)} Every stage $\mu_{t,qt}$, $t > 0$, has the principle. With $q^nt \in [1, q)$, clause (1)
makes the stage at $q^nt$ the law $L_{\sqrt\tau}$, and the stage at $t$ is its dilate
$L_{\sqrt\tau q^{-n}}$, which has the principle by Proposition~\ref{prop:stage-lemma}(1). So clause
(1) applies at every $t > 0$: $\widehat{\mu_t}(\omega) = \exp(-F^1_q(\chi(t)|\omega|))$ with
$\chi(t) := \sqrt{\tau(t)} > 0$, since $\sum_n\log(1 + \tau q^{-2n}r^2) = F^1_q(\sqrt\tau r)$. As
$F^1_q$ is strictly increasing on $[0, \infty)$, $\chi(t)$ is read off $\widehat{\mu_t}(e)$ at a
unit vector $e$, and covariance, $D_{q^n}\mu_{0,t} = \mu_{0,q^nt}$, gives
$\chi(q^nt) = q^n\chi(t)$ for $n \in \ZZ$. If $\chi(t) \le \chi(s)$ with $s < t$, then by (I1)
$\widehat{\mu_{s,t}}(\omega) = \exp\bigl(F^1_q(\chi(s)|\omega|) - F^1_q(\chi(t)|\omega|)\bigr) \ge 1$.
A transform has modulus at most $1$, so $\widehat{\mu_{s,t}} = 1$ and $\mu_{s,t} = \delta_0$,
against (ND). So $\chi$ is strictly increasing on $(0, \infty)$. It is continuous there:
$\exp(-F^1_q(\chi(t))) = \widehat{\mu_t}(e)$ is continuous in $t$ by (I2), and if
$a' > \chi(t_0)$, then $F^1_q(a') > F^1_q(\chi(t_0))$, so $\chi(t) < a'$ for $t$ near $t_0$;
likewise from below. By the intermediate values $\chi([1, q]) \supseteq [\chi(1), q\chi(1)]$, and
with the dilations $\chi$ maps $(0, \infty)$ onto itself. So $\chi \in \mathcal{G}_q$, with
$\chi(0) := 0$. By (I1) the transforms are those of Theorem~\ref{thm:matern-pyramid} at $c = 1$ in
the gauge $\chi$,
$\widehat{\mu_{s,t}}(\omega) = \exp\bigl(-[F^1_q(\chi(t)|\omega|) - F^1_q(\chi(s)|\omega|)]\bigr)$,
and the transform determines the law. The map $S_\lambda := \chi^{-1}(\lambda\chi)$, with
$S_\lambda(0) := 0$, is an increasing bijection of $[0, \infty)$, and reading the transforms at
$\lambda\omega$ gives $D_\lambda\mu_{s,t} = \mu_{S_\lambda s, S_\lambda t}$ for every
$\lambda > 0$. Conversely, for such a family telescoping gives
$\widehat{\mu_{t,qt}}(\omega) = (1 + \chi(t)^2|\omega|^2)^{-1}$, so $\mu_{t,qt} = L_{\chi(t)}$ by
Lemma~\ref{lem:bessel-profile}(1) and the uniqueness of the transform, and
Proposition~\ref{prop:stage-lemma}(1) applies.

% CHANGED (wave 3b, 2026-11-15): (3) follows the machine-checked route, two coefficients of a
% quadratic identity. The old proof read exp[F_q(tr) - F_q(sr)] as a quotient of entire functions
% of u = r^2, with zeros at -q^{2i}/(c^2 t^2) and poles at -q^{2i}/(c^2 s^2); equal to 1 + tau u on
% u >= 0, it would have no pole, which forces t = q^k s, and the telescoped polynomial of degree k
% forces k = 1; for mu^M_{0,t} the reciprocal has infinitely many zeros.
\emph{(3)} The stage $\mu^{\mathrm M}_{s,t}$, $0 \le s < t$, has the transform
$\exp\bigl(-[F_q(t|\omega|) - F_q(s|\omega|)]\bigr)$, real and even, so it is symmetric. If it
has the principle, Proposition~\ref{prop:stage-lemma}(3) along a unit vector $e$ gives
$\exp D(x) = 1 + kx^2$ for $x \ge 0$, with $D(x) := F_q(tx) - F_q(sx)$ and
$k := e\cdot Ce \ge 0$. The telescoping $F_q(qy) = \log(1 + c^2y^2) + F_q(y)$ at $y = tx$ and
$y = sx$ gives $D(qx) = D(x) + \log(1 + c^2t^2x^2) - \log(1 + c^2s^2x^2)$, so with $u = x^2$
\[
  (1 + kq^2u)(1 + c^2s^2u) = (1 + c^2t^2u)(1 + ku) , \qquad u \ge 0 .
\]
% CHANGED (2026-11-23, review round 3, s4-5.S16): wording.
Both sides are polynomials in $u$ of degree at most $2$ with the constant term $1$, so their equality at $u = 1$ and $u = 2$ makes the coefficients of $u^2$ and of $u$ agree:
$kc^2(t^2 - q^2s^2) = 0$ and $c^2(t^2 - s^2) = k(q^2 - 1)$. If $k = 0$, the second gives $t = s$,
which is excluded; so $k > 0$, and the first gives $t = qs$. At $s = 0$ this reads $t = 0$, which is
excluded: $\mu^{\mathrm M}_{0,t}$ never has the principle. Conversely
$\mu^{\mathrm M}_{s,qs} = L_{cs}$ (Theorem~\ref{thm:matern-pyramid}(1)) has the principle by
Proposition~\ref{prop:stage-lemma}(1).

% CHANGED (wave 3b, 2026-11-15): (4) follows the machine-checked route, two points on a line. The
% old proof: "The Gaussian reciprocal e^{a|omega|^2} is not a polynomial, so prop:stage-lemma(3)
% excludes it. The frame step's reciprocal prod_j(1 + c^2 t^2 omega_j^2) is a polynomial of degree
% 2d > 2 in omega, so it is not 1 + omega.C omega, and the same proposition excludes it." It did
% not say why the two laws are symmetric, which the stage lemma requires.
\emph{(4)} The Gaussian stage and the frame step have real even transforms, so they are symmetric.
If the Gaussian stage had the principle, Proposition~\ref{prop:stage-lemma}(3) along a unit vector
would give $e^{ax^2} = 1 + kx^2$ for every $x$; at $x = 1$ and $x = \sqrt2$,
$(1 + k)^2 = e^{2a} = 1 + 2k$, so $k = 0$ and $a = 0$, which is excluded. If the frame step had it,
the same proposition would give $\prod_j(1 + c^2t^2\omega_j^2) = 1 + \omega\cdot C\omega$. On the
diagonal $\omega = x(e_1 + e_2)$, which $d \ge 2$ allows, this reads
$(1 + Kx^2)^2 = 1 + Qx^2$ with $K := c^2t^2 > 0$ and $Q := (e_1 + e_2)\cdot C(e_1 + e_2)$; at
$x = 1$ and $x = \sqrt2$ it gives $2K^2 = 0$, which is excluded. The axis pass along $e_j$, the
Laplace law on the line $\RR e_j$ with transform $(1 + c^2t^2\omega_j^2)^{-1}$, is
Proposition~\ref{prop:stage-lemma}(1) with the degenerate $C = c^2t^2e_je_j^{\tp}$.
\end{proof}

% draft: "The grid is Proposition 5.11(4)", line 707; the grid case moved to chapter 5 (adr/0004).
% Moved out of the proof of (4) to here (second blind review, flag 14).
Clause (4) is stated on the continuum. Its grid case is
Proposition~\ref{prop:grid-frame-step-not-selected}: on $h\ZZ^d$, $d \ge 2$, the frame mask fails the
principle on every finite symmetric neighbourhood, while each of its axis passes has it on every
neighbourhood that contains the two neighbours along its axis.

% draft: the comment after Theorem 4.10, line 699.
Why a Gaussian step fails and the resolvent passes can be seen with a nonnegative input: two narrow
bumps at $\pm be_1$ with nothing between them. A kernel with a $C^2$ density and a nondegenerate
maximum at $0$ merges them into one maximum at the origin when $b$ is small against its width, so
that the Hessian of the density stays negative definite on $|x| \le b$, and the bumps are narrow
against $b$. There the output is positive and the input is zero, so $u > f$ at a maximum. The
Gaussian is such a kernel. The resolvent's kernel is not smooth at its centre, with a corner on the
line and a singularity in the plane (Proposition~\ref{prop:matern-generator-octave}(3)), and it
escapes. The picture is not a proof here, and clause (4) does not rest on it.

% draft: the paragraph after the proof of Theorem 4.10, line 709.
So the case for the \Matern{} pyramid in the plane is not counting, which fails for every candidate
of Section~\ref{sec:matern-counting} and survives on the Radon projections
(Theorem~\ref{thm:matern-pyramid}(3)). Non-enhancement read at a finite step remains. For an
isotropic family covariant under $q$, required of one octave stage it selects the ladder, and required of
every sliding octave stage it selects the \Matern{} pyramid, up to the gauge, without assuming full
covariance (Theorem~\ref{thm:selection-stratum-two}). Both rest on the converse of the stage lemma,
Proposition~\ref{prop:stage-lemma}(3). On $\ZZ^2$, with the principle on the four axis neighbours
and invariance under the signed permutations of the axes, the same principle selects the five-point
discrete Helmholtz mask among the masks with a summable inverse, and it rejects the frame mask
(Proposition~\ref{prop:grid-implicit-euler-principle}).

% draft: Remark 4.11(1) (section 4.3), line 713. It asserts and argues, so it is a proposition.
% CHANGED: made a proposition, with the label prop:principle-other-readings (not in the brief's
% table); the remark's two one-line arguments are its proof. The family of (1) is named by its
% transforms, and "selects F" is stated as an equivalence with the converse.
% CHANGED: clause (1) adds the hypothesis F(0) = 0, a narrowing: the family sees only differences of
% F, so F + const gives the same family, and the proof's e^{F(t|omega|)} = 1/mu_{0,t}-hat needs
% F(0) = 0 (blind review, flag 1).
% CHANGED (fifth blind pass, 2026-11-15, flags 38, 39): lem:bessel-profile leaves the uses, which
% neither statement nor proof reads; the proof gets the uses it reads. The annotation says that the
% hypothesis "F not identically 0" is implied by (ND), and why the statement keeps it.
% CHANGED (blueprint round, 2026-11-16, review flag "the principle read along other axes"): the
% title drops "axes", which collides with the spatial axes; the statement is unchanged.
\begin{proposition}[the principle read from the origin and at every stage]
\label{prop:principle-other-readings}
\lean{Pyramids.principle_other_readings_1, Pyramids.principle_other_readings_2}
% CHANGED (wave 3b, 2026-11-15): both declarations are machine-checked, so statement and proof are
% marked proved. The line marked the node not ready.
\leanok
\uses{prop:stage-lemma, def:implicit-euler-principle, def:family-dimension-d}
Let $d \ge 1$.
\begin{enumerate}
% CHANGED (2026-11-15, a strengthening)
% F need not be continuous: the proof reads F only through the kernel's transform along one ray.
% The old wording: "Let $F$ be continuous on $[0, \infty)$ with $F(0) = 0$, not".
\item \emph{(From the origin.)} Let $F$ be a real function on $[0, \infty)$ with $F(0) = 0$, not
identically $0$, and let
$(\mu_{s,t})$ be a family in dimension $d$ with
$\widehat{\mu_{s,t}}(\omega) = \exp\bigl(-[F(t|\omega|) - F(s|\omega|)]\bigr)$. If the kernel
$\mu_{0,t}$ has the principle for one $t > 0$, then $F(r) = \log(1 + \tau r^2)$ for some
$\tau > 0$, the \Matern{} family of index $1$, and then every kernel $\mu_{0,t}$ has it. Its octave
steps, with transform $(1 + \tau t^2|\omega|^2)/(1 + \tau q^2t^2|\omega|^2)$ for a ratio $q > 1$, are
not resolvents $(1 + \omega\cdot C\omega)^{-1}$.
\item \emph{(Every stage.)} No symmetric family in dimension $d$ has the principle at every stage
$\mu_{s,t}$, $0 \le s < t$.
\end{enumerate}
% CHANGED (wave 3b, 2026-11-15): the status line says the declarations are machine-checked. The old
% line: "(Cites nothing beyond Proposition~\ref{prop:stage-lemma}.)"
% CHANGED (fidelity review, 2026-11-15, R26/R50/R69): the status line lists the declarations' hypotheses beside the node's (LINKAGE rule 6).
\statusT\quad\emph{(Cites nothing beyond Proposition~\ref{prop:stage-lemma}. Both declarations are
machine-checked. In (1) the hypothesis that $F$ is not identically $0$ is implied by (ND): for
$F \equiv 0$ every $\mu_{s,t}$ is $\delta_0$. It is stated because the declaration takes it, and it
is what the proof reads for $\tau > 0$. \textbf{Hypotheses} (the node's, each with what its
declarations require): $d \ge 1$: same in \texttt{principle\_other\_readings\_1}; weaker in
\texttt{\_2}, which does not require it; (1) $F$ real on $[0, \infty)$: same, a real function on $\RR$
read at $t|\omega| \ge 0$ only; $F(0) = 0$: same; $F$ not identically $0$: same, $F(r) \ne 0$ for
some $r \ge 0$; $(\mu_{s,t})$ a family in dimension $d$: same, \texttt{IsFamily}; its transforms:
same, for $0 \le s \le t$; one $t > 0$ with the principle: same; the ratio $q > 1$: same, quantified
in the conclusion; (2) a symmetric family in dimension $d$: same, \texttt{IsSymmetricFamily}.)}
\end{proposition}

\begin{proof}
\leanok
\uses{prop:stage-lemma, def:family-dimension-d}
% CHANGED (wave 3b, 2026-11-15): the proof follows the machine-checked route. The old proof of (1)
% read the kernel as rotation invariant and took e^{F(t|omega|)} = 1 + tau'|omega|^2 on all of R^d,
% and compared the coefficients of r^4 in the octave step's identity; the checked proof reads the
% stage lemma along one unit vector and evaluates the identities at two points. The old proof of (2)
% took 0 < s < t < u, compared the terms of degree four, and said that a product of two polynomials
% vanishing identically has a factor vanishing identically; the checked proof takes s, t, u = 0, 1,
% 2, evaluates at two points of a ray, and shows the factor vanishes through C_i = B_i^T B_i.
\emph{(1)} The kernel $\mu_{0,t}$ has the transform $e^{-F(t|\omega|)}$, as $F(0) = 0$, which is
real and even, so $\mu_{0,t}$ is symmetric. By Proposition~\ref{prop:stage-lemma}(3),
$1/\widehat{\mu_{0,t}}(\omega) = 1 + \omega\cdot C\omega$ with $C$ positive semidefinite. Along a
unit vector $e$, with $k := e\cdot Ce \ge 0$, this reads $e^{F(tx)} = 1 + kx^2$ for $x \ge 0$, so
$F(r) = \log(1 + \tau r^2)$ for $r \ge 0$, with $\tau := k/t^2 \ge 0$, and $\tau > 0$ as $F$ is not
identically $0$. Then $\widehat{\mu_{0,t'}}(\omega) = (1 + \tau t'^2|\omega|^2)^{-1}$ for every
$t' > 0$, which has the principle by Proposition~\ref{prop:stage-lemma}(1). Suppose the octave
step $\mu_{t',qt'}$, with transform $(1 + \tau t'^2|\omega|^2)/(1 + \tau q^2t'^2|\omega|^2)$, were
$(1 + \omega\cdot C'\omega)^{-1}$ for a real matrix $C'$. Along $e$, with $a := \tau t'^2 > 0$ and
$K := e\cdot C'e$, $(1 + ax^2)(1 + Kx^2) = 1 + aq^2x^2$ for $x \ge 0$. At $x = 1$ and $x = \sqrt2$
this gives $a + K + aK = aq^2$ and $2a + 2K + 4aK = 2aq^2$, so $aK = 0$, hence $K = 0$, and then
$a = aq^2$, $q = 1$, a contradiction.

\emph{(2)} Suppose every stage has the principle. By Proposition~\ref{prop:stage-lemma}(3) there
are positive semidefinite $C_1$, $C_2$, $C_3$ with
$1/\widehat{\mu_{0,1}} = 1 + \omega\cdot C_1\omega$, $1/\widehat{\mu_{1,2}} = 1 + \omega\cdot
C_2\omega$ and $1/\widehat{\mu_{0,2}} = 1 + \omega\cdot C_3\omega$. Write
$Q_i(\omega) := \omega\cdot C_i\omega \ge 0$. The cascade $\mu_{0,2} = \mu_{1,2} * \mu_{0,1}$,
read at $x\omega$, gives $(1 + x^2Q_1(\omega))(1 + x^2Q_2(\omega)) = 1 + x^2Q_3(\omega)$, and at
$x = 1$ and $x = \sqrt2$, as in (1), $Q_1(\omega)Q_2(\omega) = 0$ for every $\omega$. Write
$C_i = B_i^\top B_i$, so that $Q_i(\omega) = |B_i\omega|^2$ and a zero of $Q_i$ is in the kernel of
$B_i$. If $Q_1(\omega_1) > 0$ and $Q_2(\omega_2) > 0$, then $Q_2(\omega_1) = 0 = Q_1(\omega_2)$, so
$Q_1(\omega_1 + \omega_2) = Q_1(\omega_1) > 0$ and $Q_2(\omega_1 + \omega_2) = Q_2(\omega_2) > 0$,
against the product's vanishing at $\omega_1 + \omega_2$. So $Q_1 \equiv 0$ or $Q_2 \equiv 0$: then
$\widehat{\mu_{0,1}} \equiv 1$ or $\widehat{\mu_{1,2}} \equiv 1$, that is $\mu_{0,1} = \delta_0$ or
$\mu_{1,2} = \delta_0$, against (ND).
\end{proof}

% CHANGED (review round 2, batch C, 2026-11-23): new node, answering open question (2) of the theory
% paper's section 7.2 ("whether the output on the coarser grid obeys the principle after
% decimation"), which the prose after it stated as open. Read pointwise, against the decimated
% input, the answer is no, by a finite example.
\begin{proposition}[the pointwise principle after decimation]
\label{prop:principle-after-decimation}
\lean{Pyramids.principle_after_decimation}
\leanok
\uses{prop:stage-lemma, def:implicit-euler-principle}
Let $d \ge 1$, $h > 0$ and $\tau > 0$, let $\kappa$ be the resolvent of
Proposition~\ref{prop:stage-lemma}(2) on $h\ZZ^d$, and let $\Delta_h$ be the $(2d+1)$-point
Laplacian, $\Delta_hu(x) := h^{-2}\sum_{j=1}^d\bigl(u(x + he_j) - 2u(x) + u(x - he_j)\bigr)$. Let
$u(0) := 1$, $u(\pm he_j) := 2$ for $j = 1, \dots, d$, and $u := 0$ elsewhere on $h\ZZ^d$, and put
$f := u - \tau\Delta_hu$.
\begin{enumerate}
\item $\kappa * f = u$, and $f(0) = 1 - 2d\tau/h^2 < u(0)$.
\item \emph{(After decimation.)} For every integer $q \ge 2$, $u$ vanishes at the $2d$ axis
neighbours $\pm qhe_j$ of the origin in the coarse grid $qh\ZZ^d$, so the origin is a strict
maximum of $u$ restricted to $qh\ZZ^d$ over those neighbours, while $u(0) > f(0)$. The decimated
output and the decimated input, $u$ and $f$ restricted to $qh\ZZ^d$, fail the inequality
$u(x) \le f(x)$ of Definition~\ref{def:implicit-euler-principle} at that maximum.
\item \emph{(Before decimation.)} On $h\ZZ^d$ the origin is not a maximum of $u$ over its axis
neighbours, and the pair $(u, f)$ satisfies the principle on the $2d$ axis neighbours at every
point of $h\ZZ^d$, by Proposition~\ref{prop:stage-lemma}(2).
\end{enumerate}
\statusT\quad\emph{(Cites nothing beyond Proposition~\ref{prop:stage-lemma}(2). Machine-checked, in
Lean core. The failure in (2) is that of the pointwise inequality for the decimated pair; no kernel
on the coarse grid is involved, and the statement says nothing of a principle read otherwise
across decimation. \textbf{Hypotheses} (the node's, each with what
\texttt{principle\_after\_decimation} requires): $d \ge 1$: same; $h > 0$: same, on $h\ZZ^d$ read in
the index coordinates $\ZZ^d$, $\Delta_h$ being \texttt{gridLap h}; $\tau > 0$: weaker, the
declaration takes $\tau \ge 0$ and requires $\tau > 0$ only for $f(0) < u(0)$ and the failure in
(2); the resolvent $\kappa$: same, every probability kernel with the transform of
Proposition~\ref{prop:stage-lemma}(2); $q \ge 2$ an integer: same, a natural number, the coarse grid
read in its own indices, $m \mapsto u(qm)$, with axis neighbours $\pm e_j$ standing for
$\pm qhe_j$.)}
\end{proposition}

\begin{proof}
\leanok
\uses{prop:stage-lemma}
% CHANGED (review round 2, batch C, blind register review, 2026-11-23): the identity for kappa is
% read off its transform; the sentence read "by the uniqueness of Fourier coefficients, kappa is the
% random-walk series of that proof, whose series gives the identity".
\emph{(1)} First, $\kappa - \tau\Delta_h\kappa = \delta_0$. Both sides are summable on $h\ZZ^d$,
the transform of the left side is
$\bigl(1 + 4\tau h^{-2}\sum_{j=1}^d\sin^2(h\omega_j/2)\bigr)\widehat\kappa(\omega) = 1$, and a
summable function on $h\ZZ^d$ is determined by its transform. The function $u$ has finite support, so $u$ and
$\Delta_hu$ are summable, and $\Delta_h$ passes onto either factor of a convolution of summable
functions by shifting the summation index. So
\[
\kappa * f = \kappa * u - \tau\,\kappa * \Delta_hu = \kappa * u - \tau\,(\Delta_h\kappa) * u
= (\kappa - \tau\Delta_h\kappa) * u = u .
\]
At the origin, $\Delta_hu(0) = h^{-2}\sum_{j=1}^d(2 - 2 + 2) = 2dh^{-2}$, so
$f(0) = 1 - 2d\tau/h^2$, which is less than $u(0) = 1$ as $\tau > 0$.

\emph{(2)} The $j$-th coordinate of $\pm qhe_j$ is $\pm qh$, and $q \ge 2$, so $\pm qhe_j$ is neither
the origin nor one of the points $\pm he_k$. Hence $u(\pm qhe_j) = 0 < 1 = u(0)$, while
$u(0) > f(0)$ by (1).

\emph{(3)} $u(he_1) = 2 > 1 = u(0)$. The input $f$ has finite support, so it is bounded, and
$u = \kappa * f$ by (1). Proposition~\ref{prop:stage-lemma}(2) gives $u(x) \le f(x)$ at every
maximum of $u$ over the axis neighbours and $u(x) \ge f(x)$ at every minimum.

Every step is machine-checked, with no cited fact.
\end{proof}

% draft: Remark 4.11(1)'s last sentence and Remark 4.11(2)-(3), lines 713-715. Prose: they say what
% the principle does not give; (3) states an open question.
% CHANGED (review round 2, batch C, 2026-11-23): the last sentence, "Whether the output on the
% coarser grid obeys the principle after decimation is open.", is replaced by the answer of
% prop:principle-after-decimation, read pointwise.
A pyramid's own axis is its ladder, and reading the principle along the ladder is a choice. The
principle does not count. It bounds the
value of an output maximum by the input there, and it allows the creation of
Proposition~\ref{prop:creation-plane}. On the grid it is proved
for one smoothing step, and it does not survive decimation: a local maximum of a
decimated record need not be a local maximum on the finer grid, and read pointwise against the
decimated input the principle fails there (Proposition~\ref{prop:principle-after-decimation}).
